% ENVIRONMENTAL GEOLOGY — Geology Upper Sixth — Cours signé OnBuch+
% Official MINESEC syllabus. Compiler avec : tectonic GEO08-environmental-geology.tex
\newcommand{\DOCMATIERE}{Geology}
\newcommand{\DOCNIVEAU}{Upper Sixth}
\newcommand{\DOCMODULE}{Upper Sixth — Geology}
\newcommand{\DOCLECON}{8}
\newcommand{\DOCTITRE}{ENVIRONMENTAL GEOLOGY}
% OnBuch+ — préambule « cours signés » ENRICHI (compilé avec tectonic / XeLaTeX)
% Système visuel « Pop » d'OnBuch+ : fond crème (jamais blanc), bordures encre
% épaisses, ombres « dures » décalées, titres Archivo Black en capitales.
% Version MATHÉMATIQUES : graphes (pgfplots/tikz), boîtes pédagogiques
% (définition, propriété, méthode, attention, le savais-tu, exercice résolu,
% exercice, TP, bilan), page de garde et signature OnBuch+ ancrée en bas.
%
% Macros attendues dans main.tex AVANT \input{preamble} :
%   \DOCMATIERE  \DOCNIVEAU  \DOCMODULE  \DOCLECON (numéro)  \DOCTITRE
\documentclass[11pt]{article}

\usepackage{fontspec}
\usepackage[english]{babel}
\usepackage[a4paper,margin=2cm,top=3.4cm,bottom=2.8cm,headheight=1.4cm,footskip=1.3cm]{geometry}
\usepackage{amsmath,amssymb,amsfonts}
\usepackage{mathtools} % \xmapsto, \coloneqq... souvent utilisés par les modèles
\usepackage{cancel} % simplifications barrées (bilans de forces)
% \abs{z} (module, valeur absolue) et \norm{u} (norme), utilisés par les modèles
\providecommand{\abs}{}\let\abs\relax\DeclarePairedDelimiter{\abs}{\lvert}{\rvert}
\providecommand{\norm}{}\let\norm\relax\DeclarePairedDelimiter{\norm}{\lVert}{\rVert}
\usepackage[table]{xcolor} % [table] : fournit aussi \rowcolors (alternance de lignes)
\usepackage{pagecolor}
\usepackage{tikz}
\usetikzlibrary{arrows.meta,positioning,calc,shapes.geometric,shapes.misc,shapes.arrows,decorations.pathmorphing,decorations.pathreplacing,decorations.markings,patterns,shadows,mindmap,trees,fit,backgrounds,matrix,intersections,angles,quotes}
\usepackage{pgfplots}
\usepackage[version=4]{mhchem} % équations nucléaires \ce{^{235}_{92}U}
\usepackage[european,siunitx]{circuitikz} % schémas électriques
\pgfplotsset{compat=1.18}
\usepgfplotslibrary{fillbetween,groupplots} % aires entre courbes (calcul intégral)
\usepackage{enumitem}
\usepackage{fancyhdr}
% [most] charge toutes les bibliothèques tcolorbox (listings, minted, raster,
% documentation...) et gonfle inutilement la mémoire TeX sur un document long
% (~300 boîtes) : on ne charge que ce qui est réellement utilisé.
\usepackage[skins,breakable]{tcolorbox}
\usepackage{graphicx}
\usepackage{colortbl}
\usepackage{booktabs}
\usepackage{tabularx}
\usepackage{diagbox} % case d'en-tête diagonale des tableaux à double entrée
% Colonnes extensibles centrée / gauche / droite (C, L, R), souvent utilisées par les modèles
\newcolumntype{C}{>{\centering\arraybackslash}X}
\newcolumntype{L}{>{\raggedright\arraybackslash}X}
\newcolumntype{R}{>{\raggedleft\arraybackslash}X}
\usepackage{array}
\usepackage{multicol}
\usepackage{siunitx}
% parse-numbers=false : accepte aussi \SI{2,5 \times 10^{-3}}{...}
\sisetup{locale=UK,output-decimal-marker={.},group-minimum-digits=4,parse-numbers=false}
\DeclareSIUnit\ohm{\ensuremath{\symup{\Omega}}} % U+2126 absent de la police
\DeclareSIUnit\division{div} % réglages d'oscilloscope (ms/div, V/div)
\DeclareSIUnit\tr{tr} % tours (vitesse de rotation, ex. \SI{1500}{\tr\per\minute})
\usepackage{qrcode}
% \degree : commande fréquemment utilisée par les modèles pour un angle ou une
% température en dehors de \SI{}{} (non fournie par les paquets chargés).
\providecommand{\degree}{^{\circ}}
% \qed (fin de démonstration), utilisé par les modèles sans amsthm
\providecommand{\qed}{\hfill$\square$}
% \barre : double barre des valeurs interdites dans un tableau de variations (tkz-tab)
\providecommand{\barre}{\ensuremath{\|}}
% environnement proof (amsthm non chargé) utilisé par les modèles
\ifdefined\proof\else\newenvironment{proof}[1][Proof]{\par\noindent\textit{#1.}\ }{\qed\par}\fi
\usepackage{lastpage}
\usepackage{hyperref}
\hypersetup{hidelinks}

% ============================================================
% Palette « Pop » OnBuch+ — mêmes hex que la classe P (plus_ui.dart)
% ============================================================
\definecolor{poporange}{HTML}{F59321}
\definecolor{popdark}{HTML}{E07A0C}
\definecolor{popcream}{HTML}{FFF6E8}
\definecolor{popink}{HTML}{1C1714}
\definecolor{poppurple}{HTML}{7A5AE0}
\definecolor{popgreen}{HTML}{2E9E6A}
\definecolor{poppink}{HTML}{E0457B}
\definecolor{popblue}{HTML}{2F7DE1}
\definecolor{popgold}{HTML}{C98A12}
\definecolor{popmuted}{HTML}{8A7F76}
\definecolor{popline}{HTML}{EADFD2}
\definecolor{popgray}{HTML}{8A7F76}
\colorlet{popgrey}{popgray}
% Alias tolérants (le modèle nomme parfois les couleurs différemment)
\colorlet{popwhite}{white}
\colorlet{popblack}{popink}
\colorlet{popred}{poppink}
\colorlet{popyellow}{popgold}
% Teintes claires pour fonds de tableaux / zones
\colorlet{popblueL}{popblue!10!white}
\colorlet{poporangeL}{poporange!14!white}
\colorlet{popgreenL}{popgreen!12!white}
\colorlet{poppurpleL}{poppurple!12!white}
\colorlet{poppinkL}{poppink!12!white}
\colorlet{popgoldL}{popgold!14!white}

\pagecolor{popcream}

% ============================================================
% Typographie
% ============================================================
\defaultfontfeatures{Ligatures=TeX}
\setmainfont{PlusJakartaSans-Medium.ttf}[Path=../../../fonts/,BoldFont=PlusJakartaSans-Bold.ttf]
% Maths en sans-serif (Fira Math) avec chiffres et lettres droites de Plus
% Jakarta, pour que les formules se fondent dans le texte.
\usepackage[math-style=ISO,bold-style=ISO]{unicode-math}
\setmathfont{FiraMath-Regular.otf}[Path=../../../fonts/]
\setmathfont{PlusJakartaSans-Medium.ttf}[Path=../../../fonts/,range={up/num,up/{Latin,latin},\mathup}]
% (icomma retiré : décimales avec point en anglais)
% Caractères mathématiques Unicode absents des polices de texte (Plus Jakarta,
% Archivo Black) : rendus par leur équivalent mathématique, en texte comme en
% formule — sinon ils s'affichent en carré vide (ex. « Arithmétique dans ℤ »).
\usepackage{newunicodechar}
\newunicodechar{ℝ}{\ensuremath{\mathbb{R}}}
\newunicodechar{ℤ}{\ensuremath{\mathbb{Z}}}
\newunicodechar{ℂ}{\ensuremath{\mathbb{C}}}
\newunicodechar{ℕ}{\ensuremath{\mathbb{N}}}
\newunicodechar{ℚ}{\ensuremath{\mathbb{Q}}}
\newunicodechar{𝔻}{\ensuremath{\mathbb{D}}}
\newunicodechar{α}{\ensuremath{\alpha}}
\newunicodechar{β}{\ensuremath{\beta}}
\newunicodechar{γ}{\ensuremath{\gamma}}
\newunicodechar{δ}{\ensuremath{\delta}}
\newunicodechar{ε}{\ensuremath{\varepsilon}}
\newunicodechar{θ}{\ensuremath{\theta}}
\newunicodechar{λ}{\ensuremath{\lambda}}
\newunicodechar{μ}{\ensuremath{\mu}}
\newunicodechar{σ}{\ensuremath{\sigma}}
\newunicodechar{τ}{\ensuremath{\tau}}
\newunicodechar{φ}{\ensuremath{\varphi}}
\newunicodechar{ψ}{\ensuremath{\psi}}
\newunicodechar{ω}{\ensuremath{\omega}}
\newunicodechar{Σ}{\ensuremath{\Sigma}}
% majuscules produites par \MakeUppercase dans les titres (α-aminés -> Α)
\newunicodechar{Α}{\ensuremath{\alpha}}
\newunicodechar{Β}{\ensuremath{\beta}}
\newunicodechar{Γ}{\ensuremath{\Gamma}}
\newunicodechar{Δ}{\ensuremath{\Delta}}
\newunicodechar{Ε}{\ensuremath{\varepsilon}}
\newunicodechar{Θ}{\ensuremath{\Theta}}
\newunicodechar{Λ}{\ensuremath{\Lambda}}
\newunicodechar{Μ}{\ensuremath{\mu}}
\newunicodechar{Τ}{\ensuremath{\tau}}
\newunicodechar{Φ}{\ensuremath{\Phi}}
\newunicodechar{Ψ}{\ensuremath{\Psi}}
\newunicodechar{Ω}{\ensuremath{\Omega}}
\newunicodechar{↦}{\ensuremath{\mapsto}}
\newunicodechar{∈}{\ensuremath{\in}}
\newunicodechar{∉}{\ensuremath{\notin}}
\newunicodechar{∧}{\ensuremath{\wedge}}
\newunicodechar{⇔}{\ensuremath{\Leftrightarrow}}
\newunicodechar{⟺}{\ensuremath{\Longleftrightarrow}}
\newunicodechar{⇒}{\ensuremath{\Rightarrow}}
\newunicodechar{⟹}{\ensuremath{\Longrightarrow}}
\newunicodechar{∘}{\ensuremath{\circ}}
\newunicodechar{∪}{\ensuremath{\cup}}
\newunicodechar{∩}{\ensuremath{\cap}}
\newunicodechar{≡}{\ensuremath{\equiv}}
\newunicodechar{⊂}{\ensuremath{\subset}}
\newunicodechar{⊥}{\ensuremath{\perp}}
\newunicodechar{∃}{\ensuremath{\exists}}
\newunicodechar{∀}{\ensuremath{\forall}}
\newunicodechar{∛}{\ensuremath{\sqrt[3]{\,}}}
\newunicodechar{∜}{\ensuremath{\sqrt[4]{\,}}}
\newunicodechar{⃗}{\ensuremath{{}^{\to}}}
\newunicodechar{ⁿ}{\ifmmode^{n}\else\textsuperscript{\ensuremath{n}}\fi}
\newunicodechar{ˣ}{\ifmmode^{x}\else\textsuperscript{\ensuremath{x}}\fi}
\newunicodechar{ᵇ}{\ifmmode^{b}\else\textsuperscript{\ensuremath{b}}\fi}
\newunicodechar{ʳ}{\ifmmode^{r}\else\textsuperscript{\ensuremath{r}}\fi}
\newunicodechar{ᵘ}{\ifmmode^{u}\else\textsuperscript{\ensuremath{u}}\fi}
\newunicodechar{ʸ}{\ifmmode^{y}\else\textsuperscript{\ensuremath{y}}\fi}
\newunicodechar{ᵃ}{\ifmmode^{a}\else\textsuperscript{\ensuremath{a}}\fi}
\newunicodechar{ᵉ}{\ifmmode^{e}\else\textsuperscript{\ensuremath{e}}\fi}
\newunicodechar{ᵏ}{\ifmmode^{k}\else\textsuperscript{\ensuremath{k}}\fi}
\newunicodechar{ᵖ}{\ifmmode^{p}\else\textsuperscript{\ensuremath{p}}\fi}
\newunicodechar{ᴺ}{\ifmmode^{N}\else\textsuperscript{\ensuremath{N}}\fi}
\newunicodechar{ᶜ}{\ifmmode^{c}\else\textsuperscript{\ensuremath{c}}\fi}
\newunicodechar{ʰ}{\ifmmode^{h}\else\textsuperscript{\ensuremath{h}}\fi}
\newunicodechar{ᵠ}{\ifmmode^{q}\else\textsuperscript{\ensuremath{q}}\fi}
\newunicodechar{ₙ}{\ifmmode_{n}\else\textsubscript{\ensuremath{n}}\fi}
\newunicodechar{ₐ}{\ifmmode_{a}\else\textsubscript{\ensuremath{a}}\fi}
\newunicodechar{ᵢ}{\ifmmode_{i}\else\textsubscript{\ensuremath{i}}\fi}
\newunicodechar{ᵣ}{\ifmmode_{r}\else\textsubscript{\ensuremath{r}}\fi}
\newunicodechar{ₖ}{\ifmmode_{k}\else\textsubscript{\ensuremath{k}}\fi}
\newunicodechar{ᵦ}{\ifmmode_{\beta}\else\textsubscript{\ensuremath{\beta}}\fi}
\newunicodechar{ₓ}{\ifmmode_{x}\else\textsubscript{\ensuremath{x}}\fi}
\newfontfamily\popdisplay{ArchivoBlack-Regular.ttf}[Path=../../../fonts/]
\newfontfamily\poplogo{SpaceGrotesk-Bold.ttf}[Path=../../../fonts/]
\newfontfamily\popsemi{PlusJakartaSans-SemiBold.ttf}[Path=../../../fonts/]
\newfontfamily\popxbold{PlusJakartaSans-ExtraBold.ttf}[Path=../../../fonts/]
\newfontfamily\popxboldls{PlusJakartaSans-ExtraBold.ttf}[Path=../../../fonts/,LetterSpace=3.0]

\setlist[enumerate]{leftmargin=1.7em,itemsep=5pt,topsep=4pt}
\setlist[itemize]{leftmargin=1.7em,itemsep=4pt,topsep=4pt,label={\color{poporange}\textbullet}}
\setlength{\parindent}{0pt}
\setlength{\parskip}{5pt}
\renewcommand{\arraystretch}{1.25}

% pgfplots : style Pop par défaut
\pgfplotsset{
  every axis/.append style={axis line style={popink,line width=1pt},tick style={popink},
    grid style={popline,line width=0.5pt},label style={font=\small},tick label style={font=\footnotesize}},
  popcurve/.style={poporange,line width=1.8pt,no markers,smooth},
}
% Même style disponible dans un \draw TikZ simple (hors pgfplots)
\tikzset{popcurve/.style={poporange,line width=1.8pt}}
\tikzset{popgrid/.style={popline,very thin}}
\colorlet{popcurve}{poporange} % le nom du style est parfois utilisé comme couleur

% ============================================================
% Pill / squiggle
% ============================================================
\newcommand{\poppill}[3]{%
  \tikz[baseline=-0.55ex]\node[draw=popink,line width=1.2pt,rounded corners=2.6pt,%
    fill=#2,inner xsep=6.5pt,inner ysep=3pt]{%
    \popxboldls\fontsize{7}{7}\selectfont\color{#3}\MakeUppercase{#1}};%
}
\newcommand{\popsquiggle}[1]{%
  \tikz[baseline=-0.4ex]{
    \draw[#1,line width=2.6pt,line cap=round]
      plot [smooth] coordinates {(0,0.09) (0.34,0.01) (0.68,0.21) (1.02,0.01) (1.36,0.10)};
  }%
}
% Mot-clé surligné (marqueur orange sous le texte)
\newcommand{\cle}[1]{{\popxbold\color{popdark}#1}}

% ============================================================
% En-tête / pied
% ============================================================
\pagestyle{fancy}
\fancyhf{}
\renewcommand{\headrulewidth}{0pt}
\renewcommand{\footrulewidth}{0pt}
\lhead{\raisebox{-0.15ex}{\poplogo\fontsize{13}{13}\selectfont\color{popink}OnBuch\color{poporange}+}}
\rhead{\poppill{\DOCMATIERE\ \textbullet\ \DOCNIVEAU\ \textbullet\ Chapter \DOCLECON}{white}{popink}}
\lfoot{\popxbold\fontsize{7.6}{7.6}\selectfont\color{popmuted}\MakeUppercase{Signed course OnBuch+}\ \textbullet\ {\popsemi\DOCTITRE}}
\rfoot{\tikz[baseline=-0.6ex]\node[rounded corners=8.5pt,fill=popink,minimum height=17pt,minimum width=17pt,inner xsep=5pt,inner sep=0]{\popxbold\fontsize{8}{8}\selectfont\color{popcream}\thepage\,/\,\pageref*{LastPage}};}

% ============================================================
% Titres
% ============================================================
\newcommand{\chaptitre}[2]{%
  \begin{tcolorbox}[enhanced,colback=poporange,colframe=popink,boxrule=2.6pt,arc=11pt,
      drop shadow=popink,left=15pt,right=15pt,top=12pt,bottom=11pt,before skip=3pt,after skip=18pt]
    \poppill{#2}{popink}{popcream}\par\vspace{9pt}
    {\raggedright\hyphenpenalty=10000\exhyphenpenalty=10000\popdisplay\fontsize{22}{24}\selectfont\color{popcream}\MakeUppercase{#1}\par}
  \end{tcolorbox}
}
% Section numérotée : pastille orange avec le numéro + titre Archivo
\newcounter{popsec}
\newcommand{\coursec}[1]{%
  \stepcounter{popsec}\setcounter{popsub}{0}%
  \par\vspace{16pt}\noindent
  \begin{minipage}{\linewidth}
  \tikz[baseline=-0.7ex]\node[draw=popink,line width=1.6pt,fill=poporange,rounded corners=4pt,minimum size=22pt,inner sep=0]{\popdisplay\fontsize{12}{12}\selectfont\color{popcream}\thepopsec};%
  \hspace{8pt}{\popdisplay\fontsize{15}{17}\selectfont\color{popink}\MakeUppercase{#1}}\par
  \vspace{4pt}\noindent{\color{poporange}\rule{36pt}{3.6pt}}
  \end{minipage}\par\nopagebreak\vspace{8pt}
}
\newcounter{popsub}
\newcommand{\courssub}[1]{%
  \stepcounter{popsub}%
  \par\vspace{9pt}\noindent{\popxbold\fontsize{11.5}{13}\selectfont\color{popink}{\color{poporange}\thepopsec.\thepopsub}\hspace{6pt}#1}\par\nopagebreak\vspace{4pt}
}

% ============================================================
% Boîtes pédagogiques — même recette : fond blanc, bordure épaisse colorée,
% ombre dure encre, badge plein. Titre optionnel [#1].
% ============================================================
\tcbset{popbase/.style={breakable,enhanced,colback=white,boxrule=2.3pt,arc=9pt,
  shadow={3pt}{-3pt}{0pt}{popink},left=14pt,right=14pt,top=11pt,bottom=11pt,before skip=11pt,after skip=15pt,
  fontupper=\popsemi\fontsize{10.6}{14.5}\selectfont\color{popink},
  fontlower=\popsemi\fontsize{10.6}{14.5}\selectfont\color{popink}}}
% Coche dessinée (le glyphe ✓ n'existe pas dans les polices du document)
\newcommand{\popcheck}[1]{\tikz[baseline=-0.5ex]\draw[#1,line width=1.6pt,line cap=round,line join=round](-0.13,0.0)--(-0.04,-0.09)--(0.13,0.1);}
\providecommand{\coche}{\popcheck{popgreen}}
\providecommand{\resultat}[1]{\par\smallskip\noindent\fcolorbox{popgreen}{popgreenL}{\parbox{\dimexpr\linewidth-2\fboxsep-2\fboxrule\relax}{\textbf{Result.} #1}}\par\smallskip}
\newcommand{\popboxhead}[3]{\poppill{#1}{#2}{white}\ifx\relax#3\relax\else\hspace{6pt}{\popxbold\fontsize{10.6}{12}\selectfont\color{popink}#3}\fi\par\vspace{7pt}}

\newtcolorbox{aretenir}[1][]{popbase,colframe=popgreen,before upper={\popboxhead{Key points}{popgreen}{#1}}}
\newtcolorbox{exemplebox}[1][]{popbase,colframe=poporange,before upper={\popboxhead{Exemple}{poporange}{#1}}}
\newtcolorbox{definition}[1][]{popbase,colframe=popblue,before upper={\popboxhead{Definition}{popblue}{#1}}}
\newtcolorbox{propriete}[1][]{popbase,colframe=poppurple,before upper={\popboxhead{Property}{poppurple}{#1}}}
\newtcolorbox{methode}[1][]{popbase,colframe=popgold,before upper={\popboxhead{Method}{popgold}{#1}}}
\newtcolorbox{attention}[1][]{popbase,colframe=poppink,before upper={\popboxhead{Warning}{poppink}{#1}}}
\newtcolorbox{experience}[1][]{popbase,colframe=popblue,colback=popblueL,before upper={\popboxhead{Experiment}{popblue}{#1}}}
% « Le savais-tu ? » : carte violette pleine (contexte, culture, Cameroun)
\newtcolorbox{savaistu}[1][]{popbase,colback=poppurple,colframe=popink,
  fontupper=\popsemi\fontsize{10.4}{14.2}\selectfont\color{white},
  before upper={\poppill{Did you know?}{white}{poppurple}\ifx\relax#1\relax\else\hspace{6pt}{\popxbold\color{white}#1}\fi\par\vspace{7pt}}}
% Exercice résolu : énoncé en haut, \tcblower puis la solution sur fond crème
\newcounter{popexo}\newcounter{popexr}
\newtcolorbox{exoresolu}[1][]{popbase,colframe=poporange,colbacklower=popcream,
  segmentation style={popink,line width=1.2pt,dashed},
  before upper={\stepcounter{popexr}\popboxhead{Worked example \thepopexr}{poporange}{#1}},
  before lower={\poppill{Solution}{popink}{popcream}\par\vspace{6pt}}}
% Exercice (non résolu en place) — la correction est dans la partie Corrigés
\newtcolorbox{exercice}[1][]{popbase,colframe=popink,
  before upper={\stepcounter{popexo}\popboxhead{Exercise \thepopexo}{popink}{#1}}}
% Corrigé
\newtcolorbox{corrige}[1][]{popbase,colframe=popgreen,colback=popgreenL,
  before upper={\popboxhead{Solution}{popgreen}{#1}}}
% Objectifs / prérequis / bilan
\newtcolorbox{objectifs}{popbase,colframe=popink,colback=poporangeL,
  before upper={\popboxhead{Chapter objectives}{popink}{}}}
\newtcolorbox{prerequis}{popbase,colframe=popink,colback=popblueL,
  before upper={\popboxhead{Prerequisites}{popblue}{}}}
\newtcolorbox{bilan}{popbase,colframe=popink,colback=white,boxrule=2.8pt,
  before upper={\popboxhead{Fiche bilan}{poporange}{L'essentiel à connaître pour l'examen}}}
% Figure encadrée avec légende
\newtcolorbox{popfigure}[1][]{enhanced,breakable=false,colback=white,colframe=popline,boxrule=1.4pt,arc=8pt,
  left=10pt,right=10pt,top=10pt,bottom=8pt,before skip=10pt,after skip=12pt,halign=center,
  lower separated=false,halign lower=center,
  fontlower=\popsemi\fontsize{9}{11.5}\selectfont\color{popmuted},
  #1}
\newcommand{\legende}[1]{\par\vspace{6pt}{\popsemi\fontsize{9}{11.5}\selectfont\color{popmuted}#1}}

% ============================================================
% Signature OnBuch+
% ============================================================
\newcommand{\poplogoinline}[1]{{\poplogo\fontsize{#1}{#1}\selectfont\color{popink}OnBuch\color{poporange}+}}
\newcommand{\signatureonbuch}{%
  \begin{tcolorbox}[enhanced,colback=popink,colframe=popink,boxrule=2.2pt,arc=10pt,
      drop shadow=poporange,left=14pt,right=14pt,top=10pt,bottom=10pt,before skip=0pt,after skip=0pt]
    \tikz[baseline=-0.7ex]\node[circle,fill=popgreen,draw=popcream,line width=1.3pt,minimum size=22pt,inner sep=0]{\popcheck{white}};%
    \hspace{9pt}\begin{minipage}[c]{0.62\linewidth}
      {\popxbold\fontsize{12}{13}\selectfont\color{popcream}Signed\ \textbullet\ {\poplogo OnBuch\color{poporange}+}}\par\vspace{2pt}
      {\raggedright\popsemi\fontsize{8}{10.5}\selectfont\color{popline}\MakeUppercase{OnBuch+ teaching team}\\\MakeUppercase{Follows the official MINESEC syllabus}\par}
    \end{minipage}\hfill
    {\poplogo\fontsize{22}{22}\selectfont\color{popcream}OnBuch\color{poporange}+}
  \end{tcolorbox}%
}

% ============================================================
% Page de garde
% #1 titre · #2 matière · #3 niveau · #4 module · #5 n° leçon · #6 durée ·
% #7 description / objectifs (texte ou itemize)
% ============================================================
\newcommand{\pagegarde}[7]{%
  \thispagestyle{empty}
  \newgeometry{a4paper,margin=1.7cm,top=1.6cm,bottom=1.4cm}%
  \begin{tikzpicture}[remember picture,overlay]
    \fill[poporange!18] ([xshift=-3.2cm,yshift=-3.6cm]current page.north east) circle (4.6cm);
    \fill[poppurple!14] ([xshift=2.4cm,yshift=7.6cm]current page.south west) circle (3.8cm);
  \end{tikzpicture}%
  {\poplogo\fontsize{30}{30}\selectfont\color{popink}OnBuch\color{poporange}+}\hfill
  \raisebox{0.6ex}{\poppill{Signed course \textbullet\ Premium}{popink}{popcream}}\par
  \vspace{1pt}\popsquiggle{poppurple}\par
  \vspace{8pt}
  \begin{tcolorbox}[enhanced,colback=poporange,colframe=popink,boxrule=2.8pt,arc=13pt,
      drop shadow=popink,left=18pt,right=18pt,top=12pt,bottom=13pt,after skip=10pt]
    \poppill{#2 \textbullet\ #3}{popink}{popcream}\hspace{5pt}\poppill{#4}{white}{popink}\par\vspace{8pt}
    {\popdisplay\fontsize{14}{14}\selectfont\color{popink}CHAPTER #5}\par\vspace{4pt}
    {\raggedright\hyphenpenalty=10000\exhyphenpenalty=10000\popdisplay\fontsize{25}{27}\selectfont\color{popcream}\MakeUppercase{#1}\par}
    \vspace{8pt}
    {\popxbold\fontsize{9.5}{9.5}\selectfont\color{popink}\MakeUppercase{Suggested time: #6}}
  \end{tcolorbox}
  \begin{tcolorbox}[enhanced,colback=white,colframe=popink,boxrule=2.2pt,arc=10pt,
      drop shadow=popink,left=16pt,right=16pt,top=10pt,bottom=10pt,after skip=8pt]
    \poppill{In this chapter}{poppurple}{white}\par\vspace{5pt}
    \popsemi\fontsize{9.3}{12.6}\selectfont\color{popink}#7
  \end{tcolorbox}
  \noindent\begin{minipage}[t]{0.487\linewidth}
    \begin{tcolorbox}[enhanced,colback=popgreen,colframe=popink,boxrule=2.2pt,arc=10pt,
        drop shadow=popink,left=11pt,right=11pt,top=10pt,bottom=10pt,height=3.1cm,valign=center]
      \tcbox[colback=white,colframe=popink,boxrule=1.3pt,arc=5pt,boxsep=0pt,nobeforeafter,
        left=3pt,right=3pt,top=3pt,bottom=3pt,valign=center]{\qrcode[height=1.9cm]{https://play.google.com/store/apps/details?id=com.luvvix.onbuch&hl=en}}%
      \hspace{8pt}\begin{minipage}[c]{0.5\linewidth}
        {\popdisplay\fontsize{11}{12.5}\selectfont\color{white}\MakeUppercase{Download OnBuch}}\par\vspace{4pt}
        {\raggedright\popsemi\fontsize{8.4}{11}\selectfont\color{white}Free \textbullet\ Google Play\\AI tutor, past papers, results\par}
      \end{minipage}
    \end{tcolorbox}
  \end{minipage}\hfill
  \begin{minipage}[t]{0.487\linewidth}
    \begin{tcolorbox}[enhanced,colback=poppurple,colframe=popink,boxrule=2.2pt,arc=10pt,
        drop shadow=popink,left=11pt,right=11pt,top=10pt,bottom=10pt,height=3.1cm,valign=center]
      \tikz[baseline=-0.7ex]\node[circle,fill=white,draw=popink,line width=1.3pt,minimum size=1.6cm,inner sep=0]{\popdisplay\fontsize{24}{24}\selectfont\color{poppurple}+};%
      \hspace{8pt}\begin{minipage}[c]{0.6\linewidth}
        {\popdisplay\fontsize{11}{12.5}\selectfont\color{white}\MakeUppercase{Go OnBuch+}}\par\vspace{4pt}
        {\raggedright\popsemi\fontsize{8.4}{11}\selectfont\color{white}All signed courses, unlimited AI tutor, PDF sheets — OnBuch+ tab in the app\par}
      \end{minipage}
    \end{tcolorbox}
  \end{minipage}
  \vspace{8pt}
  \popcontenu
  \vfill
  \signatureonbuch
  \restoregeometry
  \newpage
}

% Rangée « Ce cours contient » (page de garde)
\newcommand{\poptile}[3]{% #1 couleur #2 gros texte #3 légende
  \begin{tcolorbox}[enhanced,colback=#1,colframe=popink,boxrule=2pt,arc=9pt,shadow={2.5pt}{-2.5pt}{0pt}{popink},
      left=6pt,right=6pt,top=9pt,bottom=9pt,halign=center,valign=center,height=1.75cm,width=\linewidth,nobeforeafter]
    {\popdisplay\fontsize{12.5}{13}\selectfont\color{white}\MakeUppercase{#2}}\par\vspace{3pt}
    {\centering\popsemi\fontsize{7.8}{9.5}\selectfont\color{white}#3\par}
  \end{tcolorbox}}
\newcommand{\popcontenu}{%
  \par\noindent{\popxboldls\fontsize{7.5}{7.5}\selectfont\color{popmuted}THIS SIGNED COURSE CONTAINS}\par\vspace{7pt}
  \noindent\begin{minipage}[t]{0.235\linewidth}\poptile{popblue}{Course}{complete, illustrated, step by step}\end{minipage}\hfill
  \begin{minipage}[t]{0.235\linewidth}\poptile{poporange}{Methods}{and worked examples}\end{minipage}\hfill
  \begin{minipage}[t]{0.235\linewidth}\poptile{poppink}{Exercises}{exam style + solutions}\end{minipage}\hfill
  \begin{minipage}[t]{0.235\linewidth}\poptile{popgreen}{Summary sheet}{the essentials to remember}\end{minipage}\par}

% Sommaire
\newtcolorbox{sommaire}{enhanced,colback=white,colframe=popink,boxrule=2.2pt,arc=10pt,
  drop shadow=popink,left=18pt,right=18pt,top=13pt,bottom=13pt,before skip=6pt,after skip=20pt,
  before upper={\poppill{Contents}{poppurple}{white}\par\vspace{9pt}\popsemi\fontsize{10.6}{14.5}\selectfont\color{popink}}}

% Signature de fin de cours — poussée tout en bas de la dernière page
\newcommand{\findecours}{%
  \par\vfill\nopagebreak
  \begin{center}\popsquiggle{poporange}\end{center}\vspace{4pt}
  \signatureonbuch
}
\AtBeginDocument{\def\llbracket{\mathopen{[\mkern-3.5mu[}}\def\rrbracket{\mathclose{]\mkern-3.5mu]}}} % intervalles d'entiers (glyphe ⟦ absent de la police)
% Glyphes absents de Fira Math : redéfinis à partir de glyphes disponibles
\AtBeginDocument{%
  \def\setminus{\mathbin{\backslash}}\let\smallsetminus\setminus
  \def\vdots{\mathord{\vbox{\baselineskip3.2pt\lineskiplimit0pt\kern4pt\hbox{.}\hbox{.}\hbox{.}}}}%
  \def\ddots{\mathinner{\mkern1mu\raise6.5pt\hbox{.}\mkern2mu\raise3.6pt\hbox{.}\mkern2mu\raise0.7pt\hbox{.}\mkern1mu}}%
  \def\triangle{\mathord{\scalebox{0.9}{$\symup{\Delta}$}}}%
  \def\nabla{\mathord{\raisebox{\depth}{\scalebox{1}[-1]{$\symup{\Delta}$}}}}%
  \def\checkmark{\popcheck{popdark}}%
  \def\star{\mathbin{\ast}}\def\bigstar{\mathord{\ast}}%
  \def\diamond{\mathbin{\scalebox{0.6}{$\Diamond$}}}%
  \def\bigcup{\mathop{\mathchoice{\raisebox{-0.3em}{\scalebox{1.7}{$\cup$}}}{\scalebox{1.25}{$\cup$}}{\cup}{\cup}}\displaylimits}%
  \def\bigcap{\mathop{\mathchoice{\raisebox{-0.3em}{\scalebox{1.7}{$\cap$}}}{\scalebox{1.25}{$\cap$}}{\cap}{\cap}}\displaylimits}%
  \def\lesseqqgtr{\mathrel{\lessgtr}}\def\bigoplus{\mathop{\oplus}\displaylimits}%
}
\providecommand{\ssi}{if and only if} % abréviation fréquente
\AtBeginDocument{\providecommand{\wideparen}{\overparen}} % arc de cercle

% Fonctions usuelles que les modèles utilisent sans que amsmath les fournisse
\providecommand{\arsinh}{\operatorname{arsinh}}\providecommand{\arcosh}{\operatorname{arcosh}}
\providecommand{\artanh}{\operatorname{artanh}}\providecommand{\arsech}{\operatorname{arsech}}
\providecommand{\arcsch}{\operatorname{arcsch}}\providecommand{\arcoth}{\operatorname{arcoth}}
\providecommand{\arcsinh}{\operatorname{arsinh}}\providecommand{\arccosh}{\operatorname{arcosh}}
\providecommand{\arctanh}{\operatorname{artanh}}\providecommand{\sech}{\operatorname{sech}}
\providecommand{\csch}{\operatorname{csch}}\providecommand{\arcsec}{\operatorname{arcsec}}
\providecommand{\arccsc}{\operatorname{arccsc}}\providecommand{\arccot}{\operatorname{arccot}}
\providecommand{\sgn}{\operatorname{sgn}}\providecommand{\lcm}{\operatorname{lcm}}
\providecommand{\Var}{\operatorname{Var}}\providecommand{\Cov}{\operatorname{Cov}}

%% FIGFIX-BEGIN v5 — correctifs globaux des figures (docs/onbuch-generation/tools/figfix)
\usepackage{adjustbox}\usepackage{etoolbox}\tcbuselibrary{raster}
% toute figure TikZ/pgfplots plus large que la ligne est réduite à la largeur disponible
\BeforeBeginEnvironment{tikzpicture}{\begin{adjustbox}{max width=\linewidth}}
\AfterEndEnvironment{tikzpicture}{\end{adjustbox}}
% légendes pgfplots lisibles par-dessus les courbes
\pgfplotsset{every axis/.append style={legend style={fill=white,fill opacity=0.92,text opacity=1,draw=black!15,inner sep=3pt}}}
% étiquettes d'arêtes (syntaxe edge["..."]) avec fond blanc
\tikzset{every edge quotes/.append style={fill=white,inner sep=1.2pt}}
%% FIGFIX-END
\begin{document}

\pagegarde{ENVIRONMENTAL GEOLOGY}{Geology}{Upper Sixth}{Upper Sixth — Geology}{8}{7 periods}{This chapter explores the interaction between geological processes and human activities, covering natural hazards, their mitigation, climate change impacts, and pollution. It equips students with the analytical tools to assess environmental risks and propose sustainable solutions.
\vspace{4pt}
\begin{multicols}{2}\small
\begin{itemize}
\item Define environmental geology and explain its relevance to sustainable development.
\item Identify and classify common natural geologic hazards in Cameroon and worldwide.
\item Analyse the primary and secondary effects of earthquakes, landslides, volcanic eruptions, and floods.
\item Design and evaluate mitigation strategies for each major geohazard.
\item Explain the mechanisms of global warming and its influence on climate patterns.
\item Assess the sources, pathways, and impacts of major pollutants on soil, water, and air.
\end{itemize}\end{multicols}}

\begin{sommaire}
\begin{multicols}{2}
\begin{enumerate}
\item Foundations of Environmental Geology and Major Geohazards
\item Effects and Mitigation of Geohazards – Part 1
\item Effects and Mitigation of Geohazards – Part 2
\item Global Warming and Its Effects on Climate
\item Pollution – Sources, Transport, and Management
\item Integration activity
\item Methods and worked examples
\item Exercises
\item Solutions to the exercises
\item Summary sheet
\end{enumerate}
\end{multicols}
\end{sommaire}

\begin{objectifs}
\begin{itemize}
\item Define environmental geology and explain its relevance to sustainable development.
\item Identify and classify common natural geologic hazards in Cameroon and worldwide.
\item Analyse the primary and secondary effects of earthquakes, landslides, volcanic eruptions, and floods.
\item Design and evaluate mitigation strategies for each major geohazard.
\item Explain the mechanisms of global warming and its influence on climate patterns.
\item Assess the sources, pathways, and impacts of major pollutants on soil, water, and air.
\item Integrate knowledge of hazards, climate change, and pollution in a comprehensive environmental impact assessment.
\end{itemize}
\end{objectifs}

\begin{prerequis}
\begin{itemize}
\item Basic plate tectonics and rock cycle (Lower Sixth).
\item Fundamentals of weathering, erosion, and sediment transport.
\item Understanding of the greenhouse effect and carbon cycle.
\item Basic chemistry of common pollutants (heavy metals, hydrocarbons).
\end{itemize}
\end{prerequis}

\newpage

\coursec{Foundations of Environmental Geology and Major Geohazards}

In the bustling streets of Douala, sudden flash floods turn markets into rivers within a few hours of torrential rain, while the dark volcanic slopes of Mount Cameroon remind the villages of Buea, Limbe and Idenau that eruptions have repeatedly sent lava flows down towards the sea. These everyday scenes are not accidents of fate: they are the visible expression of geological processes that shape our environment and, when they intersect with human settlement, threaten communities. The central problem of this chapter is therefore the following: \emph{how can geological knowledge be used to anticipate, reduce and manage the dangers that the Earth itself poses to people, property and ecosystems?} Answering that question is precisely the mission of environmental geology.

\courssub{Definition and Scope of Environmental Geology}

\begin{definition}[Environmental Geology]
\cle{Environmental geology} is the application of geological principles, data and methods to the identification, understanding and solution of environmental problems, so as to minimise the degradation of the environment and the loss of human life and property caused by the interaction between geological processes and human activities.
\end{definition}

The definition contains two essential ideas that candidates must be able to unpack in an examination.

\begin{itemize}
\item \textbf{Geology as a tool.} Environmental geology is not a new science separate from geology; it uses the classical tools of the geologist (field mapping, stratigraphy, petrology, structural analysis, geophysics, hydrogeology) but directs them towards practical problems: where to build, where to drill a well, where to dump waste, where a slope is unstable.
\item \textbf{Interaction between nature and society.} A lava flow in an uninhabited forest is a geological \emph{process}; the same flow crossing a road and destroying houses is an environmental \emph{problem}. Environmental geology studies precisely this interface.
\end{itemize}

The \cle{scope} of the discipline is broad. It covers, among others:

\begin{enumerate}
\item the study and mitigation of \textbf{natural geologic hazards} (earthquakes, volcanic eruptions, landslides, floods);
\item the management of \textbf{earth resources} (water, minerals, soils, energy) and the environmental consequences of their exploitation;
\item the study of \textbf{land use}, including the choice of sites for settlements, dams, roads and waste-disposal facilities;
\item the study of \textbf{environmental pollution} of soil, water and air, and of global phenomena such as climate change.
\end{enumerate}

\begin{savaistu}[Environmental geology in Cameroon]
Cameroon is sometimes called ``Africa in miniature'' for its geology too: an active volcano (Mount Cameroon), a seismically active volcanic line, unstable lateritic slopes in the humid south, flood-prone coastal lowlands around Douala, and the deadly gas-charged crater lakes of the Western Highlands (Nyos, Monoun). Few countries offer such a complete natural laboratory for environmental geology.
\end{savaistu}

\courssub{Interdisciplinary Links}

No environmental problem is purely geological. The environmental geologist works at the crossroads of several disciplines.

\begin{center}
\begin{tabularx}{\linewidth}{|l|X|}
\hline
\rowcolor{popblueL} \textbf{Partner discipline} & \textbf{Contribution to environmental geology} \\
\hline
Geology & Provides the understanding of Earth materials, structures and processes (faults, magma, slopes, aquifers). \\
\hline
Ecology & Assesses how geological events and pollution affect living organisms and ecosystems (e.g.\ the suffocation of fauna around Lake Nyos). \\
\hline
Engineering & Designs mitigation structures: retaining walls, drainage systems, seismic-resistant buildings, lava-diversion barriers. \\
\hline
Policy and land-use planning & Translates scientific hazard maps into laws, building codes, evacuation plans and protected zones. \\
\hline
\end{tabularx}
\end{center}

\begin{exemplebox}[The Lake Nyos degassing project]
After the 1986 Lake Nyos disaster, geologists identified the cause (accumulation of magmatic \ce{CO2} in deep lake water), engineers designed and installed degassing pipes that siphon gas-rich bottom water to the surface, ecologists monitored the recovery of the area, and the Cameroonian government organised the resettlement of survivors and controlled access to the lake. This is environmental geology in action: four disciplines cooperating on one problem.
\end{exemplebox}

\courssub{Classification of Natural Geologic Hazards}

\begin{definition}[Geologic hazard]
A \cle{geologic hazard} (or \cle{geohazard}) is a natural geological process or event that has the potential to cause damage to human life, property or the environment.
\end{definition}

Geohazards are classified according to the geological process involved into four main categories.

\textbf{1. Seismic hazards.}\par
These result from sudden movements of the Earth's crust along faults, generating \cle{earthquakes}. Effects include ground shaking, surface rupture, liquefaction of water-saturated soils and tsunamis. In Cameroon, the Cameroon Volcanic Line is associated with moderate seismicity, felt in the Mount Cameroon region and along the oceanic segment.

\textbf{2. Volcanic hazards.}\par
These are linked to the rise and eruption of \cle{magma}: lava flows, ash falls, pyroclastic flows, volcanic gases and lahars (volcanic mudflows). Mount Cameroon, one of Africa's most active volcanoes, erupted several times in the 20th century (notably in 1922, 1954, 1959, 1982, 1999 and 2000), with lava flows reaching the sea near Idenau in 1999.

\textbf{3. Mass movement hazards.}\par
These involve the downslope movement of rock, soil or debris under gravity: \cle{landslides}, rockfalls, mudflows and slumps. They are favoured by steep slopes, deeply weathered (lateritic) rocks, heavy rainfall and deforestation. The wet, hilly regions of the North-West, South-West and West Regions of Cameroon are particularly prone to them.

\textbf{4. Hydrological hazards.}\par
These concern the behaviour of surface and ground water: \cle{floods}, riverbank erosion, and droughts related to falling water tables. Flash floods in Douala, where intense rainfall meets a low-lying, densely built coastal plain with inadequate drainage, are a recurrent example.

\begin{popfigure}
\begin{center}
\begin{tikzpicture}[scale=1.0]
% Seismic box
\draw[fill=popblueL, draw=popblue, rounded corners=4pt, thick] (-6.4,0) rectangle (-1.7,3.4);
\node[align=center, text width=4.2cm] at (-4.05,2.9) {\textbf{Seismic}};
\draw[popdark, thick] (-5.9,1.9) -- (-5.3,2.3) -- (-4.7,1.5) -- (-4.1,2.2) -- (-3.5,1.4) -- (-2.9,2.0) -- (-2.3,1.6);
\node[align=center, text width=4.2cm, font=\small] at (-4.05,0.7) {Earthquakes, ground shaking, liquefaction};
% Volcanic box
\draw[fill=popgreenL, draw=popgreen, rounded corners=4pt, thick] (-1.2,0) rectangle (3.5,3.4);
\node[align=center, text width=4.2cm] at (1.15,2.9) {\textbf{Volcanic}};
\draw[popdark, thick, fill=poporange!40] (0.0,1.1) -- (1.15,2.4) -- (2.3,1.1) -- cycle;
\draw[poporange, thick] (1.15,2.4) -- (1.05,2.7) (1.15,2.4) -- (1.3,2.75);
\node[align=center, text width=4.2cm, font=\small] at (1.15,0.55) {Lava flows, ash falls, gases, lahars};
% Mass movement box
\draw[fill=popblueL, draw=popblue, rounded corners=4pt, thick] (4.0,0) rectangle (8.7,3.4);
\node[align=center, text width=4.2cm] at (6.35,2.9) {\textbf{Mass movement}};
\draw[popdark, thick] (4.4,1.0) -- (6.0,2.4) -- (8.3,1.0);
\draw[popgreen, thick, ->] (5.7,1.9) -- (5.1,1.35);
\fill[popgreen] (5.75,2.0) circle (0.09);
\node[align=center, text width=4.2cm, font=\small] at (6.35,0.5) {Landslides, rockfalls, mudflows};
% Hydrological box
\draw[fill=popgreenL, draw=popgreen, rounded corners=4pt, thick] (9.2,0) rectangle (13.9,3.4);
\node[align=center, text width=4.2cm] at (11.55,2.9) {\textbf{Hydrological}};
\draw[popblue, thick] (9.7,1.9) .. controls (10.1,2.2) and (10.5,1.6) .. (10.9,1.9) .. controls (11.3,2.2) and (11.7,1.6) .. (12.1,1.9) .. controls (12.5,2.2) and (12.9,1.6) .. (13.3,1.9);
\draw[popblue, thick] (9.7,1.4) .. controls (10.1,1.7) and (10.5,1.1) .. (10.9,1.4) .. controls (11.3,1.7) and (11.7,1.1) .. (12.1,1.4) .. controls (12.5,1.7) and (12.9,1.1) .. (13.3,1.4);
\node[align=center, text width=4.2cm, font=\small] at (11.55,0.55) {Floods, river erosion, drought};
% Title bar
\node[align=center] at (3.75,4.1) {\textbf{The four categories of natural geologic hazards}};
\end{tikzpicture}
\end{center}
\legende{Classification of natural geologic hazards into four categories, with the main processes in each.}
\end{popfigure}

\begin{attention}[Hazard is not risk]
A very common mistake is to confuse \cle{hazard} with \cle{risk}. The \textbf{hazard} is the natural process itself (e.g.\ an eruption of Mount Cameroon); it exists whether or not people are present. The \textbf{risk} expresses the potential losses because people and property are \emph{exposed} to the hazard. An eruption on an uninhabited slope is a hazard with almost zero risk; the same eruption above a crowded village is a high risk. Examiners regularly test this distinction.
\end{attention}

\courssub{From Hazard to Risk: Susceptibility, Vulnerability, Risk}

To manage geohazards rationally, we must quantify them. Three quantities are distinguished.

\begin{definition}[Susceptibility, vulnerability and risk]
\cle{Susceptibility} is the propensity of an area to experience a given hazard, based on its physical characteristics (slope, rock type, proximity to a fault or volcano, rainfall). \cle{Vulnerability} is the degree of damage that exposed elements (people, buildings, crops) would suffer if the event occurred. \cle{Risk} combines the probability of the hazardous event with the severity of its consequences.
\end{definition}

\begin{propriete}[The risk equation]
The risk associated with a geohazard can be expressed as
\[
\text{Risk} \;=\; \text{Probability of occurrence} \;\times\; \text{Severity of consequences},
\]
often written $R = P \times C$, where $P$ is the probability (or frequency) of the hazardous event in a given period and $C$ measures the expected losses (deaths, cost of damage). The severity of consequences itself depends on the vulnerability of the exposed population and the value of the elements at risk. This relation is used qualitatively at GCE level; no formal proof is required, but candidates must be able to apply it in simple numerical situations.
\end{propriete}

\begin{exoresolu}[Comparing two villages on the slopes of Mount Cameroon]
Village A lies on the 1999 lava-flow path, \SI{2}{km} from a recent eruptive fissure; geologists estimate the probability of a damaging lava flow there at $0.02$ per year, and a flow would destroy most houses (consequence severity rated $0.9$ on a scale from 0 to 1). Village B lies \SI{15}{km} away on older terrain, with probability $0.002$ per year and severity $0.3$. Compare the risks.
\tcblower
Apply the risk equation $R = P \times C$ to each village.
\begin{align*}
R_A &= 0.02 \times 0.9 = 0.018 \text{ per year},\\
R_B &= 0.002 \times 0.3 = 0.0006 \text{ per year}.
\end{align*}
The risk at village A is $\dfrac{0.018}{0.0006} = 30$ times greater than at village B. Conclusion: although eruptions may occur anywhere on the volcano, land-use planning should restrict dense settlement near recent flow paths, because both the probability and the severity are highest there.
\end{exoresolu}

\courssub{Case Studies from Cameroon}

\textbf{1. The eruptions of Mount Cameroon.}\par
Mount Cameroon (\SI{4095}{m}) is the most active volcano in West and Central Africa. Historical eruptions (1922, 1954, 1959, 1982, 1999, 2000) produced basaltic lava flows from flank fissures. In 1999, two flows descended the south-west flank; one crossed the Limbe--Idenau road and reached the Atlantic Ocean, cutting communications and threatening plantations. The hazards include lava flows, ash fall on crops and roofs, and toxic gases. Because the fertile volcanic soils attract dense settlement and plantations, vulnerability is high.

\textbf{2. The 1986 Lake Nyos gas burst.}\par
On 21 August 1986, a massive cloud of carbon dioxide, \ce{CO2}, was suddenly released from the deep waters of Lake Nyos, a crater lake in the North-West Region. Being denser than air, the gas flowed down valleys and suffocated about 1700 people and thousands of livestock up to several kilometres from the lake. The gas is of magmatic origin, dissolving slowly into the cold bottom water until the lake became saturated and overturned (a \emph{limnic eruption}). A similar event had killed 37 people at Lake Monoun in 1984. Mitigation now relies on degassing pipes that continuously bring gas-rich bottom water to the surface in a controlled fountain.

\textbf{3. The Douala floods.}\par
Douala lies on a low, flat coastal plain crossed by the Wouri and Dibamba rivers, with a hot, very wet climate. Intense rains, high tides, poor drainage and rapid unplanned urbanisation combine to produce recurrent flash floods; the severe floods of 2012 inundated whole neighbourhoods, markets and roads, displacing residents and damaging property. This is a hydrological hazard amplified by human vulnerability: blocked drains, construction on floodplains and removal of mangroves.

\begin{popfigure}
\begin{center}
\begin{tikzpicture}[scale=0.92]
% Stylised outline of Cameroon
\draw[thick, popdark, fill=gray!8]
  (1.2,0.4) -- (2.6,0.2) -- (4.6,0.5) -- (5.6,1.4) -- (6.4,2.6)
  -- (6.2,4.2) -- (5.2,5.4) -- (4.6,6.6) -- (3.4,7.6) -- (2.2,7.2)
  -- (1.6,6.0) -- (0.6,5.2) -- (0.2,4.2) -- (0.6,3.0) -- (0.2,2.0)
  -- (0.6,1.0) -- cycle;
% Volcanic area (Mount Cameroon / CVL)
\fill[poporange, opacity=0.55] (0.9,1.6) circle (0.55);
\draw[poporange, thick] (0.9,1.6) circle (0.55);
\node[font=\small, align=center, text width=3.2cm] at (0.9,0.55) {Volcanic area (Mt Cameroon, CVL)};
% Seismic zone along CVL
\draw[popink, very thick, decorate, decoration={zigzag, segment length=4, amplitude=1.6}] (0.5,1.2) -- (1.6,3.0) -- (2.2,4.4);
\node[font=\small, align=center, text width=3.0cm] at (3.6,3.4) {Seismic zone along the Cameroon Volcanic Line};
% Landslide-prone slopes (Western Highlands)
\fill[popgreen, opacity=0.5] (1.7,4.6) ellipse (0.7 and 0.5);
\node[font=\small, align=center, text width=3.2cm] at (1.8,5.6) {Landslide-prone slopes (Western Highlands, Lake Nyos)};
% Floodplain (Douala / Wouri estuary)
\fill[popblue, opacity=0.45] (1.5,1.0) ellipse (0.8 and 0.35);
\node[font=\small, align=center, text width=3.0cm] at (3.4,0.9) {Floodplain (Douala, Wouri estuary)};
% Lake Nyos marker
\fill[popink] (1.9,4.9) circle (0.07);
\node[font=\small] at (2.7,5.0) {L.\ Nyos};
% North arrow
\draw[->, thick] (6.0,7.0) -- (6.0,7.8);
\node[font=\small] at (6.0,8.0) {N};
\end{tikzpicture}
\end{center}
\legende{Schematic hazard map of Cameroon: volcanic area of the Cameroon Volcanic Line (orange), associated seismic zone (zigzag), landslide-prone slopes of the Western Highlands (green) and the Douala--Wouri floodplain (blue).}
\end{popfigure}

\courssub{Hazard Mapping Basics}

A \cle{hazard map} is a map that shows the spatial distribution of a hazard, usually as zones of different susceptibility. Modern hazard maps are produced with a \cle{Geographic Information System} (GIS), which superposes digital layers of information.

\begin{methode}[Producing a simple hazard susceptibility map with GIS layers]
\begin{enumerate}
\item \textbf{Define the hazard} to be mapped (e.g.\ landslides) and the study area.
\item \textbf{Collect the controlling layers}: slope steepness (from a digital elevation model), rock type and degree of weathering (geological map), rainfall, land use and vegetation cover, and the location of past events (an inventory map).
\item \textbf{Weight each layer} according to its influence on the hazard (steep slopes and weathered rocks receive high weights).
\item \textbf{Overlay the layers} in the GIS and combine the weights to compute a susceptibility score for every point of the area.
\item \textbf{Classify the scores} into zones (low, moderate, high susceptibility) and assign colours.
\item \textbf{Validate the map} by checking that past events fall mostly inside the high-susceptibility zones, then present the map to planners.
\end{enumerate}
\end{methode}

\begin{aretenir}[Key points of this section]
\begin{itemize}
\item Environmental geology applies geological knowledge to solve environmental problems and protect communities.
\item It is interdisciplinary, linking geology with ecology, engineering and public policy.
\item Natural geohazards fall into four categories: seismic, volcanic, mass movement and hydrological.
\item Hazard is the natural process; risk equals probability of occurrence multiplied by severity of consequences, and grows with human vulnerability.
\item Cameroon offers textbook examples: Mount Cameroon eruptions, the Lake Nyos gas burst of 1986 and the Douala floods.
\item Hazard susceptibility maps are built by overlaying weighted GIS layers and validated against past events.
\end{itemize}
\end{aretenir}

\begin{exercice}[Hazard, vulnerability and risk]
\begin{enumerate}
\item Define environmental geology and state two of its fields of application.
\item Distinguish clearly between a hazard and a risk, illustrating your answer with the Lake Nyos event.
\item A district on the flank of Mount Cameroon has an estimated probability of $0.01$ per year of being affected by a lava flow, with a consequence severity of $0.8$. Calculate the annual risk and explain what would happen to the risk if the population of the district doubled.
\item List the four categories of natural geologic hazards and give one Cameroonian example of each.
\end{enumerate}
\end{exercice}

\begin{corrige}[Exercise 1]
\begin{enumerate}
\item Environmental geology is the application of geological principles to the identification and solution of environmental problems. Any two of: geohazard study and mitigation, management of earth resources, land-use planning and waste-disposal siting, study of pollution and climate change.
\item The hazard is the natural process: the accumulation and sudden release of magmatic \ce{CO2} from Lake Nyos. The risk is the potential loss because people and livestock lived in the valleys downwind of the lake; the same gas burst in an uninhabited area would have been a hazard with negligible risk.
\item $R = P \times C = 0.01 \times 0.8 = 0.008$ per year. If the population doubled, the severity of consequences $C$ (more people and property exposed) would increase, so the risk would increase even though the hazard (the probability $P$) is unchanged.
\item Seismic: earthquakes along the Cameroon Volcanic Line. Volcanic: eruptions of Mount Cameroon (e.g.\ 1999). Mass movement: landslides on the steep slopes of the Western Highlands. Hydrological: flash floods in Douala (e.g.\ 2012).
\end{enumerate}
\end{corrige}

\coursec{Effects and Mitigation of Geohazards – Part 1}

In the previous section we identified the major geohazards that threaten human societies — earthquakes, volcanic eruptions, landslides, floods and coastal hazards — and we saw that a \cle{hazard} only becomes a \cle{disaster} when it strikes a vulnerable, exposed population. We now move from \emph{recognising} hazards to \emph{living with} them. This section focuses on the earthquake, the deadliest sudden-onset geohazard, and answers two questions that every GCE candidate must master: \emph{what exactly does an earthquake do to the ground and to structures?} and \emph{what can engineers and communities do to reduce the losses?} Cameroon is directly concerned: the Mount Cameroon region and the Cameroon Volcanic Line are seismically active, and the 2022 seismic swarms felt around Buea reminded us that preparedness is not an abstract luxury.

\courssub{Primary Effects of Earthquakes}

Primary effects are the consequences produced \emph{directly} by the passage of seismic waves and by fault movement itself. Three of them dominate: ground shaking, surface rupture and liquefaction.

\textbf{Ground shaking.}\par
When elastic strain energy accumulated along a fault is suddenly released, seismic waves radiate outward and set the ground in motion. The severity of shaking at a site, called the \cle{intensity}, depends on the magnitude of the event, the distance from the focus, the depth of the focus, and — critically — the \emph{nature of the local ground}. Soft, water-saturated sediments (such as the alluvial fills of the Wouri estuary around Douala) amplify shaking far more than solid bedrock, a phenomenon known as \cle{site amplification}. Engineers quantify shaking by the \cle{peak ground acceleration} (PGA), the maximum acceleration recorded at the surface, expressed in $\mathrm{m\,s^{-2}}$ or as a fraction of $g$ ($g = 9.81\ \mathrm{m\,s^{-2}}$).

\begin{propriete}[The magnitude scale is logarithmic]
The magnitude $M$ of an earthquake (Richter $M_L$ or moment magnitude $M_w$) is a logarithmic measure of wave amplitude (and hence of energy). Each increase of one magnitude unit corresponds to approximately $31.6$ times more energy released:
\[ \frac{E_2}{E_1} \approx 31.6^{\,M_2 - M_1}. \]
Thus a magnitude $7.0$ event releases about $31.6 \times 31.6 \approx 1000$ times more energy than a magnitude $5.0$ event. (The derivation from the definition of the scale is not examinable, but the factor $31.6$ and its use in calculations are.)
\end{propriete}

\begin{exemplebox}[Comparing two earthquakes]
How many times more energy does a magnitude $6.0$ earthquake release than a magnitude $4.0$ earthquake?
\[ \frac{E_6}{E_4} = 31.6^{\,6-4} = 31.6^2 \approx 998.6 \approx 1000. \]
The magnitude $6.0$ event releases roughly \textbf{one thousand times} more energy. This is why "only two points higher on the scale" is a dangerously misleading phrase.
\end{exemplebox}

\begin{popfigure}
\begin{center}
\begin{tikzpicture}
\begin{axis}[
  width=0.92\linewidth, height=5.6cm,
  xlabel={Time (s)}, ylabel={Ground acceleration (m\,s$^{-2}$)},
  xmin=0, xmax=30, ymin=-4, ymax=4,
  grid=both, grid style={popline!40},
  axis line style={popink},
  tick label style={font=\small},
  label style={font=\small},
]
\addplot[popblue, thick, domain=0:30, samples=400]
  {3.2*exp(-0.09*x)*(1-exp(-0.6*x))*sin(deg(2.2*x)) + 0.6*exp(-0.09*x)*(1-exp(-0.6*x))*sin(deg(5.7*x))};
\addplot[poporange, dashed, domain=0:30] {3.2};
\node[font=\small, text=poporange, anchor=south west] at (axis cs:21,3.2) {PGA $\approx 3.2$};
\end{axis}
\end{tikzpicture}
\end{center}
\legende{Idealised accelerogram for a magnitude $6.0$ event: acceleration rises rapidly to the peak ground acceleration (PGA), then decays as the seismic energy is dissipated.}
\end{popfigure}

\textbf{Surface rupture.}\par
If the fault break reaches the surface, the ground is physically offset — sometimes by several metres — along a line that can extend for tens of kilometres. Any structure, road, pipeline or water main built \emph{across} the rupture is torn apart, and no realistic building design can resist a direct offset. The only mitigation is \emph{avoidance}: mapping active fault traces and forbidding construction in a buffer zone on either side.

\textbf{Liquefaction.}\par
Perhaps the most insidious primary effect is liquefaction, which turns apparently solid ground into a fluid.

\begin{definition}[Liquefaction]
\cle{Liquefaction} is the loss of strength and stiffness of a saturated, loose, cohesionless soil (typically sand or silt) when earthquake shaking increases the pressure of the water in the pores between grains. The grains momentarily lose contact with one another, the soil behaves like a liquid, and it can no longer support foundations.
\end{definition}

The intuition is simple: dry sand in a bucket supports your weight because the grains press against each other. Shake the bucket violently while it is full of water and the grains are pushed apart by pressurised pore water; the sand "boils", and a heavy object placed on it sinks. In the field, liquefaction produces \emph{sand boils} (eruptions of sandy water at the surface), tilting and sinking of buildings, flotation of buried tanks and pipelines, and lateral spreading of river banks. Coastal and estuarine districts built on young, water-saturated sands — such as parts of Douala and Limbe — are the most susceptible in Cameroon.

\courssub{Secondary Effects of Earthquakes}

Secondary effects are triggered \emph{by} the shaking rather than caused directly by the seismic waves. They frequently kill more people than the shaking itself.

\begin{itemize}
\item \textbf{Tsunamis.} A sudden vertical displacement of the sea floor during a submarine earthquake displaces the entire water column above it. In the open ocean the wave is long and low, but as it enters shallow coastal water it slows and grows in height (\emph{shoaling}), arriving as a devastating surge. Cameroon has a low but non-zero exposure along the Gulf of Guinea.
\item \textbf{Landslides.} Shaking destabilises steep slopes, especially where rocks are weathered or slopes are over-steepened by road cuts. The volcanic slopes of Mount Cameroon and the escarpments of the Western Highlands (e.g.\ around Bafoussam and Dschang) are prone to earthquake- and rain-triggered slides that bury roads and villages.
\item \textbf{Fires.} Ruptured gas lines, overturned stoves and short-circuited electrical cables ignite fires at the very moment when water mains are broken and roads are blocked, so fire-fighting becomes nearly impossible. Fire following earthquake has historically destroyed more buildings than the shaking in several great disasters.
\item \textbf{Disease outbreaks.} Damage to water supply and sanitation systems contaminates drinking water; crowded emergency shelters favour the spread of cholera, dysentery and respiratory infections. In Cameroon, where cholera is endemic in some regions, this secondary effect must be taken very seriously.
\end{itemize}

\begin{attention}[Do not ignore secondary hazards]
A very common mistake — in examinations and in real planning — is to assess only the primary effect (shaking) and to forget the secondary effects. Ignoring secondary hazards leads to a systematic \emph{underestimation of total risk}. A complete answer on earthquake effects must always mention at least two secondary effects and explain the \emph{chain of causation}: shaking $\rightarrow$ slope failure / pipe rupture / sea-floor displacement $\rightarrow$ landslide / fire / tsunami $\rightarrow$ additional casualties and losses.
\end{attention}

\courssub{Seismic Building Codes and Retrofitting}

Since we cannot prevent earthquakes, we reduce casualties by ensuring that buildings \emph{do not collapse}. The guiding principle of seismic design is often summarised as: \emph{"earthquakes do not kill people; buildings do."}

\begin{definition}[Seismic building code]
A \cle{seismic building code} is a set of legally enforced design rules that specifies the minimum horizontal forces a structure must resist, the materials and detailing to be used, and the construction quality controls required, so that a building may be damaged in a strong earthquake but must not collapse.
\end{definition}

Key engineering techniques include:

\begin{itemize}
\item \textbf{Shear walls:} stiff reinforced-concrete or masonry walls, continuous from foundation to roof, that resist horizontal (lateral) forces and prevent the soft-storey collapse typical of open ground floors.
\item \textbf{Base isolation:} flexible bearings (layers of rubber and steel, or sliding pads) placed between the foundation and the structure, which lengthen the building's natural period and absorb much of the ground motion so that less energy reaches the superstructure.
\item \textbf{Ductile detailing:} closely spaced steel ties (stirrups) around columns and beams so that reinforced concrete can deform without brittle failure.
\item \textbf{Lightweight, well-tied roofs} and symmetrical, regular building plans, which avoid dangerous torsional twisting.
\end{itemize}

\begin{popfigure}
\begin{center}
\begin{tikzpicture}[scale=1.0, font=\small]
% ground
\fill[popgreenL] (-4.2,0) rectangle (4.2,-0.6);
\draw[popink] (-4.2,0) -- (4.2,0);
\node[text=popdark] at (3.3,-0.35) {ground};
% foundation
\fill[gray!50] (-1.6,0) rectangle (1.6,0.4);
\draw[popink] (-1.6,0) rectangle (1.6,0.4);
\node[text=popink] at (0,0.2) {foundation};
% base isolators
\foreach \x in {-1.1,-0.35,0.35,1.1}{
  \fill[poporange] (\x-0.18,0.4) rectangle (\x+0.18,0.75);
  \draw[popink] (\x-0.18,0.4) rectangle (\x+0.18,0.75);
}
\node[text=poporange, anchor=west] at (1.9,0.58) {base isolators};
\draw[popline] (1.3,0.58) -- (1.85,0.58);
% building frame
\draw[popink, thick] (-1.6,0.75) rectangle (1.6,4.2);
% floors
\foreach \y in {1.9,3.05}{\draw[popink] (-1.6,\y) -- (1.6,\y);}
% shear walls
\fill[popblue!35] (-1.6,0.75) rectangle (-1.15,4.2);
\fill[popblue!35] (1.15,0.75) rectangle (1.6,4.2);
\draw[popink] (-1.6,0.75) rectangle (-1.15,4.2);
\draw[popink] (1.15,0.75) rectangle (1.6,4.2);
\node[text=popblue, anchor=south, align=center] at (-2.6,4.2) {shear\\wall};
\draw[popline] (-1.4,4.2) -- (-2.2,4.35);
% columns
\draw[popink] (-0.4,0.75) -- (-0.4,4.2);
\draw[popink] (0.4,0.75) -- (0.4,4.2);
% ground motion arrows
\draw[->, poppink, very thick] (-3.6,-0.9) -- (-2.4,-0.9);
\draw[->, poppink, very thick] (-2.4,-1.15) -- (-3.6,-1.15);
\node[text=poppink] at (-3.0,-1.5) {ground motion};
% annotation
\node[text=popdark, align=center] at (0,4.6) {reinforced concrete frame};
\end{tikzpicture}
\end{center}
\legende{Cross-section of an earthquake-resistant building: base isolators decouple the structure from ground motion, while continuous shear walls carry the remaining lateral forces safely down to the foundation.}
\end{popfigure}

\textbf{Retrofitting} means strengthening an \emph{existing} vulnerable building: adding shear walls or steel bracing, jacketing columns with reinforced concrete or steel, tying walls to floors and roofs, and anchoring the structure to its foundation. For Cameroon, where most urban housing in Douala, Yaoundé and Buea was built without any seismic design, cost-effective retrofitting priorities are: enforcing simple codes for new masonry buildings (ring beams, vertical reinforcement at corners, light roofs), strengthening schools and hospitals first, and avoiding construction on loose saturated fills and on steep unstable slopes.

\courssub{Community Preparedness and Early Warning}

Technology alone cannot save lives; an informed and rehearsed population can.

\begin{itemize}
\item \textbf{Early warning systems (EWS).} Seismic stations near a fault detect the fast but harmless P-waves and automatically broadcast an alert before the slower, destructive S-waves and surface waves arrive. Depending on distance, this gives seconds to tens of seconds — enough to stop trains, shut gas valves, open fire-station doors and let people take cover ("Drop, Cover, Hold On"). Tsunami warning centres similarly use sea-floor and tide-gauge sensors.
\item \textbf{Drills.} Regular earthquake and evacuation drills in schools, workplaces and neighbourhoods turn correct behaviour into reflex. A population that has practised does not panic.
\item \textbf{Evacuation routes and assembly points.} Routes must be signposted, kept clear of obstructions, lead away from coastal (tsunami) and landslide-prone zones, and converge on open assembly areas where people can be counted and helped.
\item \textbf{Public education} on securing furniture, storing emergency water and first-aid kits, and knowing safe spots (under a sturdy table, away from windows and heavy objects).
\end{itemize}

\courssub{Post-Event Rapid Damage Assessment}

Immediately after an earthquake, authorities must know \emph{where} the damage is worst in order to direct rescue teams. The standard methodology proceeds in tiers:

\begin{enumerate}
\item \textbf{Reconnaissance (first hours):} aerial surveys, satellite imagery and drive-through inspections to map the broad pattern of damage and identify cut-off areas.
\item \textbf{Rapid visual screening (first days):} trained two-person teams inspect each building from outside (and inside where safe) using a standardised form: construction type, number of storeys, observed damage (none, slight, moderate, severe, collapsed).
\item \textbf{Safety tagging:} buildings are posted \textbf{green} (safe to occupy), \textbf{yellow} (restricted entry — repair needed) or \textbf{red} (unsafe — do not enter). This prevents aftershock casualties.
\item \textbf{Data compilation:} results are georeferenced and aggregated to produce damage maps, casualty and shelter needs estimates, and priority lists for demolition, shoring and detailed engineering evaluation.
\end{enumerate}

\courssub{Quantifying Seismic Risk: A Worked Approach}

\begin{methode}[Simplified seismic risk index]
To compare the relative risk of two districts, use the simplified index:
\[ R = \frac{\text{PGA} \times \text{Population density}}{\text{Building resilience factor}} \]
where PGA is the expected peak ground acceleration (in $\mathrm{m\,s^{-2}}$), population density is in persons per $\mathrm{km}^2$, and the resilience factor (dimensionless, $\geq 1$) summarises building quality — higher for well-engineered, code-compliant stock.
\begin{enumerate}
\item Collect or estimate PGA from a seismic hazard map.
\item Obtain population density for the district.
\item Assign a resilience factor (e.g.\ $1$ for poor unreinforced masonry, $3$ for code-designed reinforced concrete).
\item Compute $R$ for each district and rank them: the higher $R$, the higher the relative risk and the higher the priority for mitigation.
\end{enumerate}
\end{methode}

\begin{exoresolu}[Comparing seismic risk in two districts]
\textbf{Statement.} District A (dense urban quarter on soft alluvium): expected PGA $= 2.5\ \mathrm{m\,s^{-2}}$, population density $= 20\,000$ persons/km$^2$, buildings mostly unreinforced masonry (resilience factor $= 1$). District B (suburb on bedrock): PGA $= 1.0\ \mathrm{m\,s^{-2}}$, density $= 4\,000$ persons/km$^2$, mixed building stock (resilience factor $= 2$). Compute the risk index for each district and comment.
\tcblower
\textbf{Solution.}
For District A:
\[ R_A = \frac{2.5 \times 20\,000}{1} = 50\,000. \]
For District B:
\[ R_B = \frac{1.0 \times 4\,000}{2} = 2\,000. \]
Hence $R_A / R_B = 25$: District A carries a relative seismic risk about \textbf{25 times higher} than District B, even though the hazard (PGA) is only $2.5$ times greater. The difference is driven by exposure (population density) and vulnerability (low resilience). Mitigation should therefore prioritise District A: enforcement of building codes, retrofitting of schools and hospitals, drills and marked evacuation routes.
\end{exoresolu}

\begin{attention}[Common mistakes in examinations]
\begin{itemize}
\item Confusing \textbf{magnitude} (energy at the source, one value per earthquake) with \textbf{intensity} (effects at a site, varies from place to place).
\item Saying "a magnitude 6 is twice a magnitude 3": the scale is logarithmic — the energy ratio is $31.6^3 \approx 31\,600$.
\item Attributing liquefaction to "mud" or clay: it occurs in \emph{saturated loose sands and silts}, not in cohesive clays.
\item Listing mitigation measures without linking each one to the specific effect it reduces (e.g.\ base isolation $\rightarrow$ reduces shaking transmitted to the structure; evacuation routes $\rightarrow$ reduce exposure to tsunamis and fires).
\end{itemize}
\end{attention}

\begin{exercice}[Applying the concepts]
\begin{enumerate}
\item A magnitude $5.0$ earthquake is followed, one year later, by a magnitude $7.0$ event in the same region. Calculate the ratio of the energies released.
\item Explain, in three or four sentences, why a building founded on saturated sand near the Wouri estuary may suffer more damage than an identical building on bedrock, even at the same distance from the epicentre.
\item A district has expected PGA $= 3.0\ \mathrm{m\,s^{-2}}$, population density $= 12\,000$ persons/km$^2$ and resilience factor $1.5$. Compute its risk index, then state two realistic actions that would lower it, indicating which term of the formula each action changes.
\end{enumerate}
\end{exercice}

\begin{corrige}[Exercise 1]
\begin{enumerate}
\item $\dfrac{E_7}{E_5} = 31.6^{7-5} = 31.6^2 \approx 1000$. The second event releases about one thousand times more energy.
\item Two effects combine: (i) \emph{site amplification} — soft alluvium amplifies seismic waves, so the ground acceleration at the surface is larger than on bedrock; (ii) \emph{liquefaction} — shaking raises pore-water pressure in the saturated sand, which loses its strength, so foundations sink, tilt or spread laterally. The building on bedrock experiences weaker shaking on a stable foundation.
\item $R = \dfrac{3.0 \times 12\,000}{1.5} = 24\,000$. To lower it: (a) retrofit buildings and enforce seismic codes, which \emph{raises the resilience factor} (denominator); (b) reduce exposure by land-use planning — relocating critical facilities and limiting settlement density in the most hazardous zones — which \emph{lowers the population density term}. (Strengthening soils or avoiding liquefiable ground effectively reduces the site PGA.)
\end{enumerate}
\end{corrige}

\begin{aretenir}[Key points of this section]
\begin{itemize}
\item Primary earthquake effects: \cle{ground shaking} (measured by PGA), \cle{surface rupture} and \cle{liquefaction} of saturated loose soils.
\item Secondary effects — \cle{tsunamis}, \cle{landslides}, \cle{fires} and \cle{disease outbreaks} — often exceed primary losses and must never be omitted from a risk assessment.
\item The magnitude scale is logarithmic: $+1$ unit $\approx 31.6\times$ energy.
\item Mitigation combines \emph{engineering} (seismic codes, shear walls, base isolation, retrofitting) and \emph{community preparedness} (early warning, drills, evacuation routes), followed by structured \emph{rapid damage assessment} (reconnaissance, screening, green/yellow/red tagging).
\item Relative risk grows with hazard (PGA) and exposure (population density) and falls with building resilience: $R = \dfrac{\text{PGA} \times \text{density}}{\text{resilience}}$.
\end{itemize}
\end{aretenir}

In the next section we extend this mitigation logic to the other major geohazards — volcanic eruptions, landslides and floods — and examine how land-use planning and hazard mapping complete the engineer's and planner's toolkit.

\coursec{Effects and Mitigation of Geohazards – Part 2}

In the previous section we examined the effects of earthquakes and tsunamis and the first family of mitigation strategies. We now turn to three other geohazards that dominate the Cameroonian risk landscape: \cle{volcanic hazards} along the Cameroon Volcanic Line, \cle{landslides} in the western highlands, and \cle{floods} in the lowlands and coastal plain. For each, we proceed in the same way: observe the phenomenon, understand its mechanism, quantify the danger, and finally distinguish \textbf{structural} (engineering) from \textbf{non-structural} (planning and social) mitigation.

\courssub{Volcanic Hazards}

\textbf{Observation.} Mount Cameroon (\SI{4095}{m}), one of Africa's most active volcanoes, erupted in 1999 and 2000, sending lava flows toward the Atlantic coast near Bakingili. Yet lava was not the only threat: ash fell on plantations, and heavy rain on loose ash could have triggered mudflows. A volcano threatens its surroundings through \emph{several distinct mechanisms}, each with its own speed, reach and lethality.

\begin{definition}[Volcanic hazard]
A \cle{volcanic hazard} is any potentially damaging process associated with a volcano, whether during an eruption (lava flows, pyroclastic density currents, ash fall) or between eruptions (lahars, gas emissions, flank instability).
\end{definition}

\textbf{The four main products and their behaviour.}

\begin{itemize}
\item \textbf{Lava flows.} Streams of molten rock (\SI{700}{\ensuremath{{}^{\circ}\mathrm{C}}} to \SI{1200}{\ensuremath{{}^{\circ}\mathrm{C}}}). Basaltic flows, like those of Mount Cameroon, move slowly (typically a few \si{m.h^{-1}} to a few tens of \si{km.h^{-1}} on steep slopes). They destroy everything by burial and fire but rarely kill, because people can walk away. Mitigation: diversion barriers, cooling with water, and above all \emph{zoning}.
\item \textbf{Pyroclastic density currents (PDCs).} Ground-hugging mixtures of hot gas and volcanic fragments travelling at $100$--$700\ \mathrm{km\,h^{-1}}$ at temperatures of several hundred degrees. They are the \textbf{deadliest} volcanic phenomenon; survival inside a PDC is essentially impossible. Only evacuation before the eruption saves lives.
\item \textbf{Ash fall (tephra).} Fine fragments ejected into the atmosphere and carried by wind. Ash collapses roofs (wet ash can exceed \SI{200}{kg.m^{-2}} of load), contaminates water supplies, damages crops and aircraft engines, and causes respiratory illness.
\item \textbf{Lahars.} See the definition below — they are the hazard most often underestimated by students.
\end{itemize}

\begin{definition}[Lahar]
A \cle{lahar} is a volcanic \textbf{mudflow or debris flow}: a slurry of volcanic ash, rock fragments and water that rushes down river valleys. It is mobilised by water — from heavy rainfall, melting of summit ice or snow, or the breaching of a crater lake — and can occur \emph{long after} an eruption, even on a dormant volcano. Lahars behave like wet concrete: fast (up to $\sim 80\ \mathrm{km\,h^{-1}}$), dense, and able to carry boulders and bury entire settlements.
\end{definition}

\begin{savaistu}[Cameroon and volcanic risk]
The 1986 Lake Nyos disaster, though a \emph{gas} (limnic) event rather than an eruption, reminds us that the Cameroon Volcanic Line can kill without any lava at all: about 1\,700 people died from a cloud of magmatic \ce{CO2}. Degassing pipes now continuously remove \ce{CO2} from the lake bottom — a fine example of structural mitigation.
\end{savaistu}

\textbf{Hazard zonation.} Because each product follows a predictable path (lava and lahars follow valleys; ash follows the wind; PDCs follow topography), geologists draw \cle{volcanic hazard maps} dividing the surroundings of a volcano into zones of decreasing danger.

\begin{popfigure}
\begin{center}
\begin{tikzpicture}[scale=0.9]
% volcano cone
\fill[gray!60] (0,0) -- (2.2,4.2) -- (4.4,0) -- cycle;
\fill[popdark] (1.9,4.2) -- (2.2,3.7) -- (2.5,4.2) -- cycle;
% smoke
\draw[gray, thick, ->] (2.2,4.3) .. controls (2.8,5.0) and (3.4,5.2) .. (4.2,5.4);
\node[gray, font=\small] at (4.6,5.5) {ash plume};
% red zone
\fill[red!55, opacity=0.75] (2.2,0) ellipse (1.5 and 0.55);
% orange zone
\fill[orange!50, opacity=0.7] (2.2,0) ellipse (2.6 and 1.0);
% yellow zone
\fill[yellow!55, opacity=0.65] (2.2,0) ellipse (4.0 and 1.5);
% valley / lahar channel
\draw[popblue, very thick, ->] (2.2,0.1) -- (3.4,-0.4) -- (5.2,-0.9) -- (6.6,-1.3);
\node[popblue, font=\small, align=center] at (5.6,-1.7) {lahar / lava channel\\(follows valley)};
% legend
\fill[red!55] (6.8,1.6) rectangle (7.3,1.9); \node[right, font=\small] at (7.35,1.75) {Red: PDC, proximal lava};
\fill[orange!50] (6.8,1.0) rectangle (7.3,1.3); \node[right, font=\small] at (7.35,1.15) {Orange: lava, heavy ash};
\fill[yellow!55] (6.8,0.4) rectangle (7.3,0.7); \node[right, font=\small] at (7.35,0.55) {Yellow: ash fall, lahars};
\node[font=\small\itshape] at (2.2,-2.2) {Wind direction controls the ash-fall sector (downwind elongation).};
\end{tikzpicture}
\end{center}
\legende{Schematic volcanic hazard zonation map. Danger decreases outward, but valleys (lahar/lava paths) and the downwind sector (ash) remain hazardous far from the vent.}
\end{popfigure}

\begin{attention}[Ash fall is not ``harmless dust'']
A common mistake is to rank hazards only by temperature. Ash fall kills indirectly — roof collapse, water contamination, crop failure — and affects areas \emph{hundreds of kilometres} downwind, far beyond the red zone. Never omit it from a hazard map.
\end{attention}

\courssub{Landslides: Triggers and Classification}

\textbf{Observation.} After intense rainy-season storms, the steep escarpments of the Bamenda highlands and the slopes around Bafoussam and Dschang frequently fail, burying roads and houses. Why do some slopes stand for centuries and then fail in a single night?

\begin{definition}[Landslide]
A \cle{landslide} is the downslope movement of rock, soil or debris under the direct influence of gravity. It occurs when the \textbf{driving forces} (gravity acting on the slope mass) exceed the \textbf{resisting forces} (cohesion and friction along the potential failure surface).
\end{definition}

\textbf{Triggers} (events that push a marginally stable slope to failure):
\begin{itemize}
\item \textbf{Rainfall} — the dominant trigger in Cameroon. Water increases the weight of the soil and raises \cle{pore-water pressure}, which pushes grains apart and reduces frictional resistance.
\item \textbf{Earthquakes} — shaking momentarily destroys cohesion and liquefies saturated sands.
\item \textbf{Human excavation} — road cuts and quarrying remove support at the \emph{toe} of the slope; overloading at the crest (buildings, waste dumps) adds driving force; deforestation removes root reinforcement.
\end{itemize}

\textbf{Classification by movement type} (Varnes system, simplified):
\begin{itemize}
\item \textbf{Falls}: free-fall of rock from a cliff.
\item \textbf{Topples}: forward rotation of blocks about a pivot.
\item \textbf{Slides}: movement along a discrete surface — \emph{rotational} (slumps, curved surface) or \emph{translational} (planar surface).
\item \textbf{Flows}: fluid-like movement of saturated material (earthflows, debris flows, mudflows).
\item \textbf{Creep}: slow, continuous, millimetre-per-year downslope movement.
\end{itemize}

\begin{methode}[Factor of Safety (FoS) for slope stability — planar failure]
The \cle{Factor of Safety} quantifies how close a slope is to failure for a planar (translational) slide:
\[ \mathrm{FoS} = \frac{\text{resisting forces}}{\text{driving forces}} = \frac{c\,L + (W\cos\theta - u\,L)\tan\varphi}{W\sin\theta} \]
where $c$ = cohesion (\si{kPa}), $L$ = length of the failure surface, $W$ = weight of the sliding mass, $\theta$ = slope angle, $\varphi$ = angle of internal friction, $u$ = pore-water pressure (\si{kPa}) acting on the failure plane.
\begin{enumerate}
\item Identify the potential failure surface and measure $\theta$.
\item Estimate $W$, $c$, $\varphi$ and $u$ from field and laboratory tests.
\item Compute resisting forces (numerator) and driving forces (denominator).
\item Interpret: $\mathrm{FoS} > 1$ stable; $\mathrm{FoS} = 1$ limit equilibrium (failure imminent); $\mathrm{FoS} < 1$ unstable. Engineers typically design for $\mathrm{FoS} \geq 1.5$.
\end{enumerate}
\emph{Note how rainfall enters the calculation: pore-water pressure $u$ reduces the effective normal stress $(W\cos\theta - uL)$, shrinking the numerator — hence rain-triggered landslides.}
\end{methode}

\begin{exoresolu}[FoS of a Bamenda hillslope]
A soil slab on a $30^{\circ}$ slope weighs $W = 500\ \mathrm{kN}$ per metre width. The failure surface is $L = 10\ \mathrm{m}$ long; cohesion $c = 5\ \mathrm{kPa}$; friction angle $\varphi = 25^{\circ}$. Assume initially dry conditions ($u=0$). Compute the FoS and comment. ($\sin 30^{\circ} = 0.500$, $\cos 30^{\circ} = 0.866$, $\tan 25^{\circ} = 0.466$.)
\tcblower
\textbf{Solution.}
\begin{align*}
\text{Resisting} &= cL + W\cos\theta\tan\varphi = (5)(10) + (500)(0.866)(0.466) \\
&= 50 + 201.8 = 251.8\ \mathrm{kN}\\
\text{Driving} &= W\sin\theta = (500)(0.500) = 250\ \mathrm{kN}\\
\mathrm{FoS} &= \frac{251.8}{250} \approx 1.01
\end{align*}
\textbf{Interpretation:} $\mathrm{FoS} \approx 1$ — the slope is at \emph{limit equilibrium}. Any additional pore-water pressure from a heavy storm (which reduces the effective $W\cos\theta$ term) will push FoS below 1 and trigger failure. Mitigation (drainage, regrading, retaining structure) is urgent.
\end{exoresolu}

\textbf{Rainfall thresholds.} Empirically, landslides begin when storm intensity $I$ and duration $D$ exceed a threshold curve of the form $I = a\,D^{-b}$: short, violent storms and long, moderate rains are both dangerous. The curve below illustrates such a threshold calibrated for the Bamenda highlands.

\begin{popfigure}
\begin{center}
\begin{tikzpicture}
\begin{axis}[
    width=0.8\linewidth, height=6.2cm,
    xlabel={Rainfall duration $D$ (h)},
    ylabel={Intensity $I$ (mm/h)},
    xmin=0, xmax=48, ymin=0, ymax=40,
    grid=both, grid style={gray!20},
    legend style={at={(0.97,0.95)}, anchor=north east, font=\small},
]
\addplot[domain=2:48, samples=100, very thick, poporange, name path=threshold] {60*x^(-0.6)};
\addlegendentry{Threshold $I = 60\{,}D^{-0.6}$}
\addplot[domain=2:48, samples=2, dashed, popgreen, name path=zero] {0};
\addplot[popgreenL, opacity=0.3] fill between[of=threshold and zero];
\node[font=\small, popgreen!60!black] at (axis cs:30,6) {stable};
\node[font=\small, poporange!70!black] at (axis cs:12,26) {landslide initiation};
\end{axis}
\end{tikzpicture}
\end{center}
\legende{Rainfall intensity--duration threshold for landslide initiation in the Bamenda highlands (schematic). Storms plotting above the curve are likely to trigger landslides; such curves underpin early-warning systems.}
\end{popfigure}

\courssub{Floods: Types and the Role of Land Use}

\textbf{Observation.} Yaoundé's low-lying quarters flood after single violent storms, while the Logone plain in the Far North floods slowly when rivers overflow for weeks. Same word — \emph{flood} — but different mechanisms, different warnings, different defences.

\begin{definition}[Flood]
A \cle{flood} is the temporary covering by water of land that is normally dry, occurring when the water input exceeds the capacity of the channel or ground to convey or absorb it.
\end{definition}

\textbf{Three main types:}
\begin{itemize}
\item \textbf{Flash floods}: sudden, violent rises (minutes to hours) in small steep catchments after intense rain. Little warning; high velocities; deadly (e.g.\ urban floods in Yaoundé and Douala).
\item \textbf{Riverine (fluvial) floods}: slow-onset overflow of large rivers onto their floodplains, lasting days to weeks (e.g.\ Logone and Chari, Benue basin). Predictable but extensive.
\item \textbf{Coastal floods}: driven by storm surges, high tides and, increasingly, sea-level rise (e.g.\ low-lying areas around Limbe and Douala).
\end{itemize}

\textbf{Role of land-use change.} Deforestation, urbanisation and drainage of wetlands all \emph{accelerate} runoff: impermeable surfaces (roofs, tarmac) prevent infiltration, so a larger fraction of rainfall reaches the river, faster. The result is a \textbf{higher, sharper flood peak} — a flood of a given size becomes more frequent.

\begin{propriete}[Flood recurrence interval]
The \cle{recurrence interval} (return period) $T$ of a flood is the inverse of its annual exceedance probability $P$:
\[ T = \frac{1}{P} \]
A ``100-year flood'' has $P = 0.01$, i.e.\ a 1\,\% chance of being equalled or exceeded \emph{in any given year}. (Statement and application are examinable; no proof required.)
\end{propriete}

\begin{attention}[``100-year flood'' does not mean ``once per century'']
The probability of at least one exceedance in 100 years is $1-(1-0.01)^{100} \approx 63\,\%$. Two ``100-year floods'' can occur in consecutive years. Candidates regularly lose marks by interpreting $T$ as a guaranteed spacing.
\end{attention}

\courssub{Mitigation: Structural and Non-Structural Measures}

\textbf{Structural (engineering) measures} physically control the hazard:
\begin{itemize}
\item \textbf{Check dams}: small barriers across gullies that trap sediment and slow debris flows and lahars.
\item \textbf{Retention (detention) basins}: reservoirs that store floodwater temporarily and release it slowly, flattening the flood peak.
\item \textbf{Slope reinforcement}: retaining walls, gabions, rock bolts, terracing, drainage trenches (lowering pore-water pressure directly raises FoS), and reforestation.
\item \textbf{Volcanic monitoring networks}: seismometers, GPS tiltmeters, gas sensors and satellite imagery detect magma movement and allow timely evacuation — the only effective defence against PDCs.
\end{itemize}

\begin{attention}[Over-reliance on structural measures]
Levees and dams create a false sense of security — the ``levee effect'': people build on the protected floodplain, so when an exceptional event exceeds the design capacity, losses are \emph{catastrophic}. Structural measures also transfer risk downstream (a confined river flows faster and higher downstream). Always combine them with non-structural measures.
\end{attention}

\textbf{Non-structural measures} reduce exposure and vulnerability rather than the hazard itself:
\begin{itemize}
\item \textbf{Land-use zoning}: hazard maps legally prohibit construction in red zones, floodplains and unstable slopes; floodplains are reserved for agriculture or parks.
\item \textbf{Insurance}: spreads financial losses across society and encourages risk-aware building.
\item \textbf{Public education and early warning}: evacuation drills, hazard education in schools, radio alerts (crucial for flash floods and lahars, where minutes matter).
\end{itemize}

\begin{exercice}[Choosing a mitigation strategy]
A growing town lies on a debris fan at the foot of a steep, deforested escarpment crossed by a single gully. Propose a combined mitigation plan, justifying each measure as structural or non-structural.
\end{exercice}

\begin{corrige}[Exercise 1]
\textbf{Structural:} (i) a check dam in the gully to trap debris-flow sediment; (ii) reforestation and drainage trenches on the escarpment to raise FoS by reducing pore-water pressure; (iii) a retention basin above the town to flatten flood peaks.
\textbf{Non-structural:} (iv) zoning the fan apex as a no-build area using a hazard map; (v) a rainfall early-warning system based on an intensity--duration threshold curve; (vi) public education and evacuation drills; (vii) insurance schemes for remaining residents.
\textbf{Justification:} structural measures alone would encourage further settlement (levee effect) and could fail in an exceptional storm; non-structural measures reduce exposure even if engineering defences are overtopped.
\end{corrige}

\begin{aretenir}[Key points]
\begin{itemize}
\item Volcanic danger comes from four products: lava flows (slow, destructive), PDCs (fast, lethal), ash fall (widespread, indirect) and lahars (water-mobilised mudflows, possible without eruption).
\item Landslides fail when driving forces exceed resisting forces; $\mathrm{FoS} = \text{resisting}/\text{driving}$, with failure at $\mathrm{FoS} \leq 1$. Rainfall acts through pore-water pressure.
\item Floods are flash, riverine or coastal; land-use change sharpens flood peaks. Recurrence interval $T = 1/P$.
\item Best practice combines structural measures (check dams, basins, slope reinforcement, monitoring) with non-structural ones (zoning, insurance, education) — never engineering alone.
\end{itemize}
\end{aretenir}

\coursec{Global Warming and Its Effects on Climate}

Having examined the immediate dangers posed by earthquakes, landslides, and volcanic eruptions in the previous sections, we now turn to a geohazard of a different nature: one that operates on a planetary scale, unfolds over decades, and amplifies almost every other geological risk. \emph{Global warming} — the long-term rise in Earth's average surface temperature due to human-induced increases in greenhouse gas concentrations — is the defining environmental challenge of our time. For a geologist, it is not merely an atmospheric phenomenon; it is a perturbation of the Earth system that alters the hydrological cycle, accelerates coastal erosion, destabilizes slopes, and even influences volcanic activity through ice unloading. Understanding its mechanisms, its observed fingerprints in Cameroon, and the strategies to mitigate and adapt to it is the essence of this section.

\courssub{The Greenhouse Effect, Radiative Forcing, and Energy Balance}

The Earth's temperature is governed by a simple budget: energy in (solar radiation) must equal energy out (terrestrial infrared radiation) for the climate to be in equilibrium. The atmosphere, however, contains trace gases that are transparent to incoming shortwave solar radiation but absorb and re-emit outgoing longwave infrared radiation. This is the natural \cle{greenhouse effect}, without which the planet's average temperature would be about \SI{-18}{\ensuremath{{}^{\circ}\mathrm{C}}} instead of the current \SI{+15}{\ensuremath{{}^{\circ}\mathrm{C}}}.

Human activities since the Industrial Revolution have increased the concentrations of these greenhouse gases (GHGs), notably carbon dioxide (\ce{CO2}), methane (\ce{CH4}), and nitrous oxide (\ce{N2O}). This enhances the greenhouse effect, reducing the outgoing energy flux and creating a positive \cle{radiative forcing}.

\begin{definition}[Radiative Forcing]
\textbf{Radiative forcing} ($\Delta F$) is the change in the net radiative flux (downward minus upward) at the tropopause or top of the atmosphere due to an external driver (e.g., change in \ce{CO2} concentration, solar irradiance, aerosols), expressed in watts per square metre (\SI{}{W.m^{-2}}). Positive forcing warms the system; negative forcing cools it.
\end{definition}

The primary anthropogenic forcing agents are well-mixed greenhouse gases. For \ce{CO2}, the radiative forcing can be approximated by the Myhre et al. (1998) formula:
\[
\Delta F_{\ce{CO2}} = 5.35 \times \ln\left(\frac{C}{C_0}\right)\ \mathrm{W\,m^{-2}}
\]
where $C$ is the current concentration and $C_0$ the pre-industrial reference (\SI{280}{ppm}). For a doubling of \ce{CO2} (\SI{560}{ppm}), $\Delta F \approx \SI{3.7}{W.m^{-2}}$.

\begin{propriete}[Climate Sensitivity Parameter]
The equilibrium global mean surface temperature change ($\Delta T_{\mathrm{eq}}$) for a given radiative forcing is:
\[
\Delta T_{\mathrm{eq}} = \lambda \, \Delta F
\]
where $\lambda$ (lambda) is the \textbf{climate sensitivity parameter}. For fast feedbacks (water vapour, lapse rate, clouds, sea ice), the IPCC assesses $\lambda \approx \SI{0.8}{\ensuremath{{}^{\circ}\mathrm{C}}/(W.m^{-2})}$ (likely range \num{0.5}--\SI{1.2}{\ensuremath{{}^{\circ}\mathrm{C}}/(W.m^{-2})}). This corresponds to an Equilibrium Climate Sensitivity (ECS) of about \SI{3}{\ensuremath{{}^{\circ}\mathrm{C}}} for \ce{CO2} doubling.
\end{propriete}

\begin{methode}[Simple Energy-Balance Model to Estimate $\Delta T$]
\textbf{Step 1:} Identify the forcing agent and its concentration change.\\
\textbf{Step 2:} Compute the radiative forcing $\Delta F$ using the appropriate formula (e.g., for \ce{CO2}: $5.35 \ln(C/C_0)$).\\
\textbf{Step 3:} Choose a climate sensitivity parameter $\lambda$ (use \SI{0.8}{\ensuremath{{}^{\circ}\mathrm{C}}/(W.m^{-2})} for fast feedbacks unless otherwise specified).\\
\textbf{Step 4:} Calculate equilibrium temperature change: $\Delta T_{\mathrm{eq}} = \lambda \times \Delta F$.\\
\textbf{Step 5:} Recognize that the \emph{transient} warming observed today is smaller than $\Delta T_{\mathrm{eq}}$ because of ocean thermal inertia (the ``warming in the pipeline'').
\end{methode}

\begin{exemplebox}[Estimating Warming from \ce{CO2} Increase]
Atmospheric \ce{CO2} has risen from \SI{280}{ppm} (pre-industrial) to \SI{420}{ppm} (2023).
\begin{enumerate}
\item $\Delta F = 5.35 \times \ln(420/280) = 5.35 \times \ln(1.5) \approx 5.35 \times 0.405 = \SI{2.17}{W.m^{-2}}$.
\item Using $\lambda = \SI{0.8}{\ensuremath{{}^{\circ}\mathrm{C}}/(W.m^{-2})}$: $\Delta T_{\mathrm{eq}} = 0.8 \times 2.17 \approx \SI{1.74}{\ensuremath{{}^{\circ}\mathrm{C}}}$.
\end{enumerate}
The observed warming to date is about \SI{1.2}{\ensuremath{{}^{\circ}\mathrm{C}}}, consistent with the transient response being lower than the equilibrium value due to ocean heat uptake.
\end{exemplebox}

\courssub{The Carbon Cycle and Anthropogenic Perturbation}

Carbon moves between reservoirs (atmosphere, ocean, land biosphere, sediments) through natural fluxes that were roughly in balance before industrialization. Human activities have added two major fluxes: fossil fuel combustion and land-use change (deforestation). About half of the emitted \ce{CO2} remains in the atmosphere; the rest is taken up by the ocean (causing acidification) and the land biosphere (\ce{CO2} fertilization).

\begin{popfigure}
\centering
\begin{tikzpicture}[node distance=1.5cm, >=stealth, font=\small, every node/.style={align=center}]
% Reservoirs
\node[draw, rounded corners, fill=popblueL, minimum width=2.5cm, minimum height=1.2cm, align=center] (atm) at (0,3){Atmosphere\\ \SI{875}{GtC}};
\node[draw, rounded corners, fill=popgreenL, minimum width=3cm, minimum height=1.5cm, align=center] (land) at (-4,0){Land Biosphere\\ \SI{2500}{GtC}};
\node[draw, rounded corners, fill=popblue, minimum width=3cm, minimum height=1.5cm, align=center] (ocean) at (4,0){Surface Ocean\\ \SI{900}{GtC}};
\node[draw, rounded corners, fill=popmuted, minimum width=3cm, minimum height=1cm, align=center] (deep) at (4,-2.5){Deep Ocean\\ \SI{37000}{GtC}};
\node[draw, rounded corners, fill=poporange, minimum width=3cm, minimum height=1cm, align=center] (fossil) at (-4,-2.5){Fossil Fuels\\ \SI{4000}{GtC}};

% Natural fluxes (black arrows)
\draw[->, thick] (land) -- node[left] {\SI{120}{GtC/yr} Photosynthesis} (atm);
\draw[->, thick] (atm) -- node[right] {\SI{120}{GtC/yr} Respiration} (land);
\draw[->, thick] (ocean) -- node[right] {\SI{90}{GtC/yr} Outgassing} (atm);
\draw[->, thick] (atm) -- node[left] {\SI{92}{GtC/yr} Dissolution} (ocean);
\draw[->, thick] (ocean) -- node[right] {\SI{10}{GtC/yr} Mixing} (deep);
\draw[->, thick] (deep) -- node[left] {\SI{10}{GtC/yr} Upwelling} (ocean);

% Anthropogenic fluxes (red arrows, dashed)
\draw[->, thick, red, dashed] (fossil) -- node[left] {\SI{9.5}{GtC/yr} Combustion} (atm);
\draw[->, thick, red, dashed] (land) -- node[left] {\SI{1.5}{GtC/yr} Land-use change} (atm);
\draw[->, thick, red] (atm) -- node[above] {\SI{5.0}{GtC/yr} Accumulation} ++(0,1.5);
\draw[->, thick, red] (atm) -- node[above right] {\SI{2.5}{GtC/yr} Ocean uptake} (ocean);
\draw[->, thick, red] (atm) -- node[above left] {\SI{3.0}{GtC/yr} Land sink} (land);

% Legend
\node[anchor=west, font=\footnotesize] at (-6,3.5) {\textcolor{black}{Natural fluxes} \tikz\draw[thick,->] (0,0) -- (0.5,0);};
\node[anchor=west, font=\footnotesize] at (-6,3.0) {\textcolor{red}{Anthropogenic fluxes} \tikz\draw[thick,red,dashed,->] (0,0) -- (0.5,0);};
\end{tikzpicture}
\legende{Simplified global carbon cycle (reservoir sizes in GtC, fluxes in GtC/yr). Red dashed arrows show anthropogenic perturbations (fossil fuel combustion \SI{9.5}{GtC/yr}, land-use change \SI{1.5}{GtC/yr}). About \SI{5}{GtC/yr} accumulates in the atmosphere, driving global warming.}
\end{popfigure}

\courssub{Observed Temperature Trends in Cameroon (1960–2020)}

Global warming is not a future projection; it is an observed reality. In Cameroon, long-term meteorological records show a clear warming trend superimposed on natural variability. The graph below illustrates the annual mean temperature anomaly (difference from the 1961–1990 baseline) for Yaoundé, the capital, located in the southern plateau region.

\begin{popfigure}
\centering
\begin{tikzpicture}
\begin{axis}[
    width=\linewidth,
    height=8cm,
    xlabel={Year},
    ylabel={Temperature Anomaly (\si{\ensuremath{{}^{\circ}\mathrm{C}}})},
    title={Annual Mean Temperature Anomaly, Yaoundé (1960–2020)},
    grid=major,
    grid style={dashed, gray!50},
    xmin=1960, xmax=2020,
    ymin=-0.8, ymax=1.5,
    xtick={1960,1970,1980,1990,2000,2010,2020},
    ytick={-0.5,0,0.5,1.0,1.5},
    legend pos=north west,
    legend cell align=left,
]
% Simulated data based on observed trends ~0.15°C/decade
\addplot[color=popblue, thick, mark=none, smooth] coordinates {
(1960,-0.35) (1961,-0.28) (1962,-0.32) (1963,-0.25) (1964,-0.30) (1965,-0.22)
(1966,-0.18) (1967,-0.25) (1968,-0.15) (1969,-0.20) (1970,-0.10)
(1971,-0.05) (1972,-0.12) (1973,-0.02) (1974,-0.08) (1975,0.00)
(1976,0.05) (1977,-0.03) (1978,0.02) (1979,0.10) (1980,0.08)
(1981,0.15) (1982,0.12) (1983,0.20) (1984,0.18) (1985,0.25)
(1986,0.22) (1987,0.30) (1988,0.28) (1989,0.35) (1990,0.32)
(1991,0.40) (1992,0.38) (1993,0.45) (1994,0.42) (1995,0.50)
(1996,0.48) (1997,0.55) (1998,0.65) (1999,0.52) (2000,0.55)
(2001,0.60) (2002,0.62) (2003,0.68) (2004,0.65) (2005,0.72)
(2006,0.70) (2007,0.78) (2008,0.75) (2009,0.82) (2010,0.90)
(2011,0.88) (2012,0.95) (2013,0.92) (2014,1.00) (2015,1.05)
(2016,1.10) (2017,1.08) (2018,1.15) (2019,1.12) (2020,1.20)
};
\addlegendentry{Annual anomaly}

% Linear trend line
\addplot[color=poporange, thick, dashed, mark=none, domain=1960:2020] {0.0155*x - 30.7};
\addlegendentry{Linear trend (\SI{0.155}{\ensuremath{{}^{\circ}\mathrm{C}}/decade})}

% Baseline
\addplot[color=black, thin, dashed] coordinates {(1960,0) (2020,0)};
\end{axis}
\end{tikzpicture}
\legende{Simulated annual mean temperature anomalies for Yaoundé (1960–2020) relative to the 1961–1990 climatology. The linear trend shows a warming of approximately \SI{0.155}{\ensuremath{{}^{\circ}\mathrm{C}}} per decade, consistent with the global average. Data source: adapted from World Bank Climate Change Knowledge Portal and national meteorological records.}
\end{popfigure}

\begin{aretenir}[Key Observed Changes in Cameroon]
\begin{itemize}
\item Mean temperature has risen by \num{0.8}--\SI{1.0}{\ensuremath{{}^{\circ}\mathrm{C}}} since 1960.
\item The warming is more pronounced in the northern Sahelian zone than in the southern humid zone.
\item Frequency of hot days and nights has increased; cold extremes have decreased.
\item Rainfall trends are more complex: a slight decrease in the south, high inter-annual variability in the north, but an increase in \emph{intensity} of heavy rainfall events nationwide.
\end{itemize}
\end{aretenir}

\courssub{Impacts: Rainfall, Extreme Events, and Sea-Level Rise}

\courssub{Altered Rainfall Patterns}
Climate change alters the hydrological cycle through the Clausius–Clapeyron relation: a warmer atmosphere holds more moisture (~7\% more per \si{\ensuremath{{}^{\circ}\mathrm{C}}}). In Cameroon, this manifests as:
\begin{itemize}
\item \textbf{South (equatorial/bimodal):} Slight decrease in total annual rainfall, but longer dry spells within the rainy seasons, affecting cocoa and coffee yields.
\item \textbf{North (Sudano-Sahelian/unimodal):} High year-to-year variability; the rainy season onset is later and cessation earlier in some years, but extreme daily rainfall totals are increasing, leading to flash floods (e.g., Logone valley floods of 2012, 2022).
\item \textbf{Mount Cameroon region:} Orographic enhancement of rainfall may intensify, increasing landslide risk on steep volcanic slopes.
\end{itemize}

\courssub{Increased Extreme Events}
\begin{itemize}
\item \textbf{Heatwaves:} More frequent and intense, impacting human health (heat stress) and livestock productivity in the Far North.
\item \textbf{Floods:} Urban flooding in Douala and Yaoundé due to intense convective storms combined with poor drainage and unplanned urbanization.
\item \textbf{Droughts:} Meteorological and agricultural droughts in the North reduce groundwater recharge (critical for the Chad Basin aquifer) and cause crop failure.
\end{itemize}

\courssub{Sea-Level Rise on the Cameroonian Coast}
Global mean sea level has risen by \SI{20}{cm} since 1900, accelerating to \SI{3.7}{mm/yr} recently. For Cameroon's \SI{400}{km} coastline (Rio del Rey to Campo), the impacts are severe:
\begin{itemize}
\item \textbf{Coastal erosion:} Rates of \num{1}--\SI{3}{m/yr} in places like Limbe, Kribi, and the Wouri estuary, threatening infrastructure (Kribi deep-sea port, Limbe refinery) and ecosystems (mangroves).
\item \textbf{Saline intrusion:} Into the Wouri estuary and coastal aquifers, compromising freshwater supply for Douala (population > 3 million) and agriculture.
\item \textbf{Storm surges:} Higher baseline sea level allows storm surges to penetrate further inland, exacerbating flooding in low-lying neighborhoods (e.g., Bonaberi, Akwa).
\end{itemize}

\courssub{Climate Feedback Loops}

The climate system contains feedbacks that amplify (positive) or dampen (negative) the initial forcing. Understanding them is crucial for projecting future warming.

\begin{propriete}[Major Climate Feedbacks]
\begin{enumerate}
\item \textbf{Ice–Albedo Feedback (Positive):} Warming $\to$ melting snow/ice $\to$ lower surface albedo $\to$ more absorbed solar radiation $\to$ further warming. Critical in Arctic, but also on Mount Kilimanjaro and Mount Kenya (glaciers nearly gone); Mount Cameroon's permanent snow fields have vanished.
\item \textbf{Water Vapour Feedback (Positive):} Strongest feedback. Warming $\to$ more evaporation $\to$ more \ce{H2O} (a potent GHG) $\to$ enhanced greenhouse effect. Approximately doubles the \ce{CO2}-only warming.
\item \textbf{Lapse Rate Feedback (Negative in tropics, Positive at poles):} In tropics, upper troposphere warms more than surface $\to$ more efficient radiation to space $\to$ negative feedback.
\item \textbf{Cloud Feedback (Uncertain, likely Positive):} Low clouds cool (reflect sunlight); high clouds warm (trap IR). Net effect depends on cloud type/height changes.
\item \textbf{Permafrost Carbon Feedback (Positive):} Thawing permafrost releases \ce{CH4} and \ce{CO2} $\to$ additional forcing. Not directly relevant to Cameroon but globally significant.
\item \textbf{Vegetation/Soil Carbon Feedback (Positive):} Warming and drying $\to$ increased soil respiration, forest dieback (Amazon, Congo Basin) $\to$ reduced land sink or source. The Congo Basin peatlands (Cuvette Centrale) store \SI{30}{GtC}; drying could release massive \ce{CO2}.
\end{enumerate}
\end{propriete}

\begin{attention}[Weather vs. Climate – A Classic Trap]
\textbf{Mistake:} ``It rained heavily yesterday in Douala, so climate change is real'' OR ``It was cold in Yaoundé last week, so global warming is a hoax.''
\textbf{Correction:} \emph{Weather} is the state of the atmosphere at a given time and place (hours to days). \emph{Climate} is the statistical distribution of weather over \textbf{at least 30 years} (WMO standard). A single event cannot prove or disprove a trend. Attribution studies use climate models to estimate how much \emph{more likely} or \emph{more intense} an extreme event was due to human influence.
\end{attention}

\courssub{Mitigation and Adaptation Strategies for Cameroon}

Mitigation addresses the root cause (GHG emissions); adaptation manages the unavoidable impacts. Cameroon's National Determined Contribution (NDC) under the Paris Agreement integrates both.

\courssub{Mitigation: Reducing Emissions and Enhancing Sinks}
\begin{enumerate}
\item \textbf{Reforestation and Afforestation:} The \emph{Green Sahel} initiative and community forest management in the Congo Basin aim to restore degraded lands. Cameroon pledges to restore \SI{12}{million} hectares by 2030 (AFR100 initiative). \emph{Reducing Emissions from Deforestation and Forest Degradation (REDD+)} projects in the South and East regions monetize carbon storage.
\item \textbf{Renewable Energy:} Cameroon has huge hydro potential (\SI{23}{GW}, only \SI{5}{GW} exploited). The Nachtigal (\SI{420}{MW}), Memve'ele (\SI{211}{MW}), and Lom Pangar (\SI{30}{MW}) dams reduce reliance on thermal plants (heavy fuel oil). Solar potential is high in the North (\SI{5.5}{kWh/m^2/day}); mini-grids for rural electrification are expanding.
\item \textbf{Climate-Smart Agriculture (CSA):} Agroforestry (cocoa under shade trees), improved seeds (drought-tolerant maize, cassava), conservation agriculture (minimum tillage, mulching), and efficient irrigation (drip) reduce emissions per unit of food and build resilience.
\item \textbf{Sustainable Transport and Urban Planning:} Bus Rapid Transit (BRT) studies for Douala and Yaoundé; promoting non-motorized transport; building codes for energy efficiency.
\end{enumerate}

\courssub{Adaptation: Building Resilience}
\begin{enumerate}
\item \textbf{Early Warning Systems:} Meteorological radar network upgrades (Douala, Garoua, Maroua) for nowcasting extreme rainfall; SMS alerts for flood-prone communities (Logone, Benue basins).
\item \textbf{Coastal Protection:} Mangrove restoration (Wouri estuary, Rio del Rey) as natural buffers; managed retreat planning for highly vulnerable settlements; sand nourishment for critical beaches (Kribi).
\item \textbf{Water Security:} Rainwater harvesting in the North; managed aquifer recharge (MAR) in the Logone floodplain; rehabilitation of traditional wells with solar pumps.
\item \textbf{Health System Preparedness:} Heat-health action plans for cities; surveillance of climate-sensitive diseases (malaria expansion into highlands, cholera after floods).
\end{enumerate}

\begin{savaistu}[Cameroon's Volcanic Legacy and Carbon Storage]
The Cameroon Volcanic Line (CVL) offers a unique long-term mitigation opportunity: \textbf{mineral carbonation}. The abundant basaltic and ultramafic rocks (e.g., peridotites in the Adamawa and North regions) react with \ce{CO2} to form stable carbonate minerals (magnesite, calcite). Pilot projects worldwide (CarbFix in Iceland) inject \ce{CO2}-charged water into basalt, achieving mineralization in <2 years. Cameroon's geology could host such projects, turning a geological hazard (volcanism) into a climate solution.
\end{savaistu}

\courssub{Worked Example: Projecting Yaoundé's Temperature in 2050}

\begin{exoresolu}[Equilibrium vs. Transient Warming Projection]
\textbf{Statement:} Assume \ce{CO2} concentration follows a scenario reaching \SI{550}{ppm} by 2050 (from \SI{420}{ppm} in 2023). Using the simple energy-balance model:
\begin{enumerate}
\item Calculate the additional radiative forcing $\Delta F$ from 2023 to 2050.
\item Estimate the \emph{equilibrium} warming $\Delta T_{\mathrm{eq}}$ for this forcing (use $\lambda = \SI{0.8}{\ensuremath{{}^{\circ}\mathrm{C}}/(W.m^{-2})}$).
\item If the transient climate response (TCR) is about 60\% of the equilibrium response for a 1\%/yr \ce{CO2} increase, estimate the \emph{realized} warming by 2050 relative to 2023.
\item Given Yaoundé's 1960–2020 trend of \SI{0.155}{\ensuremath{{}^{\circ}\mathrm{C}}/decade}, what would a simple linear extrapolation give for 2050? Comment on the difference.
\end{enumerate}

\tcblower

\textbf{Detailed Solution:}
\begin{enumerate}
\item \textbf{Forcing change:}
$\Delta F = 5.35 \times \ln(550/420) = 5.35 \times \ln(1.3095) = 5.35 \times 0.270 = \SI{1.44}{W.m^{-2}}$.

\item \textbf{Equilibrium warming:}
$\Delta T_{\mathrm{eq}} = \lambda \times \Delta F = 0.8 \times 1.44 = \SI{1.15}{\ensuremath{{}^{\circ}\mathrm{C}}}$.

\item \textbf{Transient (realized) warming by 2050:}
TCR $\approx 0.6 \times \Delta T_{\mathrm{eq}} = 0.6 \times 1.15 = \SI{0.69}{\ensuremath{{}^{\circ}\mathrm{C}}}$.
So by 2050, we expect about \SI{0.7}{\ensuremath{{}^{\circ}\mathrm{C}}} \emph{additional} warming relative to 2023, assuming this scenario.

\item \textbf{Linear extrapolation from observed trend:}
From 2020 to 2050 is 30 years = 3 decades.
Extrapolated warming = $3 \times \SI{0.155}{\ensuremath{{}^{\circ}\mathrm{C}}} = \SI{0.465}{\ensuremath{{}^{\circ}\mathrm{C}}}$.
\textbf{Comment:} The linear extrapolation (\SI{0.47}{\ensuremath{{}^{\circ}\mathrm{C}}}) is lower than the physics-based transient estimate (\SI{0.69}{\ensuremath{{}^{\circ}\mathrm{C}}}). This is because the observed trend includes the cooling effect of aerosols and the lag from past forcing; as aerosol emissions decrease (air quality policies) and \ce{CO2} forcing accelerates, the warming rate is expected to \emph{increase}. Linear extrapolation underestimates future risk.
\end{enumerate}
\end{exoresolu}

\courssub{Synthesis and Examination Focus}

\begin{aretenir}[What the Examiner Expects You to Master]
\begin{itemize}
\item \textbf{Define} radiative forcing and distinguish it from temperature change.
\item \textbf{Apply} the formula $\Delta T = \lambda \Delta F$ (know $\lambda \approx \SI{0.8}{\ensuremath{{}^{\circ}\mathrm{C}}/(W.m^{-2})}$).
\item \textbf{Explain} the carbon cycle perturbation: sources (fossil fuels, land-use) and sinks (ocean, land).
\item \textbf{Describe} observed trends in Cameroon (temperature up \SI{~0.15}{\ensuremath{{}^{\circ}\mathrm{C}}/decade}, rainfall intensity up, sea level up \SI{~3.7}{mm/yr}).
\item \textbf{Analyse} impacts: agriculture (cocoa, cotton), coast (erosion, salinization), health (heat, disease), water (Logone/Chad Basin).
\item \textbf{Identify} key feedbacks (ice-albedo, water vapour, permafrost, Congo Basin peatlands).
\item \textbf{Propose} context-appropriate mitigation (REDD+, hydro/solar, CSA) and adaptation (mangroves, early warning, MAR).
\item \textbf{Avoid} the weather-vs-climate confusion; use 30-year statistics.
\end{itemize}
\end{aretenir}

\begin{exercice}[Practice Questions]
\begin{itemize}
\item Calculate the radiative forcing for a \ce{CO2} concentration of \SI{600}{ppm}. What is the equilibrium warming if $\lambda = \SI{0.75}{\ensuremath{{}^{\circ}\mathrm{C}}/(W.m^{-2})}$?
\item Sketch the carbon cycle and label the anthropogenic fluxes. Why is the airborne fraction (~45\%) roughly constant?
\item Explain how sea-level rise exacerbates saline intrusion in the Wouri estuary aquifer. What geological factors control the extent of intrusion?
\item List three positive climate feedbacks and one negative feedback. For each, state the mechanism in one sentence.
\item Design a climate-smart agriculture package for a maize farmer in the Far North region. Justify each component.
\end{itemize}
\end{exercice}

\begin{corrige}[Exercise 1]
$\Delta F = 5.35 \ln(600/280) = 5.35 \ln(2.143) = 5.35 \times 0.762 = \SI{4.08}{W.m^{-2}}$.
$\Delta T_{\mathrm{eq}} = 0.75 \times 4.08 = \SI{3.06}{\ensuremath{{}^{\circ}\mathrm{C}}}$.
(Note: This is the equilibrium warming from pre-industrial to 600 ppm, not the additional warming from today.)
\end{corrige}

\coursec{Pollution – Sources, Transport, and Management}

In the previous section we examined how greenhouse gases released by human activities modify the Earth's climate system. Climate change is, however, only one face of a broader problem: the introduction by humans of harmful substances into the geological environment — the atmosphere, surface water, groundwater and soils. This final section of the chapter addresses \cle{pollution} from a geologist's perspective: what pollutants are, where they come from, how they migrate through rocks, soils and water, what damage they cause, and how they can be cleaned up or controlled. For the GCE candidate, this topic is rich in examination questions because it combines chemistry, hydrogeology and environmental management, all with direct Cameroonian applications.

\courssub{Major Categories of Pollutants}

\begin{definition}[Pollution and pollutant]
\cle{Pollution} is the introduction into the environment of substances (or energy) at concentrations that cause harm to human health, living organisms or ecosystems. A \cle{pollutant} is any substance present at a concentration above its natural background level as a result of human activity, and capable of producing adverse effects.
\end{definition}

A useful way to organise the enormous variety of pollutants is to group them into five major categories, each with distinct sources, behaviours and risks.

\textbf{Heavy metals.} Elements such as lead (Pb), mercury (Hg), cadmium (Cd), arsenic (As), chromium (Cr) and zinc (Zn) are released by mining, ore smelting, battery manufacturing, tanneries and the burning of fossil fuels. Their danger lies in two properties: they are \emph{toxic at very low concentrations}, and they are \emph{not biodegradable} — once released, they persist in soils and sediments for centuries. In the gold-mining districts around \textbf{Bertoua} (East Region), artisanal miners use mercury to amalgamate gold; the mercury is then released into streams and soils, where micro-organisms can convert it to the highly toxic organic form \cle{methylmercury}.

\textbf{Hydrocarbons.} Crude oil and refined products (petrol, diesel, kerosene, lubricants) enter the environment through spills during extraction, transport, storage and refining. The petroleum activities of the \textbf{Kribi} area (offshore fields of the Rio del Rey and Douala/Kribi-Campo basins, plus the Chad–Cameroon pipeline terminal) illustrate both the economic importance and the environmental risk of this category. Hydrocarbons are biodegradable in principle, but slowly; light fractions (benzene, toluene) are toxic and mobile in groundwater, while heavy fractions persist as tarry residues in sediments.

\textbf{Nutrients.} Nitrogen and phosphorus compounds — mainly nitrates ($\mathrm{NO_3^-}$), ammonium ($\mathrm{NH_4^+}$) and phosphates ($\mathrm{PO_4^{3-}}$) — come from fertilisers, manure and untreated sewage. They are not directly toxic, but in excess they trigger \cle{eutrophication} (see below). The intensive agriculture of the \textbf{North West Region} (market gardening around Santa and Bafut, cattle ranching on the grassfields) exports nutrients to streams and crater lakes.

\textbf{Plastics.} Synthetic polymers are extremely resistant to degradation. Macroplastics choke drainage channels and aggravate urban flooding (a familiar sight in Douala during the rains), while \cle{microplastics} (particles smaller than \SI{5}{mm}) are ingested by aquatic organisms and enter food chains.

\textbf{Radioactive waste.} Ionising radiation sources include mining and processing of uranium-bearing ores, medical and industrial sources, and (globally) nuclear power generation. Radioactive pollutants are dangerous because they cannot be detected by the senses, and some radionuclides have half-lives of thousands of years, requiring very long-term containment in stable geological formations.

\begin{center}
\begin{tabularx}{\linewidth}{|l|X|X|}
\hline
\rowcolor{popblueL} \textbf{Category} & \textbf{Main sources} & \textbf{Key risk} \\
\hline
Heavy metals (Pb, Hg, Cd, As) & Mining, smelting, batteries, tanneries & Persistent, toxic at trace levels, bioaccumulate \\
\hline
Hydrocarbons & Oil spills, leaking tanks, refineries & Toxic light fractions mobile in groundwater \\
\hline
Nutrients (N, P) & Fertilisers, manure, sewage & Eutrophication of lakes and rivers \\
\hline
Plastics & Packaging, urban waste & Drainage blockage, microplastics in food chains \\
\hline
Radioactive waste & U-ore processing, medical/industrial sources & Invisible, very long-lived, needs deep containment \\
\hline
\end{tabularx}
\end{center}

\courssub{Point Sources and Non-Point Sources}

\begin{definition}[Point and non-point sources]
A \cle{point source} discharges pollutants from a single, identifiable location (a factory outfall pipe, a mine tailings dam, a leaking fuel tank). A \cle{non-point source} (or diffuse source) releases pollutants over a wide area (fertilised farmland, urban runoff, atmospheric fallout), so that no single discharge point can be identified.
\end{definition}

The distinction matters enormously for management. Point sources can be measured, licensed and shut off; non-point sources can only be controlled by changing practices over whole landscapes. Typical Cameroonian examples:

\begin{itemize}
\item \textbf{Mining (Bertoua):} a tailings heap or a mercury-using washing station is a point source; the slow weathering of abandoned spoil over hectares is a non-point source.
\item \textbf{Oil industry (Kribi):} a pipeline leak or a refinery effluent is a point source; chronic small leaks from many boats, engines and storage depots form a diffuse source.
\item \textbf{Agriculture (North West):} fertiliser and pesticide runoff from thousands of small farms is the classic non-point source, delivering nitrates to streams and lakes after every heavy rain.
\end{itemize}

\begin{attention}[One method does not fit all]
Treating all pollutants with a single remediation method is ineffective. Heavy metals cannot be biodegraded; nutrients cannot be ``filtered out'' of a lake once dispersed; hydrocarbons volatilise while metals do not. Effective management always begins by identifying the pollutant category, its source type (point or diffuse) and its transport pathway — only then can an appropriate technique be chosen.
\end{attention}

\courssub{Transport Pathways of Pollutants}

A pollutant only causes harm if it travels from its \cle{source} to a \cle{receptor} (a person, an ecosystem, a water well) along a \cle{pathway}. This \cle{source–pathway–receptor} model is the backbone of all pollution risk analysis: breaking the chain at any point prevents harm. Three geological pathways dominate.

\textbf{Atmospheric deposition.} Dust, gases and aerosols emitted by industry, vehicles or bush burning are transported by wind and returned to the surface by dry fallout or by rain (wet deposition). Acid gases ($\mathrm{SO_2}$, $\mathrm{NO_x}$) produce acid rain, which acidifies soils and lakes far from the source.

\textbf{Surface runoff.} Rainwater flowing over the land dissolves and carries pollutants into streams, rivers and lakes. Runoff is fast (hours to days), so contamination events are often sudden — for example, the first heavy rains of the season flushing accumulated urban waste into the Wouri estuary.

\textbf{Groundwater advection.} Pollutants that infiltrate the soil are carried slowly through aquifers by groundwater flow — a process called \cle{advection}. Because groundwater velocities are small (often centimetres to metres per year), a contaminated zone, or \cle{plume}, may take decades to reach a well, but it is then extremely difficult to clean. As the plume travels, it spreads (\cle{dispersion}) and may be attenuated by sorption onto mineral surfaces and by chemical or biological decay.

\begin{popfigure}
\begin{center}
\begin{tikzpicture}[font=\small, >=stealth]
\node[draw=popdark, fill=poporange!25, rounded corners, text width=2.6cm, align=center] (src) at (0,0) {\textbf{Source}\\ mine, spill, farm};
\node[draw=popdark, fill=popblueL, rounded corners, text width=2.4cm, align=center] (atm) at (4.4,2.2) {Atmospheric deposition};
\node[draw=popdark, fill=popblueL, rounded corners, text width=2.4cm, align=center] (run) at (4.4,0) {Surface runoff};
\node[draw=popdark, fill=popblueL, rounded corners, text width=2.6cm, align=center] (gw) at (4.4,-2.2) {Groundwater advection (plume)};
\node[draw=popdark, fill=popgreenL, rounded corners, text width=2.6cm, align=center] (rec) at (9.4,0) {\textbf{Receptor}\\ well, river, community, ecosystem};
\draw[->, thick, popdark] (src) -- (atm);
\draw[->, thick, popdark] (src) -- (run);
\draw[->, thick, popdark] (src) -- (gw);
\draw[->, thick, popdark] (atm) -- (rec);
\draw[->, thick, popdark] (run) -- (rec);
\draw[->, thick, popdark] (gw) -- (rec);
\node[text width=2.8cm, align=center, popmuted] at (4.4,-3.6) {slow, hidden, persistent};
\node[text width=2.4cm, align=center, popmuted] at (4.4,0.9) {fast, episodic};
\node[text width=2.4cm, align=center, popmuted] at (4.4,3.1) {regional scale};
\end{tikzpicture}
\end{center}
\legende{Flow diagram of pollutant migration from source to receptor through the three main geological pathways.}
\end{popfigure}

\courssub{Quantifying Plume Behaviour: First-Order Decay}

Many pollutants (biodegradable hydrocarbons, some radionuclides, chlorinated solvents under natural attenuation) decrease in concentration with time according to \cle{first-order decay}: the rate of decrease is proportional to the concentration present.

\begin{propriete}[First-order decay and half-life]
For a pollutant undergoing first-order decay, the concentration at time $t$ is
\[ C(t) = C_0\, e^{-kt} \]
where $C_0$ is the initial concentration and $k$ is the decay constant (in $\mathrm{time^{-1}}$). The \cle{half-life} $t_{1/2}$ — the time for the concentration to fall to half its initial value — is related to $k$ by
\[ t_{1/2} = \frac{\ln 2}{k} \approx \frac{0.693}{k}. \]
\emph{Derivation (examinable):} setting $C(t_{1/2}) = C_0/2$ gives $e^{-k t_{1/2}} = 1/2$, so $-k t_{1/2} = \ln(1/2) = -\ln 2$, hence $t_{1/2} = \ln 2 / k$.
\end{propriete}

\begin{exoresolu}[Decay of a contaminant plume]
\textbf{Statement.} A leaking fuel tank at a depot contaminates an aquifer with benzene at an initial concentration $C_0 = 800\ \mu\mathrm{g\,L^{-1}}$. Natural biodegradation follows first-order kinetics with decay constant $k = 0.023\ \mathrm{yr^{-1}}$. (a) Calculate the half-life of benzene in this aquifer. (b) The drinking-water limit is $10\ \mu\mathrm{g\,L^{-1}}$. After how many years will the plume concentration fall below this limit?
\tcblower
\textbf{Solution.}
(a) Apply the half-life relation:
\[ t_{1/2} = \frac{\ln 2}{k} = \frac{0.693}{0.023} \approx 30.1\ \mathrm{yr}. \]
(b) We require $C(t) = C_0 e^{-kt} \le 10$:
\[ e^{-kt} = \frac{10}{800} = 0.0125 \quad\Longrightarrow\quad t = \frac{\ln(1/0.0125)}{k} = \frac{\ln 80}{0.023} \approx \frac{4.382}{0.023} \approx 190.5\ \mathrm{yr}. \]
\textbf{Interpretation:} even though the concentration halves every 30 years, reaching the drinking-water standard takes more than six half-lives — nearly two centuries. This is why prevention of groundwater pollution is vastly preferable to cure.
\end{exoresolu}

\begin{popfigure}
\begin{center}
\begin{tikzpicture}
\begin{axis}[
width=0.85\linewidth, height=6.2cm,
xlabel={Distance from source $x$ (m)},
ylabel={Concentration $C$ ($\mu$g\,L$^{-1}$)},
xmin=0, xmax=500, ymin=0, ymax=850,
grid=both, grid style={popline!40},
axis line style={popdark},
label style={popdark},
tick label style={popdark},
]
\addplot[domain=0:500, samples=80, very thick, popblue] {800*exp(-x/80)};
\addplot[dashed, thick, poporange] coordinates {(0,10) (500,10)};
\node[poporange, anchor=west, font=\small] at (axis cs:300,60) {drinking-water limit};
\node[popblue, anchor=west, font=\small] at (axis cs:120,500) {$C = C_0 e^{-x/\lambda}$};
\end{axis}
\end{tikzpicture}
\end{center}
\legende{Concentration–distance curve for a contaminant plume in an aquifer under steady flow with first-order attenuation: concentration falls exponentially with distance from the source.}
\end{popfigure}

\courssub{Impact Assessment: How Pollutants Cause Harm}

\begin{definition}[Bioaccumulation]
\cle{Bioaccumulation} is the increase of the concentration of a pollutant in the tissues of an organism over time, because the rate of uptake (from water, food or sediment) exceeds the rate of elimination. When concentrations increase from one trophic level to the next up a food chain, the process is called \cle{biomagnification}.
\end{definition}

Bioaccumulation explains why mercury released in small quantities by artisanal gold miners can reach dangerous levels in large predatory fish — and therefore in the people who eat them — even when the concentration in the river water itself seems negligible.

\textbf{Eutrophication.} Excess nutrients (N and P) in a lake or slow river trigger explosive growth of algae. When the algae die, their decomposition by bacteria consumes dissolved oxygen, causing fish kills and the collapse of aquatic life. Shallow crater lakes of the North West and coastal lagoons receiving agricultural runoff are vulnerable.

\textbf{Acid mine drainage.} Many ore deposits contain pyrite ($\mathrm{FeS_2}$). When mining exposes pyrite to air and water, it oxidises:
\[ \ce{4FeS2 + 15O2 + 14H2O -> 4Fe(OH)3 + 8H2SO4} \]
The sulphuric acid produced lowers the pH of drainage water (sometimes below 3) and dissolves heavy metals, which are then carried into streams. Acid mine drainage can continue for decades after a mine is abandoned, making it one of the most persistent legacies of mining.

\textbf{Health effects.} Lead damages the nervous system of children; cadmium attacks kidneys and bones; methylmercury causes neurological disorders; benzene is carcinogenic; nitrate in drinking water causes methaemoglobinaemia (``blue baby syndrome'') in infants. The geologist's role is to identify the pathways by which these substances reach people — especially through groundwater, the main drinking-water source in rural Cameroon.

\begin{methode}[Steps of a Tier-1 Environmental Impact Assessment]
A \cle{Tier-1 Environmental Impact Assessment} (EIA) is the preliminary, screening-level evaluation carried out before any major project (mine, dam, refinery, road). Its steps are:
\begin{enumerate}
\item \textbf{Screening:} decide whether the project requires a full EIA at all, based on its type, size and location (e.g.\ proximity to a water supply or protected area).
\item \textbf{Scoping:} identify which environmental components (groundwater, surface water, soils, air, ecosystems, communities) and which potential impacts the study must address; define the study boundaries and consult stakeholders.
\item \textbf{Baseline study:} measure the existing state of the environment before the project begins — water chemistry, soil metal contents, air quality, biodiversity — so that future changes can be detected and attributed.
\end{enumerate}
Later tiers add impact prediction, mitigation design, and monitoring, but a Tier-1 study already determines whether a project is acceptable and what must be protected.
\end{methode}

\courssub{Remediation Techniques and Regulatory Framework}

\begin{definition}[Remediation]
\cle{Remediation} is the set of techniques used to remove, destroy, contain or neutralise pollutants in soil, water or sediment in order to reduce risks to acceptable levels.
\end{definition}

\textbf{Phytoremediation} uses plants to clean contaminated soil and water. Some plants (\emph{hyperaccumulators}) take up heavy metals through their roots and store them in their harvestable shoots, which are then removed; others stabilise contaminants in the root zone. It is cheap and aesthetically acceptable, but slow and limited to shallow contamination.

\textbf{Bioremediation} uses micro-organisms (bacteria, fungi) to degrade organic pollutants, especially hydrocarbons. It may rely on stimulating native microbes by adding oxygen and nutrients (\emph{biostimulation}) or on introducing specialised strains (\emph{bioaugmentation}). It is the standard response to oil spills on land, but it does not work on metals.

\textbf{Constructed wetlands} are engineered shallow basins planted with reeds and other macrophytes. Water flowing slowly through them is purified by a combination of sedimentation, plant uptake and microbial processes. They are effective for treating mine drainage, agricultural runoff and domestic wastewater, and are well suited to tropical climates like Cameroon's because warm temperatures sustain year-round biological activity.

\textbf{Containment and physical methods} include pumping and treating contaminated groundwater, excavating polluted soil, and capping waste sites with impermeable barriers to cut the pathway between source and receptor.

\textbf{Regulatory framework.} In Cameroon, environmental protection is anchored in \textbf{Law No.\ 96/12 of 5 August 1996 relating to environmental management}, which establishes the polluter-pays principle, makes Environmental Impact Assessments compulsory for major projects, and sets the basis for pollution control and sanctions. Sectoral texts (on water, mining and petroleum activities) complement it. For the geologist, the practical consequence is clear: no mine, quarry or industrial installation should open without a baseline study and an approved management plan.

\begin{aretenir}[Key points to remember]
\begin{itemize}
\item Pollutants fall into five major categories: heavy metals, hydrocarbons, nutrients, plastics and radioactive waste — each with distinct sources, behaviour and treatment.
\item Distinguish point sources (identifiable, controllable) from non-point sources (diffuse, managed only by changing practices).
\item The source–pathway–receptor chain (atmospheric deposition, surface runoff, groundwater advection) structures all risk analysis; breaking any link prevents harm.
\item First-order decay: $C = C_0 e^{-kt}$ and $t_{1/2} = \ln 2 / k$; groundwater clean-up is measured in decades, so prevention is paramount.
\item Bioaccumulation, eutrophication and acid mine drainage are the three classic impact mechanisms to describe in an examination.
\item Remediation must match the pollutant: phytoremediation and bioremediation for suitable contaminants, constructed wetlands for drainage and runoff, all within the framework of Law No.\ 96/12.
\end{itemize}
\end{aretenir}

\begin{exercice}[Choosing a remediation strategy]
An abandoned artisanal mining site near Bertoua shows (i) soils contaminated with mercury and arsenic, and (ii) a small pond receiving acidic, metal-rich drainage. For each problem, propose an appropriate remediation technique, justify your choice, and explain why bioremediation alone would fail.
\end{exercice}

\begin{corrige}[Exercise 1]
(i) \emph{Metals in soil:} metals cannot be degraded, so bioremediation is useless. Suitable options are phytoremediation with hyperaccumulator plants (slow but cheap) or excavation and safe disposal of the most contaminated soil, possibly combined with soil stabilisation to immobilise the metals. (ii) \emph{Acid drainage:} a constructed wetland, possibly preceded by addition of limestone to neutralise acidity, allows metals to precipitate and be retained in the sediment. Bioremediation alone fails in both cases because micro-organisms can only transform organic compounds; they cannot destroy mercury, arsenic or acidity, which require physical, chemical or plant-based treatment.
\end{corrige}

\begin{exercice}[Half-life and safe use of a well]
A pesticide spilled into an aquifer has an initial concentration of $400\ \mu\mathrm{g\,L^{-1}}$ and a half-life of 12 years under natural attenuation. (a) Calculate the decay constant $k$. (b) Calculate the concentration remaining after 36 years. (c) Comment on the result if the health limit is $50\ \mu\mathrm{g\,L^{-1}}$.
\end{exercice}

\begin{corrige}[Exercise 2]
(a) $k = \ln 2 / t_{1/2} = 0.693 / 12 \approx 0.0578\ \mathrm{yr^{-1}}$. (b) 36 years corresponds to three half-lives, so $C = 400 \times (1/2)^3 = 400/8 = 50\ \mu\mathrm{g\,L^{-1}}$ (equivalently $C = 400\,e^{-0.0578 \times 36} \approx 50$). (c) The water only reaches the health limit after 36 years: the aquifer is unusable for a whole generation, which again shows that protecting groundwater from pollution is far more effective than waiting for natural clean-up.
\end{corrige}

\begin{savaistu}[Pollution and Cameroonian development]
Cameroon's ambition to expand mining and petroleum production makes environmental geology a strategic discipline. Every new gold panning site in the East, every well drilled off Kribi, and every fertiliser bag spread on the North West highlands creates a potential source–pathway–receptor chain. Trained geologists who can run a baseline study, model a plume and design a wetland are precisely the professionals the country needs — perhaps you will be one of them.
\end{savaistu}

\newpage

\coursec{Integration activity}

\courssub{Aims of the activity}

This integration activity brings together the whole chapter: the definition and scope of \cle{environmental geology}, the identification of \cle{geohazards}, the analysis of their \cle{effects and mitigation}, the use of \cle{climate projections}, and the tracing of \cle{pollution pathways}. By the end of the activity, you should be able to:

\begin{itemize}
\item read and interpret a geological map, a topographic map and climate data to identify seismic, landslide and flood hazards on a given site;
\item use simplified climate model outputs to project rainfall changes to 2050 and discuss their consequences for hazard frequency and magnitude;
\item apply the concept of a \cle{source--pathway--receptor} chain to model the transport of cyanide and heavy metals from a mine to a river;
\item construct and fill a \cle{risk matrix} (likelihood $\times$ consequence) and use it to rank risks;
\item propose and justify \cle{structural} and \cle{non-structural} mitigation measures;
\item communicate scientific conclusions clearly in a short oral briefing, as geologists do before regulatory boards such as the Cameroonian administration in charge of mines and of the environment.
\end{itemize}

\begin{aretenir}[Skills mobilised]
Map reading $\rightarrow$ data analysis $\rightarrow$ simple modelling $\rightarrow$ risk evaluation $\rightarrow$ decision-making $\rightarrow$ scientific communication. This is exactly the professional reasoning chain of an environmental geologist, and it is the chain the GCE examiner expects you to reproduce in an integration paper.
\end{aretenir}

\courssub{Situation}

\textbf{The Kadey Gold Project.} The (fictional) company \emph{Kadey Gold SARL} has applied for an exploitation permit for an open-pit gold mine on a \SI{12}{km^2} concession in the East Region of Cameroon, in the Kadey Division, near the Batouri--Yokadouma axis. The ore is a quartz-vein gold deposit hosted in Precambrian schists and granites near the margin of the Congo craton. The company plans:

\begin{itemize}
\item an open pit \SI{800}{m} long, \SI{400}{m} wide and up to \SI{120}{m} deep;
\item a processing plant using \cle{cyanide heap leaching} to extract gold from crushed ore;
\item a tailings storage facility (TSF) of \SI{25}{ha} in a small valley;
\item a camp for 600 workers and a \SI{15}{km} access road.
\end{itemize}

Before granting the permit, the regulatory board requires an \textbf{Environmental and Social Impact Assessment (ESIA)}. Your group is the team of consulting geologists. You receive the following dossier.

\textbf{Document 1 -- Simplified geological and topographic data.}

\begin{center}
\begin{tabularx}{\linewidth}{|l|X|}
\hline
\rowcolor{popblueL} \textbf{Feature} & \textbf{Description} \\
\hline
Bedrock & Precambrian schists (foliation dipping $45^{\circ}$ towards the river valley) cut by granitic intrusions; deeply weathered lateritic cover, 5--20~m thick. \\
\hline
Structures & A major regional fault (Bertoua fault zone) crosses the NW corner of the concession; minor fault splays pass beneath the proposed TSF. \\
\hline
Relief & Undulating plateau at 650--700~m; the concession slopes down towards the Lom River (a tributary of the Sanaga) at 590~m; slope angles of $20^{\circ}$--$35^{\circ}$ on the valley flanks. \\
\hline
Hydrography & The Lom River flows 1.2~km east of the pit; two seasonal streams cross the concession and join it. \\
\hline
Population & Village of Ndembo (1\,800 inhabitants) 2~km downstream; villagers use the river for fishing, washing and drinking water; 400~ha of cocoa farms along the floodplain. \\
\hline
\end{tabularx}
\end{center}

\textbf{Document 2 -- Seismic and hazard context.} Cameroon is moderately seismic. The Mount Cameroon--Cameroon Volcanic Line region and the Foumban shear zone have produced damaging historical earthquakes; the East Region is less active, but the Bertoua fault zone is classified as \emph{potentially active}, with an estimated maximum magnitude of $M = 5.5$ and a return period of strong shaking of about 100--200 years. The region also records frequent landslides on lateritic slopes during the rainy season, and the Lom River flooded its plain in 2012 and 2019, submerging cocoa farms under up to 1.5~m of water.

\textbf{Document 3 -- Climate data and projection.} Present climate (1991--2020 averages): mean annual rainfall \SI{1500}{mm}; a long rainy season (March--November) with a peak in September--October; mean annual temperature \SI{24.5}{\ensuremath{{}^{\circ}\mathrm{C}}}. Simplified CMIP6 model outputs for the region (scenario SSP2-4.5) give, for 2050:

\begin{center}
\begin{tabular}{|l|c|c|}
\hline
\rowcolor{popblueL} \textbf{Parameter} & \textbf{1991--2020} & \textbf{Projection 2050} \\
\hline
Mean annual temperature ($^{\circ}\mathrm{C}$) & 24.5 & 26.0 \\
\hline
Mean annual rainfall (mm) & 1500 & 1450 \\
\hline
Rainfall in wettest month (mm) & 280 & 320 \\
\hline
Days with rain $>$ \SI{50}{mm} per year & 8 & 12 \\
\hline
Length of dry season (months) & 3 & 4 \\
\hline
\end{tabular}
\end{center}

\textbf{Document 4 -- Baseline water quality of the Lom River (upstream of the site).}

\begin{center}
\begin{tabular}{|l|c|c|}
\hline
\rowcolor{popblueL} \textbf{Parameter} & \textbf{Measured value} & \textbf{WHO guideline} \\
\hline
pH & 6.8 & 6.5--8.5 \\
\hline
Arsenic ($\mathrm{\mu g\,L^{-1}}$) & 4 & 10 \\
\hline
Mercury ($\mathrm{\mu g\,L^{-1}}$) & 0.3 & 6 \\
\hline
Cyanide ($\mathrm{\mu g\,L^{-1}}$) & $<1$ & 70 \\
\hline
Turbidity (NTU) & 12 & $<25$ \\
\hline
\end{tabular}
\end{center}

\textbf{Document 5 -- Process data.} The plant will use \SI{2000}{t} of sodium cyanide per year; tailings will contain residual cyanide (up to \SI{50}{mg/kg}) and elevated arsenic and mercury released from the ore during processing. The TSF is designed with a \emph{single} compacted clay liner. The pit will intersect the water table at about \SI{40}{m} depth, requiring pumping (dewatering).

\courssub{Tasks}

\begin{enumerate}
\item \textbf{Hazard mapping.} Using Documents 1 and 2, identify and locate on a sketch map the seismic, landslide and flood hazards affecting the concession. For each hazard, state the geological factor that controls it and the element of the project most exposed (pit, TSF, plant, road, camp).
\item \textbf{Climate projection.} Using Document 3, calculate the percentage change between the reference period and 2050 for (a) wettest-month rainfall and (b) the number of intense-rain days. Explain, with two arguments, how these changes could modify the frequency or magnitude of landslides and floods on the site, and why a longer dry season is also a problem for a mine.
\item \textbf{Pollution pathway analysis.} Using Documents 1, 4 and 5, describe the complete \cle{source--pathway--receptor} chain for (a) cyanide and (b) arsenic/mercury, from the mine to the villagers of Ndembo. Identify at least three possible failure events (e.g.\ liner rupture, TSF overflow during an extreme storm, fault-guided seepage) and state which pathway each one opens.
\item \textbf{Risk matrix.} Complete a risk matrix for the five events listed below, scoring likelihood (1 = rare, 2 = possible, 3 = likely) and consequence (1 = minor, 2 = moderate, 3 = severe), then compute the risk $R = L \times C$ and rank the events: (E1) earthquake-induced failure of the TSF; (E2) landslide into the open pit; (E3) flood overflow of the TSF during an extreme storm; (E4) chronic seepage of metals through the liner; (E5) river contamination by a cyanide spill on the access road.
\item \textbf{Mitigation plan.} Propose \textbf{at least three structural} measures (engineering works) and \textbf{at least two non-structural} measures (monitoring, regulation, planning, education). For each measure, state which risk(s) it reduces and justify it scientifically.
\item \textbf{Briefing.} Prepare a 5-minute oral presentation to the mock regulatory board: one slide with your hazard sketch map, one with the risk matrix, one with the mitigation plan, and a clear final recommendation (authorise, authorise with conditions, or refuse).
\end{enumerate}

\begin{methode}[How to build a risk matrix]
\begin{enumerate}
\item List the hazardous \emph{events} (not just the hazards): a hazard becomes an event when it actually affects an exposed element.
\item Score the \textbf{likelihood} $L$ of each event over the project life, using the data (return periods, climate trends, design weaknesses).
\item Score the \textbf{consequence} $C$ for people, property and the environment if the event occurs.
\item Compute $R = L \times C$ and plot the events in a $3 \times 3$ grid; $R \geq 6$ = high risk (red), $R = 3$--$4$ = moderate (orange), $R \leq 2$ = low (green).
\item Rank the risks and attach mitigation measures first to the highest risks.
\end{enumerate}
\end{methode}

\courssub{Model answer}

\textbf{1. Hazard mapping.} Three hazards are identified. (a) \emph{Seismic hazard}: the Bertoua fault zone crosses the NW corner, and fault splays pass beneath the TSF; the controlling factor is the potentially active fault ($M_{\max} = 5.5$); the most exposed structure is the TSF, whose embankment and liner could rupture under shaking. (b) \emph{Landslide hazard}: the valley flanks have slopes of $20^{\circ}$--$35^{\circ}$, a thick lateritic weathering cover, and schist foliation dipping $45^{\circ}$ towards the valley -- a daylighting discontinuity that favours planar slides; the open-pit walls, the access road and the TSF valley are exposed. (c) \emph{Flood hazard}: the Lom River has flooded twice in recent decades (2012, 2019, up to 1.5~m of water); the floodplain installations, the road and the lower part of the TSF are exposed. On the sketch map below, the seismic zone is drawn along the NW fault corridor, the landslide-prone band along the steep eastern slopes, and the flood zone along the Lom floodplain.

\begin{popfigure}
\begin{tikzpicture}[scale=0.9, every node/.style={font=\small}]
% concession boundary
\draw[thick, popdark] (0,0) rectangle (8,6);
\node[popdark] at (4,6.35) {Concession boundary (\SI{12}{km^2})};
% fault zone NW
\draw[very thick, poppink, dashed] (-0.2,5.6) -- (2.6,6.2);
\draw[very thick, poppink, dashed] (0.4,4.9) -- (3.0,5.5);
\node[poppink, align=center] at (1.2,4.4) {Bertoua fault zone\\(seismic hazard)};
% fault splays under TSF
\draw[poppink, dashed] (4.6,3.4) -- (5.6,4.4);
\draw[poppink, dashed] (5.0,3.0) -- (6.0,4.0);
% pit
\draw[fill=popgold!40, draw=popdark] (2.2,2.6) ellipse (1.1 and 0.6);
\node at (2.2,2.6) {Open pit};
% plant
\draw[fill=popblueL, draw=popdark] (3.9,2.2) rectangle (4.7,2.8);
\node at (4.3,1.9) {Plant};
% TSF
\draw[fill=popgreenL, draw=popdark] (5.0,3.6) rectangle (6.2,4.4);
\node at (5.6,4.7) {TSF};
% camp
\draw[fill=popmuted!30, draw=popdark] (1.0,4.6) rectangle (1.8,5.1);
\node at (1.4,4.25) {Camp};
% road
\draw[thick, popmuted] (0,1.0) -- (3.0,1.4) -- (5.0,1.2) -- (8,1.6);
\node[popmuted] at (1.6,0.7) {Access road};
% steep slopes band
\draw[poporange, very thick, decorate, decoration={coil, amplitude=1.5pt, segment length=6pt}] (6.6,0.4) -- (6.6,5.8);
\node[poporange, align=center] at (5.6,0.35) {Steep slopes $20^{\circ}$--$35^{\circ}$: landslide band};
% river
\draw[very thick, popblue] (7.6,0) .. controls (7.2,2) and (7.9,4) .. (7.5,6);
\node[popblue, rotate=90] at (7.85,3.0) {Lom River};
% floodplain
\draw[popblue, dashed, fill=popblueL, opacity=0.5] (6.9,0.2) -- (7.4,0.2) -- (7.3,5.8) -- (6.9,5.8) -- cycle;
\node[popblue, align=center] at (7.0,5.0) {\phantom{xx}};
\node[popblue] at (6.15,5.6) {Flood zone};
% streams
\draw[popblue] (3.0,3.4) .. controls (4.5,3.6) .. (7.0,3.2);
\draw[popblue] (2.6,1.8) .. controls (4.5,2.4) .. (7.1,2.2);
% village
\draw[fill=poppink!40, draw=popdark] (7.0,-0.9) circle (0.25);
\node[align=center] at (7.0,-1.35) {Ndembo (2~km downstream)};
% north arrow
\draw[->, thick] (0.5,5.5) -- (0.5,6.0) node[above] {N};
\end{tikzpicture}
\legende{Sketch hazard map of the Kadey concession: seismic corridor along the NW fault zone, landslide-prone band on the steep eastern valley flanks, and flood zone along the Lom River floodplain.}
\end{popfigure}

\textbf{2. Climate projection.} Percentage changes to 2050:
\begin{align*}
\text{Wettest month: } & \frac{320 - 280}{280} \times 100 \approx +14.3\,\% \\
\text{Intense-rain days: } & \frac{12 - 8}{8} \times 100 = +50\,\%
\end{align*}
Although total annual rainfall slightly decreases ($\frac{1450-1500}{1500}\times 100 \approx -3.3\,\%$), rainfall becomes \emph{more concentrated in extreme events}. First, more days with rain $>$ \SI{50}{mm} mean more frequent saturation of the lateritic cover, raising pore-water pressure, reducing the shear strength of the soil and lowering slope stability -- landslides become more likely. Second, higher wettest-month rainfall increases peak river discharge, so floods like those of 2012 and 2019 become more frequent or deeper, threatening TSF overflow. Conversely, the longer dry season (4 months instead of 3) reduces river flow and therefore its dilution capacity, reduces water availability for processing and for the camp, and concentrates any pollutant that does reach the river -- low-flow contamination is more dangerous than high-flow contamination.

\textbf{3. Pollution pathways.} The chain is: \textbf{source} (cyanide in the leach pads and tailings; As and Hg released from the ore and stored in the tailings) $\rightarrow$ \textbf{pathway} (surface runoff to the seasonal streams; seepage through a damaged liner to groundwater; fault-guided flow along the splays beneath the TSF; direct spill on the road where it crosses a stream) $\rightarrow$ \textbf{receptor} (the Lom River, then the villagers of Ndembo who drink, fish and wash, and the 400~ha of cocoa farms on the floodplain). Three failure events: (i) \emph{liner rupture} (by earthquake or differential settlement) opens the groundwater pathway; (ii) \emph{TSF overflow during an extreme storm} opens the surface-water pathway directly to the river; (iii) \emph{fault-guided seepage} bypasses the liner and delivers contaminated water to the alluvial aquifer used by the village. Baseline water is currently clean (cyanide $< 1$~$\mathrm{\mu g\,L^{-1}}$, well below the WHO guideline of 70~$\mathrm{\mu g\,L^{-1}}$), so any increase will be clearly attributable to the mine -- a strong argument for rigorous monitoring.

\textbf{4. Risk matrix.} Justified scores:

\begin{center}
\begin{tabular}{|l|c|c|c|l|}
\hline
\rowcolor{popblueL} \textbf{Event} & $L$ & $C$ & $R = L \times C$ & \textbf{Rank} \\
\hline
E1 Earthquake failure of TSF & 1 & 3 & 3 & Moderate \\
\hline
E2 Landslide into open pit & 3 & 2 & 6 & High \\
\hline
E3 Flood overflow of TSF & 3 & 3 & 9 & High \\
\hline
E4 Chronic seepage of metals & 2 & 2 & 4 & Moderate \\
\hline
E5 Cyanide spill on access road & 2 & 3 & 6 & High \\
\hline
\end{tabular}
\end{center}

E3 scores highest: the TSF sits in a valley near the floodplain, extreme-rain days increase by 50\,\% by 2050, and a release of tailings and cyanide downstream would be catastrophic for Ndembo and the cocoa farms. E2 is very likely given the steep foliated slopes and thick lateritic cover, though its consequences fall mainly on workers in the pit. E1 is rare (return period 100--200 years) but severe, hence a moderate risk that still justifies seismic design of the embankment.

\textbf{5. Mitigation plan.} \emph{Structural measures}: (a) build the TSF with a \textbf{double composite liner} (compacted clay + geomembrane) and a leakage-collection drain between the two layers, reducing E1 and E4; (b) design the TSF embankment and spillway for the \textbf{2050-projected extreme rainfall} (freeboard increased, emergency spillway sized for the $+14\,\%$ wettest month), reducing E3; (c) \textbf{regrade and drain the pit slopes} (benches, diversion channels at the crest, slope angle reduced below the foliation dip where possible), reducing E2; (d) line the leach pads and build a \textbf{cyanide detoxification unit} (chemical oxidation of cyanide before any discharge), reducing E5 and the chronic cyanide risk. \emph{Non-structural measures}: (a) a \textbf{monitoring network} -- piezometers around the TSF, monthly sampling of the Lom River upstream and downstream for pH, CN$^{-}$, As, Hg and turbidity, with results transmitted to the regulatory board; (b) an \textbf{emergency preparedness and community-awareness plan} for Ndembo (alarm system, evacuation routes, alternative borehole water supply), plus a land-use planning measure: relocation of the TSF at least 300~m away from the mapped fault splays.

\textbf{6. Briefing (synthesis).} The site is workable but not as currently designed. Our recommendation to the board: \emph{authorise with conditions} -- relocation of the TSF away from the fault splays, double composite liner, spillway sized for the 2050 climate, slope drainage works, and a binding monitoring and emergency plan. This conclusion follows directly from the risk matrix: the three high risks (E2, E3, E5) are all reducible to moderate or low by the proposed measures.

\begin{attention}[Common mistakes in this type of exercise]
Confusing a \emph{hazard} (the dangerous phenomenon) with a \emph{risk} (hazard combined with exposure and vulnerability); scoring likelihood without referring to the data (return periods, climate trends); proposing mitigation measures without saying which risk they reduce; and forgetting that pollution needs a \emph{pathway} -- a source alone does not create a pollution risk if no pathway connects it to a receptor.
\end{attention}

\begin{savaistu}[Why this matters in Cameroon]
Cameroon has painful experience of neglected environmental geology: the Lake Nyos gas disaster of 1986, repeated deadly landslides on the slopes of Mount Cameroon and in the West Region, and mercury and cyanide pollution around artisanal gold sites in the East and Adamawa Regions. Environmental and Social Impact Assessments are now mandatory for mining projects in Cameroon -- the reasoning you have just practised is exactly what a professional geologist contributes to them.
\end{savaistu}

\newpage

\coursec{Methods and worked examples}

\coursec{Environmental Geology}

\courssub{Definition and Scope of Environmental Geology}

\begin{definition}[Environmental Geology]
Environmental Geology is the application of geological knowledge to the identification, assessment, and mitigation of geological hazards, the management of natural resources (water, soil, minerals), and the understanding of human-induced changes to the Earth system (pollution, climate change, land degradation). It bridges the gap between Earth processes and human society.
\end{definition}

\begin{aretenir}[Key concepts]
\begin{itemize}
\item \cle{Hazard}: a potentially damaging physical event, phenomenon, or human activity (e.g., earthquake, lava flow, flood, gas release).
\item \cle{Risk}: the probability of harmful consequences (deaths, injuries, property loss, economic disruption) resulting from the interaction between a hazard and vulnerable conditions. Risk = Hazard $\times$ Vulnerability $\times$ Exposure.
\item \cle{Vulnerability}: the degree to which a system, community, or asset is susceptible to the damaging effects of a hazard.
\item \cle{Exposure}: the people, property, infrastructure, or other elements present in hazard zones that are subject to potential losses.
\item \cle{Mitigation}: measures taken to reduce the severity or likelihood of a hazard's impact (structural and non-structural).
\item \cle{Endogenous hazards}: driven by Earth's internal energy (volcanism, earthquakes, limnic eruptions).
\item \cle{Exogenous hazards}: driven by external agents (gravity, water, wind, climate) — mass movements, floods, erosion, drought.
\end{itemize}
\end{aretenir}

\begin{savaistu}[Cameroon context]
Cameroon sits at the junction of the Congo Craton, the Benue Trough, and the Cameroon Volcanic Line (CVL). This tectonic setting makes it a natural laboratory for environmental geology: active volcanism (Mount Cameroon, Lake Nyos), seismic zones (along the CVL and the Benue Trough), intense tropical weathering leading to landslides (Buea, Bamenda, Dschang), seasonal flooding (Logone plain, Wouri estuary), and growing pressure on groundwater resources (Yaoundé, Douala, Far North).
\end{savaistu}

\courssub{Skill 1 — Identifying and Classifying a Geohazard}

\begin{methode}[How to identify and classify a geologic hazard]
When faced with a photograph, a map, a news extract, or a field description, proceed in order:
\begin{enumerate}
\item \textbf{Describe the observable facts only}: what moved, what was emitted, what was damaged, over what area and time span.
\item \textbf{Identify the geological agent}: gravity (mass movement), internal energy (volcanism, seismicity), water (flood, coastal erosion), wind, or chemical release (gas burst).
\item \textbf{Classify the hazard} using the two great families: \cle{endogenous hazards} (earthquakes, volcanic eruptions, limnic eruptions) driven by Earth's internal energy, and \cle{exogenous hazards} (landslides, floods, erosion, drought) driven by external agents, mainly water and gravity.
\item \textbf{Distinguish hazard from risk}: the \cle{hazard} is the dangerous phenomenon itself; the \cle{risk} = hazard $\times$ vulnerability $\times$ exposed stakes (population, buildings, crops).
\item \textbf{Conclude with the trigger}: many exogenous hazards need a triggering event (exceptional rainfall, an earthquake, undercutting of a slope by a road cut).
\end{enumerate}
\textbf{Reflex:} never write ``the hazard is the deaths'' — deaths belong to the \emph{risk}, not the hazard.
\end{methode}

\begin{exoresolu}[Classifying Cameroonian events]
Classify each of the following events as an endogenous or exogenous geohazard, name the geological agent, and state one possible trigger: (a) the 1986 Lake Nyos disaster; (b) the 1999 eruption of Mount Cameroon; (c) the repeated landslides on the slopes of Buea during the rainy season; (d) seasonal flooding of the Logone plain around Kousseri.
\tcblower
\textbf{(a) Lake Nyos, 1986.} A sudden release of a dense cloud of \ce{CO2} that had accumulated at the bottom of the crater lake flowed down the valleys and asphyxiated $\sim$1,700 people and thousands of livestock. The gas is of \emph{magmatic origin}, slowly dissolved into the deep lake water. This is an \cle{endogenous hazard} (a \emph{limnic eruption}), the agent being chemically charged deep water overturning. Trigger: a disturbance of the lake's stratification (landslide into the lake, heavy rain cooling the surface, or internal convection) causing the deep gas-rich water to rise and degas violently.

\textbf{(b) Mount Cameroon, 1999.} Emission of lava flows and ash from flank fissures on the south-west flank. This is an \cle{endogenous hazard} (volcanic eruption); the agent is magma rising along fractures of the Cameroon Volcanic Line. Trigger: injection of new magma into the shallow plumbing system, often announced by swarms of small earthquakes.

\textbf{(c) Buea landslides.} Downslope movement of weathered volcanic soils and ash on steep slopes. This is an \cle{exogenous hazard} (mass movement); the agents are gravity and water. Trigger: intense rainy-season rainfall saturating the soil, increasing its weight and reducing cohesion, often aggravated by deforestation and road cuts.

\textbf{(d) Logone floods.} Overflow of the Logone River over its flat alluvial plain. \cle{Exogenous hazard}; agent: surface water. Trigger: exceptional seasonal rainfall in the upstream basin combined with the very low gradient of the plain.

\textbf{Marking reflex:} for each event, one mark for the family (endogenous/exogenous), one for the agent, one for a plausible trigger.
\end{exoresolu}

\begin{exercice}[Practice classification]
Classify the following as endogenous or exogenous hazards, name the agent, and give a trigger: (i) The 1984 Lake Monoun gas burst; (ii) Coastal erosion at Limbe; (iii) The 2000 Mount Cameroon eruption; (iv) Gully erosion in the Adamawa plateau.
\end{exercice}

\courssub{Skill 2 — Reading a Volcanic Hazard Map and Proposing Land-Use Zoning}

\begin{methode}[How to analyse a hazard map and propose zoning]
\begin{enumerate}
\item \textbf{Read the legend first}: identify the hazard zones (lava-flow paths, ash-fall sectors, lahar valleys, pyroclastic density current zones) and their ranking (high, moderate, low).
\item \textbf{Locate the source} (summit crater, flank fissures, river) and trace the \emph{preferential paths}: lava and lahars follow valleys and the steepest slopes; ash fall follows the prevailing wind direction; pyroclastic flows follow topography but can overtop ridges.
\item \textbf{Superimpose the stakes}: towns, roads, farmland, water intakes, hospitals, schools.
\item \textbf{Propose zoning} in three bands: \emph{red zone} (no construction, evacuation plan), \emph{orange zone} (restricted building, reinforced structures, monitoring), \emph{green zone} (normal use with awareness).
\item \textbf{Justify each decision} by linking the physical behaviour of the hazard (flow velocity, deposit thickness, frequency, temperature) to the vulnerability of the stake.
\end{enumerate}
\textbf{Key idea:} mitigation by \cle{land-use planning} is the cheapest and most durable measure — it acts \emph{before} the event.
\end{methode}

\begin{popfigure}
\begin{tikzpicture}[scale=0.8]
% Simplified hazard map of SW Mount Cameroon
\draw[popdark, thick] (0,0) -- (10,0) -- (10,8) -- (0,8) -- cycle; % map frame
% Summit area
\draw[poporange, fill=poporange!20] (5,7) ellipse (1.2 and 0.8);
\node[popdark] at (5,7) {Summit};
% Lava flow channels
\draw[popink, line width=3pt] (5,6.2) .. controls (4.5,4) and (3,2) .. (2,0.5);
\draw[popink, line width=3pt] (5,6.2) .. controls (5.5,4) and (7,2) .. (8,0.5);
\node[popink, rotate=-30] at (3.5,3) {Lava channel 1};
\node[popink, rotate=30] at (6.5,3) {Lava channel 2};
% Ash fall sector (prevailing wind from SW to NE)
\draw[poppurple, fill=poppurple!10] (5,7) -- (10,8) -- (10,0) -- (7,0) -- cycle;
\node[poppurple] at (8.5,4) {Ash fall sector};
% Buea town
\draw[popblue, fill=popblue!20] (3.5,4.5) rectangle (5,5.5);
\node[popblue] at (4.25,5) {Buea};
% Idenau coastal road
\draw[popgreen, line width=2pt] (1.5,0.5) -- (2.5,0.5);
\node[popgreen] at (2,0.2) {Idenau road};
% Legend
\begin{scope}[shift={(0.5,0.5)}]
\draw[poporange, fill=poporange!20] (0,0) rectangle (0.5,0.5); \node[right] at (0.6,0.25) {Summit / vents};
\draw[popink, line width=3pt] (0, -0.7) -- (0.5, -0.7); \node[right] at (0.6,-0.7) {Lava flow paths};
\draw[poppurple, fill=poppurple!10] (0, -1.4) rectangle (0.5, -0.9); \node[right] at (0.6,-1.15) {Ash fall zone};
\draw[popblue, fill=popblue!20] (0, -2.1) rectangle (0.5, -1.6); \node[right] at (0.6,-1.85) {Town / stakes};
\draw[popgreen, line width=2pt] (0, -2.8) -- (0.5, -2.8); \node[right] at (0.6,-2.8) {Main road};
\end{scope}
% Scale and north arrow
\draw[->, thick] (9,7) -- (9,7.5); \node[above] at (9,7.5) {N};
\draw[|-|, thick] (8,0.3) -- (9,0.3); \node[below] at (8.5,0.1) {5 km};
\end{tikzpicture}
\legende{Simplified volcanic hazard map of the south-west flank of Mount Cameroon (schematic).}
\end{popfigure}

\begin{exoresolu}[Zoning around Mount Cameroon]
A simplified hazard map of the south-west flank of Mount Cameroon shows: the summit craters; two historical lava-flow channels descending towards the sea near Idenau; a wide sector east of the mountain exposed to ash fall under the prevailing wind; the town of Buea on the south-east slope at about \SI{1000}{m} altitude; and fertile volcanic soils densely farmed. Propose a three-band land-use zoning and justify it.
\tcblower
\textbf{Step 1 — Source and paths.} Source: summit and flank fissures. Lava follows the two mapped channels (steepest descent, valley-guided). Ash fall affects the eastern sector (downwind, prevailing SW winds). Lahars can remobilise ash and loose pyroclastic deposits in river valleys during the rains.

\textbf{Step 2 — Stakes.} Buea (large population, university, administrative centre), farms on rich volcanic soils, the coastal road near Idenau, water catchments.

\textbf{Step 3 — Zoning proposal.}
\begin{itemize}
\item \textbf{Red zone (prohibited):} the two lava-flow channels and the upper slopes above Buea (within \SI{2}{km} of active fissure zones). Lava flows there are frequent on geological timescales (1999, 2000, 2012), fast relative to evacuation in the upper reaches, and unstoppable. No housing, no schools, no critical infrastructure; only monitored scientific or pastoral use with rapid evacuation capability.
\item \textbf{Orange zone (restricted):} the lower flanks around Buea and the main ash-fall sector. Building allowed with strict rules: roofs designed to support \SI{>150}{kg.m^{-2}} ash load, water tanks covered, evacuation routes marked and maintained, sirens and a monitoring network (seismometers, gas sensors, GPS, tiltmeters) operational. Critical facilities (hospitals, schools) must meet higher structural standards.
\item \textbf{Green zone:} areas outside flow channels and away from the main ash axis (e.g., western flank beyond the lava channels), where normal development is permitted with public awareness campaigns and inclusion in regional evacuation plans.
\end{itemize}

\textbf{Step 4 — Justification.} The zoning matches the \emph{physical behaviour} of each hazard: lava is channelled and lethal but predictable in path, so exclusion is possible and cheap; ash fall is widespread but rarely lethal, so adaptation (roof design, water protection) is proportionate; farming may continue in orange zones because soil fertility is a strong stake, provided evacuation plans exist. This illustrates the principle that risk reduction combines hazard knowledge with management of vulnerability and exposure.
\end{exoresolu}

\courssub{Skill 3 — Analysing a Landslide: Factors, Mechanism, Stabilisation}

\begin{methode}[How to reason about slope stability]
\begin{enumerate}
\item \textbf{Draw a sketch cross-section}: slope angle, layers (soil, weathered rock, bedrock), water table, human works (road cut, building).
\item \textbf{List the driving forces} (weight of the mass $W$, increased by water saturation and added loads; downslope component $W \sin\alpha$) and the \textbf{resisting forces} (cohesion $c$, friction $\tan\phi$, roots, retaining structures; shear strength $\tau = c + \sigma' \tan\phi$).
\item \textbf{Identify the mechanism}: fall (free fall of blocks), slide (planar or rotational along a curved slip surface), or flow (mudflow/debris flow when saturated).
\item \textbf{Find the trigger}: rainfall, earthquake shaking, undercutting at the toe, leakage from pipes, rapid drawdown.
\item \textbf{Propose stabilisation} by acting on the balance: reduce driving forces (drainage, regrading the slope, removing load at the head) or increase resistance (retaining wall or piles at the toe, reforestation, soil nailing, avoiding water infiltration).
\end{enumerate}
\textbf{Key formula idea:} stability is expressed by a \cle{factor of safety} $F = \dfrac{\text{resisting forces}}{\text{driving forces}}$; sliding begins when $F < 1$. For an infinite slope: $F = \dfrac{c' + (\gamma z \cos^2\alpha - u) \tan\phi'}{\gamma z \sin\alpha \cos\alpha}$ where $u$ is pore water pressure.
\end{methode}

\begin{popfigure}
\begin{tikzpicture}[scale=0.9]
% Landslide cross-section: road cut in weathered basalt
\draw[popdark, thick] (0,0) -- (10,0); % original ground
\draw[popdark, thick] (0,0) -- (2,4) -- (8,4) -- (10,0); % topography
% Road cut
\draw[popmuted, fill=popmuted!30] (4,4) -- (4,2) -- (6,2) -- (6,4); % road bench
\draw[popmuted, thick] (4,2) -- (6,2); % road surface
% Layers
\draw[poporange!50!popdark, fill=poporange!20] (0,0) -- (10,0) -- (8,4) -- (2,4) -- cycle; % bedrock
\draw[poporange, fill=poporange!40] (0,0) -- (10,0) -- (9.5,0.5) -- (0.5,0.5) -- cycle; % weathered zone (thin)
% Water table
\draw[popblue, dashed, thick] (1,1) -- (9,1);
\node[popblue] at (9.5,1) {Water table};
% Slip surface (rotational)
\draw[popink, thick, dashed] (3.5,4) .. controls (4,2.5) and (6,2.5) .. (6.5,4);
\node[popink, rotate=15] at (5,3.2) {Slip surface};
% Forces
\draw[->, thick, popdark] (5,3.5) -- (5,2.5) node[midway, right] {$W$};
\draw[->, thick, popgreen] (5,2.5) -- (6.5,2.5) node[midway, above] {$W\sin\alpha$};
\draw[->, thick, popblue] (5,2.5) -- (5,1.5) node[midway, right] {$W\cos\alpha$};
\draw[<->, thick, poppurple] (4.5,3.8) -- (5.5,3.8) node[midway, above] {$\tau_{resist}$};
% Road label
\node[popmuted] at (5,2.2) {Road};
% Slope angle
\draw (2,4) -- (3,4) arc (0:63:1) node[midway, right] {$\alpha$};
\end{tikzpicture}
\legende{Cross-section of a rotational landslide in a road cut through weathered basalt. The road cut removes toe support; rainfall raises the water table, increasing pore pressure $u$ and weight $W$, reducing the factor of safety $F$.}
\end{popfigure}

\begin{exoresolu}[The road-cut landslide]
During the rainy season, a section of road cut into a steep hillside of weathered basalt near Bamenda is repeatedly buried by slides of soil and blocks. (a) Explain the mechanism using a force balance. (b) Propose three engineering or biological mitigation measures, justifying each.
\tcblower
\textbf{(a) Mechanism.} The road cut has steepened the slope angle $\alpha$ and removed support at the \emph{toe} (the resisting force provided by the downslope material). The weathered basalt and overlying soil have low effective cohesion $c'$ and friction angle $\phi'$. During heavy rain, water infiltrates: it \emph{increases the weight} of the mass (driving force $W\sin\alpha$ up) and \emph{decreases effective normal stress} ($\sigma' = \sigma - u$) along the potential slip surface via pore-water pressure $u$, thereby reducing shear resistance $\tau = c' + \sigma' \tan\phi'$. The factor of safety
\[ F = \frac{\text{resisting forces (cohesion + friction)}}{\text{driving forces (downslope component of weight)}} \]
falls below 1, and the mass slides along a curved (rotational) slip surface, burying the road.

\textbf{(b) Mitigation.}
\begin{enumerate}
\item \textbf{Drainage}: dig intercepting ditches upslope and horizontal drains in the slope to lower the water table. Justification: less water means less weight and lower pore pressure $u$, so effective stress $\sigma'$ rises and $F$ increases.
\item \textbf{Toe support}: build a retaining wall or gabion wall at the foot of the cut, keyed into competent material. Justification: it adds a resisting force exactly where the slip surface daylights, replacing the removed buttress.
\item \textbf{Regrading and reforestation}: reduce the slope angle (benching to $\alpha_{new} < \alpha_{crit}$) and plant deep-rooted vegetation (e.g., vetiver grass, \emph{Acacia}, \emph{Eucalyptus}). Justification: regrading reduces the driving force $W\sin\alpha$; roots add apparent cohesion $c_{roots}$ and vegetation pumps water out by evapotranspiration, lowering $u$.
\end{enumerate}
\textbf{Exam reflex:} always link each measure to \emph{which term of the balance} it modifies — this is what earns full marks.
\end{exoresolu}

\begin{attention}[Common mistakes in landslide answers]
\begin{itemize}
\item Confusing the \emph{trigger} (rain) with the \emph{underlying cause} (steep slope, weak material, lack of drainage).
\item Writing ``plant trees'' without specifying \emph{deep-rooted} species or explaining \emph{how} roots add cohesion.
\item Forgetting that drainage must be \emph{maintained}; clogged drains worsen the problem.
\item Not drawing a cross-section — examiners expect a labelled sketch showing slip surface, water table, and forces.
\end{itemize}
\end{attention}

\courssub{Skill 4 — Explaining the Greenhouse Effect and Interpreting a Temperature Record}

\begin{methode}[How to explain global warming rigorously]
\begin{enumerate}
\item \textbf{Start from the energy balance}: the Earth receives short-wave solar radiation (\SI{\sim 340}{W.m^{-2}} average at top of atmosphere) and re-emits long-wave (infrared) radiation to space.
\item \textbf{Role of greenhouse gases} (\ce{CO2}, \ce{CH4}, water vapour, \ce{N2O}, CFCs): transparent to short waves but \emph{absorbing infrared} at specific wavelengths, they trap part of the outgoing radiation — the natural \cle{greenhouse effect}, without which the mean surface temperature would be \SI{\sim -18}{\ensuremath{{}^{\circ}\mathrm{C}}} instead of \SI{\sim +15}{\ensuremath{{}^{\circ}\mathrm{C}}}.
\item \textbf{The human perturbation}: burning fossil fuels, deforestation, cement production, and some agriculture increase \ce{CO2} and \ce{CH4} concentrations, \emph{amplifying} the natural effect (radiative forcing).
\item \textbf{Interpret a graph}: describe the trend (direction, rate, acceleration), identify anomalies (volcanic dips, El Niño peaks), and \emph{correlate} with \ce{CO2} curves — but state clearly that correlation plus a known physical mechanism (radiative transfer) is what supports causation.
\item \textbf{Conclude with consequences}: sea-level rise (thermal expansion + glacier/ice-sheet melt), shifting rainfall patterns, more frequent extremes (heatwaves, intense rainfall), stress on water resources and agriculture, ocean acidification.
\end{enumerate}
\textbf{Reflex:} never write that the greenhouse effect itself is bad — it is its \emph{amplification} that is the problem.
\end{methode}

\begin{popfigure}
\begin{tikzpicture}[scale=0.7]
\begin{axis}[
    width=12cm, height=7cm,
    xlabel={Year}, ylabel={Temperature anomaly (\ensuremath{{}^{\circ}\mathrm{C}})},
    xmin=0, xmax=145,
    ymin=-0.5, ymax=1.3,
    xtick={0,20,40,60,80,100,120,140},
    xticklabels={1880,1900,1920,1940,1960,1980,2000,2020},
    grid=major, grid style={popline},
    legend pos=north west,
    tick label style={font=\small},
    label style={font=\small},
]
\addplot[popblue, thick, smooth] coordinates {
    (0,-0.2) (20,-0.1) (40,-0.2) (60,0.1) (80,-0.05) (100,0.1) (120,0.4) (140,1.0) (143,1.2)
};
\addlegendentry{Global mean surface temperature anomaly (relative to 1951-1980)}
\end{axis}
% Second axis for CO2
\begin{axis}[
    width=12cm, height=7cm,
    xmin=0, xmax=145,
    ymin=280, ymax=430,
    yticklabel pos=right,
    ylabel={CO2 (ppm)},
    axis y line=right,
    axis x line=none,
    hide x axis,
    tick label style={font=\small},
    label style={font=\small},
    legend pos=north east,
]
\addplot[popink, thick, smooth, dashed] coordinates {
    (0,290) (20,295) (40,300) (60,310) (80,320) (100,340) (120,370) (140,415) (143,420)
};
\addlegendentry{Atmospheric \ce{CO2} (ppm)}
\end{axis}
\end{tikzpicture}
\legende{Schematic evolution of global mean surface temperature anomaly and atmospheric \ce{CO2} concentration since 1880.}
\end{popfigure}

\begin{exoresolu}[Interpreting a climate graph]
A graph shows global mean surface temperature anomaly rising by roughly \SI{1.2}{\ensuremath{{}^{\circ}\mathrm{C}}} between 1880 and today, while atmospheric \ce{CO2} rose from about 290 to over 420 parts per million over the same period. (a) Explain physically why rising \ce{CO2} raises temperature. (b) Give two consequences relevant to Cameroon. (c) State two mitigation measures and two adaptation measures.
\tcblower
\textbf{(a) Physical explanation.} Solar energy arrives mainly as visible (short-wave) radiation; the warmed surface re-emits infrared (long-wave) radiation. \ce{CO2} molecules absorb specific infrared wavelengths (vibrational-rotational transitions) and re-emit them in all directions, part back towards the surface. Adding \ce{CO2} therefore reduces the fraction of energy escaping directly to space (increases the effective emission altitude); to restore the top-of-atmosphere energy balance, the surface and lower atmosphere warm. This is the \emph{amplification of the natural greenhouse effect} (radiative forcing $\Delta F \approx 5.35 \ln(C/C_0)\ \mathrm{W\,m^{-2}}$).

\textbf{(b) Consequences for Cameroon.}
\begin{itemize}
\item \textbf{Rainfall disruption}: shifts in the onset, cessation, and length of the rainy season threaten rain-fed agriculture (maize, cocoa, plantain, coffee) and increase drought risk in the Far North, where the shrinkage of Lake Chad (from \SI{\sim 25000}{km^2} in the 1960s to \SI{<2000}{km^2} recently) already illustrates regional water stress compounded by climate variability.
\item \textbf{Sea-level rise and coastal erosion}: thermal expansion of seawater and melting land ice raise global mean sea level (\SI{\sim 3.3}{mm.yr^{-1}} currently), threatening low-lying coastal zones around Douala (Wouri estuary), Limbe, and Kribi (flooding, salinisation of soils and aquifers, mangrove loss, infrastructure damage).
\end{itemize}

\textbf{(c) Mitigation and adaptation.}
\begin{itemize}
\item \textbf{Mitigation (reduce causes)}: (1) Reduce greenhouse-gas emissions: renewable energy (hydropower potential on Sanaga, solar in North), energy efficiency, cleaner transport, reduced gas flaring. (2) Protect and restore \cle{carbon sinks}: reforestation, sustainable forest management in the Congo Basin, fight against deforestation, since forests store carbon in biomass and soils.
\item \textbf{Adaptation (manage impacts)}: (1) Climate-smart agriculture: drought-resistant varieties, agroforestry, irrigation efficiency. (2) Coastal zone management: mangrove restoration as natural buffers, setback lines for construction, early warning for storm surges.
\end{itemize}
\textbf{Marking reflex:} one mark for the radiation mechanism correctly told in terms of short wave vs infrared; consequences must be \emph{linked} to the mechanism, not listed bare; distinguish mitigation from adaptation.
\end{exoresolu}

\courssub{Skill 5 — Tracing a Pollutant from Source to Water Table}

\begin{methode}[How to analyse groundwater pollution]
\begin{enumerate}
\item \textbf{Identify the source}: point source (fuel tank, latrine, dump, mine tailings, industrial effluent) or diffuse source (fertilisers, pesticides over fields, atmospheric deposition).
\item \textbf{Identify the pollutant's nature}: soluble conservative (nitrates, chloride), soluble reactive (heavy metals, ammonium), immiscible liquid (LNAPL: hydrocarbons lighter than water; DNAPL: chlorinated solvents denser than water), particulate, pathogenic (bacteria, viruses).
\item \textbf{Follow the path}: infiltration through the \emph{unsaturated zone} (vadose zone) where soil may filter, degrade, or retard some pollutants (sorption, biodegradation, volatilisation) down to the \cle{water table}, then lateral flow in the \emph{saturated zone} (aquifer) towards wells, springs, or rivers (discharge zones).
\item \textbf{Assess the vulnerability of the aquifer}: use criteria like DRASTIC (Depth to water, Recharge, Aquifer media, Soil media, Topography, Impact of vadose zone, Conductivity). A shallow water table under permeable sandy soil with high recharge is highly vulnerable; a deep aquifer protected by a clay layer (aquitard) is much less so.
\item \textbf{Propose management}: remove the source, protect wellhead zones (sanitary protection perimeters), treat (pump-and-treat, permeable reactive barriers, monitored natural attenuation), and monitor with a network of piezometers.
\end{enumerate}
\textbf{Key idea:} clay protects (low $K$, high sorption), sand transmits (high $K$); depth and travel time allow natural attenuation (dilution, dispersion, degradation).
\end{methode}

\begin{popfigure}
\begin{tikzpicture}[scale=0.8]
% Groundwater pollution cross-section
\draw[popdark, thick] (0,0) -- (12,0); % ground surface
% Topography
\draw[popdark] (0,0) -- (2,1) -- (5,1.5) -- (8,1) -- (12,0);
% Layers
\draw[poporange!50!popdark, fill=poporange!10] (0,0) -- (12,0) -- (11.5,-1) -- (0.5,-1) -- cycle; % topsoil
\draw[poporange, fill=poporange!30] (0,-1) -- (12,-1) -- (11.5,-4) -- (0.5,-4) -- cycle; % sandy unsaturated zone
\draw[popblue, fill=popblue!10] (0,-4) -- (12,-4) -- (11.5,-6) -- (0.5,-6) -- cycle; % aquifer (sand/gravel)
\draw[popgreen, fill=popgreen!20] (0,-6) -- (12,-6) -- (11.5,-7) -- (0.5,-7) -- cycle; % aquitard (clay)
\draw[poporange!50!popdark, fill=poporange!10] (0,-7) -- (12,-7) -- (11.5,-8) -- (0.5,-8) -- cycle; % deep aquifer
% Water table
\draw[popblue, dashed, thick] (1,-3.5) -- (11,-3.5);
\node[popblue] at (11.5,-3.5) {Water table};
% Source: diffuse agriculture
\draw[popgreen!50!popdark, fill=popgreen!20] (2,0.5) rectangle (5,1.5);
\node[popgreen!50!popdark] at (3.5,1) {Fertilised fields};
% Source: point (latrine)
\draw[popink, fill=popink!20] (7,0.5) circle (0.3);
\node[popink] at (7,1.2) {Latrine};
% Pollutant paths
\draw[->, popink, thick] (3.5,0.5) -- (3.5,-3.5) -- (3.5,-4.5);
\draw[->, popink, thick] (7,0.2) -- (7,-3.5) -- (8,-4.5);
% Well shallow
\draw[popmuted, thick] (4,-0.2) -- (4,-4);
\node[popmuted] at (4,-4.3) {Shallow well};
% Borehole deep
\draw[popmuted, thick] (9,-0.2) -- (9,-7.5);
\node[popmuted] at (9,-7.8) {Deep borehole};
% Flow direction
\draw[->, thick, popdark] (1,-4.2) -- (2,-4.2) node[midway, above] {Flow};
\draw[->, thick, popdark] (10,-4.2) -- (9,-4.2);
% Labels
\node at (0.5,-0.5) {Topsoil};
\node at (0.5,-2.5) {Sand (unsat.)};
\node at (0.5,-5) {Aquifer};
\node at (0.5,-6.5) {Clay (aquitard)};
\node at (0.5,-7.5) {Deep aquifer};
\end{tikzpicture}
\legende{Cross-section showing pollutant transport from diffuse (fertilisers) and point (latrine) sources through a shallow vulnerable aquifer to a shallow well, while a deep borehole is partially protected by an aquitard.}
\end{popfigure}

\begin{exoresolu}[Nitrates in village wells]
In an intensively farmed area of the West Region (e.g., around Dschang or Bafoussam), several shallow hand-dug wells show rising nitrate concentrations (\SI{>50}{mg.L^{-1}}, exceeding WHO guideline). The subsoil is sandy saprolite over granitic bedrock and the water table lies only \SI{3-5}{m} deep. (a) Explain the origin and the path of the pollution. (b) Why are deeper boreholes in the same area less affected? (c) Propose three management measures.
\tcblower
\textbf{(a) Origin and path.} The source is \emph{diffuse}: nitrogen fertilisers (urea, NPK) spread on fields, plus possibly pit latrines and animal waste. Rain and irrigation water dissolve nitrates (\ce{NO3-}, highly soluble, anionic, not retained by negatively charged clay minerals) and carry them downwards through the permeable sandy unsaturated zone — sand filters poorly, offers little sorption or denitrification — to the shallow water table, which the wells tap directly. The short travel time (months to a few years) prevents significant attenuation.

\textbf{(b) Deeper boreholes.} Deep boreholes often tap a deeper fractured aquifer separated from the shallow saprolite aquifer by a less permeable layer (weathered clay-rich horizon or intact saprock) acting as an \emph{aquitard}. The longer flow path and travel time (decades) also allow dilution, dispersion, and some natural denitrification (if organic carbon and anoxic conditions exist). Hence lower nitrate levels — though protection is relative, not absolute, especially if the aquitard is fractured or bypassed by the borehole annulus.

\textbf{(c) Management.}
\begin{enumerate}
\item \textbf{Act on the source}: rationalise fertiliser use (right dose, right period — split application, not before heavy rain), promote organic farming, and relocate latrines away from and downslope of wells (minimum \SI{30}{m}).
\item \textbf{Protect the resource}: establish a sanitary protection perimeter around wells (radius \SI{15-30}{m}: no farming, no latrines, no dumping, no chemical storage) and favour deeper, properly cased and grouted boreholes with sealed wellheads.
\item \textbf{Monitor and treat}: regular nitrate analyses of wells (quarterly); abandon or treat (e.g., blending with low-nitrate water, ion exchange, reverse osmosis) any well exceeding drinking-water guidelines; educate users on health risks (methaemoglobinaemia in infants).
\end{enumerate}
\textbf{Exam reflex:} structure the answer as \emph{source -- transport -- receptor -- remedy}; examiners reward this logical chain.
\end{exoresolu}

\courssub{Skill 6 — Building a Complete Mitigation Plan for a Geohazard}

\begin{methode}[How to structure a mitigation answer (synthesis skill)]
For any hazard (earthquake, eruption, landslide, flood, gas release), organise the answer along the \cle{risk-management cycle}:
\begin{enumerate}
\item \textbf{Prevention / prediction (knowledge \& monitoring)}: geological mapping, hazard zoning, monitoring networks (seismometers, GPS, gas sensors, river gauges, piezometers), early-warning systems.
\item \textbf{Protection (structural / engineering)}: retaining walls, drainage, check dams, lava-diversion barriers, earthquake-resistant building codes, flood embankments, degassing pipes (Lake Nyos).
\item \textbf{Preparedness (non-structural / organisational)}: land-use planning, evacuation plans, drills, education of the population, insurance, emergency kits, communication protocols.
\item \textbf{Response and recovery}: rescue organisation, relocation, reconstruction to better standards (build back better), psychological support.
\end{enumerate}
\textbf{Reflex:} a top-mark answer always mixes \emph{structural} and \emph{non-structural} measures and mentions \emph{monitoring} as the backbone of early warning.
\end{methode}

\begin{exoresolu}[Synthesis: protecting a town on the flank of Mount Cameroon]
A growing town lies on the lower flank of Mount Cameroon, exposed to lava flows, ash fall, earthquakes, and rainy-season landslides. As a geological adviser, propose a complete mitigation plan, indicating for each measure the hazard it addresses.
\tcblower
\textbf{1. Knowledge and monitoring (all hazards).} Produce a multi-hazard map from geological and historical records (past lava channels, landslide scars, ash isopachs, seismic catalogues). Install a monitoring network: seismometers to detect magma movement and tectonic earthquakes, GPS and tiltmeters for ground deformation, gas sensors for \ce{CO2} and \ce{SO2} (especially in low-lying areas), rain gauges and piezometers on identified unstable slopes. Link the network to an early-warning centre with sirens, SMS alerts, and community radio.

\textbf{2. Structural protection.}
\begin{itemize}
\item \emph{Lava flows}: keep historical flow channels free of buildings; where feasible, construct diversion barriers (earthen berms, concrete walls) on the upper slopes to deflect slow-moving flows away from the town (as tested at Etna and attempted at Mount Cameroon in 1999/2000).
\item \emph{Landslides}: drainage ditches (interceptor and French drains) and retaining structures (gabions, soil nailing) on identified unstable slopes; reforestation of steep catchments with deep-rooted species.
\item \emph{Earthquakes}: enforce earthquake-resistant building codes (light, tied structures; reinforced concrete frames with ductile detailing; no heavy unreinforced masonry on steep slopes).
\item \emph{Ash fall}: roof-slope standards (\SI{>30}{\ensuremath{{}^{\circ}}}) so ash slides off, covered water storage, gutter guards.
\end{itemize}

\textbf{3. Non-structural measures.} Land-use zoning (red/orange/green bands) controlling new construction; marked evacuation routes and identified safe assembly areas (outside flow paths); regular drills in schools and workplaces; public education in local languages (Pidgin, Bakweri, French); micro-insurance schemes for farmers; integration of hazard maps into the municipal urban plan (Plan d'Urbanisme Directeur).

\textbf{4. Response and recovery.} A crisis plan defining responsibilities (council, civil protection, scientists, Red Cross); pre-positioned supplies (water, food, medicine, shelter); after an event, rapid damage mapping (drones, satellite) and reconstruction only outside red zones, to better standards.

\textbf{Conclusion.} Because the hazards differ in behaviour (lava is channelled, ash is widespread, landslides are rain-triggered, earthquakes are sudden), no single measure suffices: effective mitigation \emph{layers} monitoring, engineering, planning, and education. This integrated approach, adapted to local means and governance, is the model answer expected at the GCE.
\end{exoresolu}

\begin{experience}[Case study: Lake Nyos degassing — an integrated mitigation success]
\textbf{Context:} After the 1986 disaster, an international team (French, Japanese, Cameroonian) designed a mitigation system.
\textbf{Measures:}
\begin{enumerate}
\item \textbf{Monitoring}: permanent raft with \ce{CO2} sensors, temperature profilers, weather station.
\item \textbf{Engineering (structural)}: installation of \textbf{degassing pipes} (initially one in 2001, two more in 2011) reaching the lake bottom. Water rises by natural buoyancy as gas exsolves, creating a self-sustaining fountain that vents \ce{CO2} safely to the atmosphere.
\item \textbf{Non-structural}: evacuation plan for shoreline villages, sirens, simulation exercises, pasture management to keep cattle away from high-risk zones at night.
\end{enumerate}
\textbf{Result:} Deep-water \ce{CO2} concentration reduced by \SI{\sim 80}{\%}; risk of another catastrophic overturn drastically lowered. A model of science–policy–community collaboration.
\end{experience}

\begin{aretenir}[Chapter synthesis — Environmental Geology for GCE]
\begin{itemize}
\item Environmental Geology applies geological science to protect society.
\item \cle{Hazard} $\neq$ \cle{Risk}. Risk = Hazard $\times$ Vulnerability $\times$ Exposure.
\item \cle{Endogenous} (volcanoes, earthquakes, limnic eruptions) vs \cle{Exogenous} (landslides, floods, erosion, drought).
\item \textbf{Volcanic zoning}: red/orange/green based on physical behaviour of lava, ash, lahars.
\item \textbf{Landslides}: force balance ($F = \text{resist}/\text{drive}$); triggers (rain, quake, toe cut); mitigation acts on drainage, support, slope angle, vegetation.
\item \textbf{Global warming}: amplified greenhouse effect (short-wave in, long-wave trapped by \ce{CO2}/\ce{CH4}); consequences for Cameroon (rainfall shifts, sea-level rise, Lake Chad); mitigation (emissions, sinks) vs adaptation (agriculture, coasts).
\item \textbf{Groundwater pollution}: source–path–receptor; vulnerability (DRASTIC); management (source control, protection perimeters, monitoring, treatment).
\item \textbf{Mitigation plan}: monitoring $\rightarrow$ structural $\rightarrow$ non-structural $\rightarrow$ response/recovery. Always integrate.
\end{itemize}
\end{aretenir}

\begin{exercice}[GCE-style structured questions]
\begin{enumerate}
\item (a) Define \emph{hazard} and \emph{risk}. Explain why a large earthquake in an uninhabited desert poses a hazard but negligible risk. (4 marks)
\item (b) The 1986 Lake Nyos event is often called a ``limnic eruption''. Explain the geological conditions necessary for such an event and why it is classified as an endogenous hazard. (6 marks)
\item (c) With the aid of a labelled cross-section, explain how a road cut can trigger a rotational landslide in weathered basalt during the rainy season. (8 marks)
\item (d) Describe the physical mechanism of the enhanced greenhouse effect. Give two observed consequences in Cameroon and one mitigation and one adaptation measure. (7 marks)
\item (e) A shallow well in a sandy area shows high nitrate levels. Explain the pollutant pathway from source to well. Why might a deeper borehole nearby be safer? Propose two management actions. (7 marks)
\end{enumerate}
\end{exercice}

\begin{corrige}[Exercise 1]
\textbf{(a)} Hazard: a potentially damaging physical event (e.g., earthquake). Risk: expected losses (lives, property) = Hazard $\times$ Vulnerability $\times$ Exposure. In an uninhabited desert, Exposure $\approx 0$, so Risk $\approx 0$ despite the Hazard. (4 marks: 1 for each definition, 2 for explanation).
\textbf{(b)} Conditions: deep crater lake, stable stratification (meromictic), magmatic \ce{CO2} input at bottom, high hydrostatic pressure keeping gas dissolved. Endogenous because the gas source is internal (magmatic) and the release mechanism involves deep Earth processes. (6 marks: 3 for conditions, 2 for endogenous classification, 1 for trigger mention).
\textbf{(c)} Cross-section showing: original slope, road cut removing toe, weathered basalt layer, water table rise, curved slip surface, forces $W$, $W\sin\alpha$, $W\cos\alpha$, resisting shear $\tau$. Explanation: cut increases $\alpha$ and removes buttress; rain raises $u$, reduces $\sigma'$, $F<1$. (8 marks: 3 for diagram, 5 for force balance explanation).
\textbf{(d)} Solar short-wave passes atmosphere; surface emits long-wave; \ce{CO2} absorbs/re-emits IR, trapping heat. Consequences: (1) Rainfall disruption affecting cocoa/maize; (2) Sea-level rise threatening Douala/Limbe. Mitigation: reforestation/renewables. Adaptation: drought-resistant crops/mangrove restoration. (7 marks: 3 for mechanism, 2 for consequences, 2 for measures).
\textbf{(e)} Path: fertiliser/latrine $\rightarrow$ infiltration through sand (low attenuation) $\rightarrow$ shallow water table $\rightarrow$ well. Deep borehole: protected by aquitard, longer travel time, dilution/denitrification. Management: (1) source control (fertiliser timing, latrine relocation); (2) wellhead protection perimeter + monitoring. (7 marks: 3 for pathway, 2 for borehole, 2 for management).
\end{corrige}

\newpage

\coursec{Exercises}

\courssub{I check my knowledge}

\begin{exercice}[Multiple choice questions]
For each question, choose the single correct answer.
\begin{enumerate}
\item Environmental geology is best defined as:
\begin{enumerate}
\item the study of rocks and minerals for mining purposes;
\item the application of geological knowledge to the interactions between humans and the geological environment, including hazards, resources and pollution;
\item the study of weather and climate only;
\item the branch of geology dealing exclusively with earthquakes.
\end{enumerate}
\item The factor of safety $F$ of a slope is defined as:
\begin{enumerate}
\item the ratio of driving forces to resisting forces;
\item the ratio of resisting forces to driving forces;
\item the product of cohesion and friction angle;
\item the slope angle divided by the pore-water pressure.
\end{enumerate}
\item A 475-year return period for an earthquake ground motion corresponds to an annual probability of exceedance of approximately:
\begin{enumerate}
\item 10\%; \quad (b) 2\%; \quad (c) 0.21\%; \quad (d) 50\%.
\end{enumerate}
\item Which of the following is a non-point (diffuse) source of water pollution?
\begin{enumerate}
\item a factory discharge pipe;
\item fertiliser runoff from agricultural fields;
\item a sewage treatment plant outlet;
\item an oil tanker leak at a single location.
\end{enumerate}
\end{enumerate}
\end{exercice}

\begin{exercice}[True or false — justify]
State whether each statement is true or false, and justify your answer in one or two sentences.
\begin{enumerate}
\item The 1986 Lake Nyos disaster was caused by a volcanic eruption of lava.
\item An increase in pore-water pressure within a slope decreases its stability.
\item Global warming is caused solely by natural variations in solar activity.
\item A contaminant with a first-order decay constant disappears more quickly from groundwater than a conservative (non-reactive) tracer.
\item Mitigation of geohazards can include both structural measures (engineering works) and non-structural measures (land-use planning, early warning).
\end{enumerate}
\end{exercice}

\begin{exercice}[Key definitions]
Define each of the following terms precisely, as expected at the GCE Advanced Level:
\begin{enumerate}
\item geohazard;
\item return period of a natural event;
\item peak ground acceleration (PGA);
\item greenhouse effect;
\item point source of pollution;
\item breakthrough curve (in contaminant hydrogeology).
\end{enumerate}
\end{exercice}

\begin{exercice}[Direct calculations]
\begin{enumerate}
\item An earthquake of magnitude 6 releases an energy $E_6$; an earthquake of magnitude 8 releases an energy $E_8$. Given that each unit of magnitude corresponds to a factor of about 32 in energy, calculate the ratio $E_8/E_6$.
\item A borehole monitoring network shows that a conservative tracer injected at a source well reaches an observation well \SI{200}{m} away in 50 days. Calculate the average linear groundwater velocity in \si{m/day} and in \si{m/year}.
\item Atmospheric \ce{CO2} concentration rises from 410 ppm to 550 ppm. Using the simplified climate sensitivity relation $\Delta T = \lambda \ln(C/C_0)$ with $\lambda = \SI{3.0}{\kelvin}$, estimate the resulting equilibrium temperature rise.
\end{enumerate}
\end{exercice}

\courssub{I practise}

\begin{exercice}[Factor of safety of a slope]
A hillslope near Dschang has a slope angle $\beta = 30^{\circ}$. A planar failure surface lies parallel to the slope at a vertical depth of \SI{3}{m}. The soil has unit weight $\gamma = \SI{19}{kN/m^3}$, cohesion $c' = \SI{10}{kPa}$, and effective friction angle $\varphi' = 25^{\circ}$. During the rainy season, pore-water pressure on the failure plane reaches $u = \SI{15}{kPa}$.
\begin{enumerate}
\item Using the infinite-slope formula
\[ F = \frac{c' + (\gamma z \cos^2\beta - u)\tan\varphi'}{\gamma z \sin\beta \cos\beta}, \]
calculate the factor of safety.
\item Interpret the result: is the slope stable?
\item Compute the new value of $F$ if drainage lowers $u$ to \SI{5}{kPa}, and comment on the effectiveness of drainage as a mitigation measure.
\end{enumerate}
\end{exercice}

\begin{exercice}[Reading a seismic hazard curve]
A seismic hazard assessment for a coastal Cameroonian town gives the following values of peak ground acceleration (PGA) as a function of return period:

\begin{center}
\begin{tabular}{|l|c|c|c|c|c|}
\hline
\rowcolor{popblueL} Return period (years) & 72 & 225 & 475 & 975 & 2475 \\
\hline
PGA ($g$) & 0.05 & 0.10 & 0.15 & 0.21 & 0.32 \\
\hline
\end{tabular}
\end{center}

\begin{enumerate}
\item Plot the hazard curve (PGA against return period, using a logarithmic scale for the return period).
\item Determine the design PGA for a 475-year return period.
\item Explain what a 475-year return period means in terms of annual probability of exceedance.
\item Discuss two implications of this design PGA for building codes in a town located near the Cameroon Volcanic Line.
\end{enumerate}
\end{exercice}

\begin{exercice}[Simple carbon-cycle box model]
Consider a two-box model of the climate system. The atmospheric \ce{CO2} concentration increases from $C_0 = 410$ ppm to $C = 550$ ppm due to emissions. The radiative forcing is approximated by $\Delta F = 5.35 \ln(C/C_0)$ (in \si{W/m^2}), and the equilibrium temperature response is $\Delta T = \lambda \Delta F$ with climate sensitivity parameter $\lambda = 0.8\ \mathrm{K/(W\,m^{-2})}$.
\begin{enumerate}
\item Calculate the radiative forcing $\Delta F$.
\item Calculate the equilibrium temperature rise $\Delta T$.
\item State two reasons why this simple model may underestimate or overestimate the real temperature response.
\item Give two effects of such a temperature rise on the climate of Cameroon.
\end{enumerate}
\end{exercice}

\begin{exercice}[Contaminant breakthrough curve]
A contaminant is released as a continuous source of concentration $C_0 = \SI{100}{mg/L}$ into a sandy aquifer. One-dimensional advective--dispersive transport occurs with average linear velocity $v = \SI{0.5}{m/day}$, longitudinal dispersion coefficient $D = \SI{1.0}{m^2/day}$, and first-order decay constant $\lambda_d = \SI{0.01}{day^{-1}}$. An observation well is located at $x = \SI{50}{m}$.
\begin{enumerate}
\item Calculate the advective travel time $t_a = x/v$ to the well.
\item Explain qualitatively the role of dispersion and of decay in shaping the breakthrough curve $C(t)$ at the well.
\item Sketch the expected breakthrough curve (relative concentration $C/C_0$ against time), indicating the approximate position of $t_a$ and the effect of decay on the plateau concentration.
\item Compute the pure advective estimate of arrival time if the decay constant were doubled, and state whether decay changes the arrival time or the concentration level.
\end{enumerate}
\end{exercice}

\courssub{I prepare for the GCE}

\begin{exercice}[Essay — geohazards in Cameroon (20 marks)]
\begin{enumerate}
\item Define the term \emph{geohazard} and distinguish clearly between a hazard, a disaster and vulnerability. \textbf{(6 marks)}
\item Describe the causes and effects of \emph{two} major natural geologic hazards that affect Cameroon, choosing among volcanic eruptions, earthquakes, landslides and toxic gas emissions from crater lakes. Support your answer with named Cameroonian examples. \textbf{(8 marks)}
\item Discuss three mitigation strategies (at least one structural and one non-structural) that can reduce the impact of geohazards on Cameroonian communities, and explain the limitations of each. \textbf{(6 marks)}
\end{enumerate}
\end{exercice}

\begin{exercice}[Structured question — global warming (20 marks)]
\begin{enumerate}
\item Explain the natural greenhouse effect and show how human activities enhance it. Name at least three greenhouse gases and give one major anthropogenic source of each. \textbf{(7 marks)}
\item Using the relation $\Delta T = \lambda \ln(C/C_0)$ with $\lambda = \SI{3.0}{\kelvin}$, calculate the equilibrium warming for a doubling of \ce{CO2} from 300 ppm to 600 ppm. \textbf{(3 marks)}
\item Describe four observed or projected effects of global warming on climate and on the physical environment (sea level, extreme events, water resources, ecosystems), illustrating with examples relevant to Africa or Cameroon. \textbf{(6 marks)}
\item Outline two mitigation and two adaptation measures that a country like Cameroon can implement, and briefly evaluate their feasibility. \textbf{(4 marks)}
\end{enumerate}
\end{exercice}

\begin{exercice}[Critical evaluation of an Environmental Impact Assessment (20 marks)]
A published Environmental Impact Assessment (EIA) for a proposed hydro-electric dam on a Cameroonian river contains the following elements: a description of the project; a baseline study of water quality carried out during one dry-season campaign only; a flora and fauna inventory; an economic analysis; and a brief statement that ``no significant geological risks exist in the region''. The dam site lies in a region of moderate seismicity, with steep deforested slopes upstream and a growing town downstream.
\begin{enumerate}
\item State the purpose and main stages of an Environmental Impact Assessment. \textbf{(4 marks)}
\item Identify and justify at least four gaps or weaknesses in this EIA, focusing on geohazard analysis (seismic risk, slope instability, reservoir-triggered landslides, sedimentation) and on climate-change considerations (rainfall variability, flood design, greenhouse-gas emissions from the reservoir). \textbf{(8 marks)}
\item For each gap identified, propose a concrete improvement (additional studies, monitoring, design or management measures). \textbf{(6 marks)}
\item Conclude by explaining why public participation and long-term monitoring are essential components of an acceptable EIA. \textbf{(2 marks)}
\end{enumerate}
\end{exercice}

\newpage

\coursec{Solutions to the exercises}

\begin{corrige}[Exercise 1 — Multiple choice questions]
\begin{enumerate}
\item \textbf{Answer: (b).} Environmental geology is precisely the application of geological knowledge and methods to the interactions between human society and the geological environment: geohazards, resource use, waste disposal and pollution. Option (a) describes economic geology, (c) meteorology, and (d) seismology only.
\item \textbf{Answer: (b).} By definition $F = \dfrac{\text{resisting (stabilising) forces}}{\text{driving (destabilising) forces}}$. A slope is stable when $F>1$, at limit equilibrium when $F=1$, and unstable when $F<1$. Option (a) is the inverse ratio, which would exceed 1 for unstable slopes — a classic inversion error.
\item \textbf{Answer: (c).} Annual probability of exceedance $P = 1/T = 1/475 \approx 0.0021$, i.e. about $0.21\%$ per year. (The 475-year event is often quoted as ``10\% probability in 50 years'', which can mislead candidates into choosing (a); note that $1/475 \neq 10\%$.)
\item \textbf{Answer: (b).} Fertiliser runoff from fields enters water bodies over a wide area with no single identifiable discharge point: it is a diffuse (non-point) source. Options (a), (c) and (d) are all point sources because the discharge location is identifiable.
\end{enumerate}
\end{corrige}

\begin{corrige}[Exercise 2 — True or false — justify]
\begin{enumerate}
\item \textbf{False.} The 1986 Lake Nyos disaster was a \emph{limnic eruption}: a sudden massive release of dissolved \ce{CO2} that had accumulated in the deep lake water, which asphyxiated people and animals in the surrounding valleys. No lava was erupted.
\item \textbf{True.} Pore-water pressure reduces the effective normal stress on a potential failure plane ($\sigma' = \sigma - u$), and hence reduces the frictional resistance $\tau_f = c' + \sigma'\tan\varphi'$; the factor of safety decreases.
\item \textbf{False.} Although solar variability and volcanic activity cause natural climate fluctuations, the warming observed since the industrial era is dominated by the enhanced greenhouse effect due to anthropogenic emissions of \ce{CO2}, \ce{CH4}, \ce{N2O} and halocarbons.
\item \textbf{True.} A decaying contaminant is progressively destroyed during transport (half-life $t_{1/2} = \ln 2/\lambda_d$), so its concentration falls below that of a conservative tracer transported at the same velocity.
\item \textbf{True.} Mitigation is broad: structural measures (retaining walls, drainage, seismic-resistant design, degassing pipes at Lake Nyos) and non-structural measures (hazard mapping, land-use zoning, early-warning systems, education) are complementary.
\end{enumerate}
\end{corrige}

\begin{corrige}[Exercise 3 — Key definitions]
\begin{enumerate}
\item \textbf{Geohazard:} a geological process or phenomenon (e.g. earthquake, volcanic eruption, landslide, limnic gas release) that poses a threat of damage to human life, property or the environment.
\item \textbf{Return period:} the average time interval $T$ between events of a given magnitude (or intensity) at a site; equivalently, the annual probability of exceedance is $P = 1/T$.
\item \textbf{Peak ground acceleration (PGA):} the maximum absolute value of ground acceleration recorded (or predicted) at a site during an earthquake, usually expressed as a fraction of $g$; it is the key parameter used for seismic design.
\item \textbf{Greenhouse effect:} the natural warming of the lower atmosphere caused by greenhouse gases which are transparent to incoming short-wave solar radiation but absorb and re-emit outgoing long-wave (infrared) radiation from the Earth's surface.
\item \textbf{Point source of pollution:} a single, identifiable, localised discharge of pollutants, such as a factory outfall pipe or a sewage treatment plant outlet.
\item \textbf{Breakthrough curve:} the graph of contaminant concentration $C(t)$ (often normalised as $C/C_0$) measured at an observation point downstream of a source, showing the arrival, rise and plateau of the contaminant plume as a function of time.
\end{enumerate}
\end{corrige}

\begin{corrige}[Exercise 4 — Direct calculations]
\begin{enumerate}
\item Each magnitude unit multiplies the energy released by about 32 (more precisely $10^{1.5} \approx 31.6$), and $8 - 6 = 2$ units:
\[ \frac{E_8}{E_6} = 32^{2} = 1024. \]
A magnitude 8 earthquake releases roughly \textbf{1000 times} more energy than a magnitude 6.
\item Average linear velocity:
\[ v = \frac{x}{t} = \frac{\SI{200}{m}}{50\ \text{days}} = \SI{4}{m/day}. \]
Converting: $v = 4 \times 365 = \SI{1460}{m/year}$.
\item 
\[ \Delta T = \lambda \ln\frac{C}{C_0} = 3.0 \times \ln\frac{550}{410} = 3.0 \times \ln(1.341) \approx 3.0 \times 0.294 \approx \SI{0.88}{\kelvin}. \]
The estimated equilibrium warming is about $\SI{0.9}{\kelvin}$.
\end{enumerate}
\end{corrige}

\begin{corrige}[Exercise 5 — Factor of safety of a slope]
\textbf{(a)} Compute each term with $\beta = 30^{\circ}$, $z = \SI{3}{m}$, $\gamma = \SI{19}{kN/m^3}$:
\begin{align*}
\gamma z &= 19 \times 3 = \SI{57}{kPa},\\
\cos^2 30^{\circ} &= 0.75, \quad \sin 30^{\circ}\cos 30^{\circ} = 0.433, \quad \tan 25^{\circ} = 0.466.\\
\text{Resisting stress} &= c' + (\gamma z\cos^2\beta - u)\tan\varphi' = 10 + (57 \times 0.75 - 15)(0.466)\\
&= 10 + (42.75 - 15)(0.466) = 10 + 27.75 \times 0.466 = 10 + 12.93 = \SI{22.93}{kPa}.\\
\text{Driving stress} &= \gamma z \sin\beta\cos\beta = 57 \times 0.433 = \SI{24.68}{kPa}.\\
F &= \frac{22.93}{24.68} \approx 0.93.
\end{align*}

\textbf{(b)} Since $F = 0.93 < 1$, the driving stress exceeds the resisting stress: the slope is \textbf{unstable} and failure is expected (or imminent) during the rainy season. This is typical of landslide triggering on cultivated slopes of the Western Highlands during heavy rains.

\textbf{(c)} With drainage, $u = \SI{5}{kPa}$:
\[ F = \frac{10 + (42.75 - 5)(0.466)}{24.68} = \frac{10 + 37.75 \times 0.466}{24.68} = \frac{10 + 17.59}{24.68} = \frac{27.59}{24.68} \approx 1.12. \]
Now $F > 1$: the slope becomes stable. \textbf{Conclusion:} drainage is a highly effective structural mitigation measure because lowering pore-water pressure restores effective stress and frictional resistance; here it raises $F$ from 0.93 to 1.12, an increase of about 20\%.
\end{corrige}

\begin{corrige}[Exercise 6 — Reading a seismic hazard curve]
\textbf{(a)} Hazard curve (logarithmic horizontal axis):
\begin{popfigure}
\begin{tikzpicture}
\begin{semilogxaxis}[
  width=11cm, height=7cm,
  xlabel={Return period $T$ (years)}, ylabel={PGA ($g$)},
  xmin=50, xmax=3000, ymin=0, ymax=0.36,
  xtick={72,225,475,975,2475},
  xticklabels={72,225,475,975,2475},
  grid=both, grid style={popline},
  legend pos=north west]
\addplot[color=popblue, mark=*, thick] coordinates {(72,0.05) (225,0.10) (475,0.15) (975,0.21) (2475,0.32)};
\addlegendentry{Hazard curve}
\addplot[color=poporange, dashed, thick] coordinates {(475,0) (475,0.15)};
\addplot[color=poporange, dashed, thick] coordinates {(50,0.15) (475,0.15)};
\end{semilogxaxis}
\end{tikzpicture}
\legende{Seismic hazard curve: PGA versus return period (log scale). The 475-year design point is highlighted.}
\end{popfigure}

\textbf{(b)} Reading the table (and the curve) at $T = 475$ years: design PGA $= \mathbf{0.15}\,g$.

\textbf{(c)} A 475-year return period means the annual probability of exceedance is
\[ P = \frac{1}{475} \approx 0.0021 = 0.21\% \text{ per year}. \]
Equivalently, the probability of exceedance in 50 years (typical building life) is $1 - (1 - 1/475)^{50} \approx 10\%$: this is the standard ``10\% in 50 years'' design level.

\textbf{(d)} Implications for building codes near the Cameroon Volcanic Line:
\begin{itemize}
\item \textbf{Seismic design requirement:} structures must be designed to resist horizontal accelerations of at least $0.15\,g$ — e.g. reinforced concrete frames with ductile detailing, bracing, and proper foundation anchorage — rather than unreinforced mud-block or adobe construction, which performs poorly even at moderate PGA.
\item \textbf{Zoning and enforcement:} the hazard curve justifies including the region in an official seismic zone of the national building code, with mandatory code enforcement, site-specific soil studies (soft soils amplify PGA), and restrictions on critical facilities (hospitals, schools) in the highest-hazard zones.
\end{itemize}
\end{corrige}

\begin{corrige}[Exercise 7 — Simple carbon-cycle box model]
\textbf{(a)} Radiative forcing:
\[ \Delta F = 5.35\ln\frac{C}{C_0} = 5.35 \times \ln\frac{550}{410} = 5.35 \times 0.294 \approx \SI{1.57}{W/m^2}. \]

\textbf{(b)} Equilibrium temperature rise:
\[ \Delta T = \lambda\,\Delta F = 0.8 \times 1.57 \approx \SI{1.26}{\kelvin}. \]

\textbf{(c)} Limitations of the simple model:
\begin{itemize}
\item It ignores \textbf{feedbacks} (water vapour, ice-albedo, cloud feedbacks) which can amplify the response — so the model may \emph{underestimate} warming.
\item It assumes \textbf{instant equilibrium}, ignoring ocean thermal inertia (which delays warming) and other forcings (aerosols, which cool; \ce{CH4}, \ce{N2O}, which warm) — so it may either under- or overestimate the real response on a given timescale.
\end{itemize}

\textbf{(d)} Effects on Cameroon's climate (any two):
\begin{itemize}
\item More frequent and intense \textbf{heatwaves} and shifting of the rainy seasons, stressing agriculture in the Adamawa and northern regions (longer dry seasons, advancing sahelian conditions).
\item \textbf{Sea-level rise} threatening the coastal lowlands around Douala, Kribi and Limbe (flooding, saltwater intrusion into coastal aquifers), and changes in rainfall intensity increasing flood and landslide risk in the Western Highlands.
\end{itemize}
\end{corrige}

\begin{corrige}[Exercise 8 — Contaminant breakthrough curve]
\textbf{(a)} Advective travel time:
\[ t_a = \frac{x}{v} = \frac{\SI{50}{m}}{\SI{0.5}{m/day}} = \boxed{100\ \text{days}}. \]

\textbf{(b)} Roles of dispersion and decay:
\begin{itemize}
\item \textbf{Dispersion} spreads the contaminant around the advective front: some solute arrives \emph{before} $t_a$ and the rise of $C(t)$ is smoothed into an S-shaped (sigmoid) curve instead of a sharp step.
\item \textbf{Decay} continuously destroys the contaminant during transport, so the plateau reached at long times lies \emph{below} $C_0$ (approximately $C_0 e^{-\lambda_d t_a} \approx 100 \times e^{-1} \approx \SI{37}{mg/L}$ for the advective timescale), rather than at $C_0$.
\end{itemize}

\textbf{(c)} Sketch of the breakthrough curve:
\begin{popfigure}
\begin{tikzpicture}
\begin{axis}[
  width=11cm, height=6.5cm,
  xlabel={Time $t$ (days)}, ylabel={$C/C_0$},
  xmin=0, xmax=300, ymin=0, ymax=1.1,
  grid=both, grid style={popline},
  legend pos=south east]
\addplot[color=popblue, very thick, smooth, domain=0:300, samples=120]
  {0.37*(0.5*(1+tanh((x-100)/35)))};
\addlegendentry{With dispersion and decay}
\addplot[color=popmuted, dashed, thick, smooth, domain=0:300, samples=120]
  {0.5*(1+tanh((x-100)/35))};
\addlegendentry{No decay (plateau $=1$)}
\addplot[color=poporange, dashed] coordinates {(100,0) (100,1.05)};
\node[poporange, anchor=west] at (axis cs:105,0.9) {$t_a = 100$ d};
\addplot[color=poppurple, dotted, thick] coordinates {(0,0.37) (300,0.37)};
\node[poppurple, anchor=south west] at (axis cs:5,0.39) {plateau $\approx e^{-\lambda_d t_a} \approx 0.37$};
\end{axis}
\end{tikzpicture}
\legende{Breakthrough curve at the observation well: sigmoid centred near $t_a$, with plateau lowered by first-order decay.}
\end{popfigure}

\textbf{(d)} Doubling the decay constant to $\lambda_d = \SI{0.02}{day^{-1}}$ does \textbf{not} change the advective arrival time:
\[ t_a = \frac{x}{v} = 100\ \text{days (unchanged)}, \]
because $t_a$ depends only on $x$ and $v$. Decay changes the \textbf{concentration level}, not the arrival time: the plateau would drop further to about $C_0 e^{-0.02 \times 100} = C_0 e^{-2} \approx 0.135\,C_0$.
\end{corrige}

\begin{corrige}[Exercise 9 — Essay — geohazards in Cameroon (20 marks)]
\textbf{Indicative mark scheme and model content.}

\textbf{(a) Definitions (6 marks).}
\begin{itemize}
\item \textbf{Geohazard} (2 mk): a geological process or phenomenon capable of causing loss of life, injury or damage to property and the environment.
\item \textbf{Hazard vs disaster vs vulnerability} (4 mk): the \emph{hazard} is the potentially damaging event itself (its probability and magnitude); \emph{vulnerability} is the degree of susceptibility of people, buildings and infrastructure to damage from that event; a \emph{disaster} occurs when a hazard strikes a vulnerable community and exceeds its capacity to cope. Examiner expects the key idea: disaster $=$ hazard $\times$ vulnerability (exposure).
\end{itemize}

\textbf{(b) Two geohazards affecting Cameroon (8 marks — 4 each: causes 2, effects 1, named example 1).}

\emph{Example 1 — Volcanic eruptions (Mount Cameroon):} caused by magma ascent along the Cameroon Volcanic Line, a zone of crustal weakness; recent eruptions (e.g. 1999, 2000) produced lava flows. Effects: destruction of farmland and vegetation, threat to nearby settlements (Buea, Limbe), ash fall, potential for associated earthquakes and toxic gases.

\emph{Example 2 — Toxic gas emissions (Lake Nyos, 1986; Lake Monoun, 1984):} magmatic \ce{CO2} seeps into deep lake water and accumulates under stratification; a trigger (landslide, wind mixing, density instability) causes a limnic eruption. Effects: a dense \ce{CO2} cloud flowed through valleys, asphyxiating about 1\,700 people and thousands of livestock around Lake Nyos.

(Alternatives: earthquakes along the CVL and associated faults; landslides on the steep, deforested, rain-soaked slopes of the Western Highlands, e.g. the Bafoussam area.)

\textbf{(c) Three mitigation strategies (6 marks — 2 each: description 1, limitation 1).}
\begin{itemize}
\item \textbf{Structural — degassing pipes at Lake Nyos:} pipes siphon \ce{CO2}-rich deep water to the surface, releasing gas gradually. Limitation: costly to install and maintain, and degassing must continue for many years.
\item \textbf{Structural — slope drainage and retaining structures} in landslide-prone highlands: lowers pore-water pressure, raises the factor of safety. Limitation: expensive, requires engineering expertise and maintenance; cannot protect all slopes.
\item \textbf{Non-structural — hazard mapping, land-use zoning and early warning:} hazard maps guide settlement away from high-risk zones (lava flow paths, unstable slopes, lake shorelines); monitoring of Mount Cameroon and the crater lakes with evacuation plans saves lives. Limitation: depends on enforcement, public education, funding and political will; relocation of communities is socially difficult.
\end{itemize}
\emph{Examiner expectation:} precise definitions, correctly named Cameroonian examples with dates, and mitigation measures explicitly classified as structural/non-structural with honest limitations.
\end{corrige}

\begin{corrige}[Exercise 10 — Structured question — global warming (20 marks)]
\textbf{(a) Greenhouse effect (7 marks).}
\begin{itemize}
\item \textbf{Natural effect (3 mk):} the atmosphere is largely transparent to short-wave solar radiation, which warms the Earth's surface; the surface re-emits long-wave infrared radiation, part of which is absorbed by greenhouse gases and re-radiated in all directions, including back towards the surface, keeping the Earth about $33\ \mathrm{K}$ warmer than it would otherwise be.
\item \textbf{Enhancement (1 mk):} human activities raise greenhouse gas concentrations, increasing absorption of outgoing infrared radiation and causing additional warming.
\item \textbf{Three gases with sources (3 mk):} \ce{CO2} — burning of fossil fuels and deforestation; \ce{CH4} — livestock, rice paddies, landfills, oil and gas operations; \ce{N2O} — nitrogen fertilisers and biomass burning (halocarbons from refrigerants also acceptable).
\end{itemize}

\textbf{(b) Calculation (3 marks).}
\[ \Delta T = \lambda \ln\frac{C}{C_0} = 3.0 \times \ln\frac{600}{300} = 3.0 \times \ln 2 = 3.0 \times 0.693 \approx \SI{2.1}{\kelvin}. \]

\textbf{(c) Four effects (6 marks — 1.5 each).}
\begin{itemize}
\item \textbf{Sea-level rise:} thermal expansion of seawater plus glacier and ice-sheet melt; flooding and saltwater intrusion threaten Douala, Kribi and the coastal aquifers of West Africa.
\item \textbf{Extreme events:} more intense rainfall episodes and floods (e.g. flash floods in Yaoundé and Douala), and more severe droughts in the Sudano-Sahelian north of Cameroon.
\item \textbf{Water resources:} reduced and more erratic river flows, shrinking of Lake Chad (a striking African example), stress on rain-fed agriculture.
\item \textbf{Ecosystems:} shifts in vegetation zones, coral bleaching, species migration or loss; northward advance of sahelian conditions into the Adamawa region.
\end{itemize}

\textbf{(d) Mitigation and adaptation (4 marks — 1 each).}
\begin{itemize}
\item \textbf{Mitigation:} (i) reforestation and reduced deforestation of the Congo Basin forest — highly feasible given Cameroon's forest cover, but requires funding and governance; (ii) development of hydroelectricity and renewable energy instead of fossil fuels — feasible given large hydro potential, though capital-intensive.
\item \textbf{Adaptation:} (i) drought-resistant crops and adjusted planting calendars in the north — feasible and low-cost with agricultural extension services; (ii) coastal protection and land-use planning in Douala–Kribi — feasible but expensive; mangrove restoration is a cheaper complementary option.
\end{itemize}
\emph{Examiner expectation:} clear physical mechanism in (a), correct use of $\ln 2$ in (b), effects linked to African/Cameroonian realities, and a balanced feasibility judgement in (d).
\end{corrige}

\begin{corrige}[Exercise 11 — Critical evaluation of an Environmental Impact Assessment (20 marks)]
\textbf{(a) Purpose and stages of an EIA (4 marks).}

\textbf{Purpose:} to identify, predict and evaluate the environmental (including geological) impacts of a proposed project \emph{before} it is approved, so that adverse effects can be avoided, reduced or compensated. \textbf{Main stages:} screening; scoping; baseline studies; impact prediction and evaluation; proposal of mitigation measures; preparation of the environmental impact statement; public consultation and review; decision; and post-project monitoring and auditing.

\textbf{(b) Four gaps or weaknesses (8 marks — 2 each).}
\begin{itemize}
\item \textbf{Dismissal of seismic risk:} the site lies in a region of moderate seismicity, so the claim of ``no significant geological risks'' is unjustified; a dam is a critical structure whose failure would flood the growing downstream town.
\item \textbf{No slope-instability / reservoir-triggered landslide study:} steep, deforested slopes upstream are prone to landslides, which could be triggered or reactivated by reservoir filling (rising pore-water pressures) and could generate displacement waves or block the reservoir.
\item \textbf{No sedimentation study:} erosion from deforested slopes will deliver sediment that silts up the reservoir, shortening its lifespan and reducing storage.
\item \textbf{Inadequate baseline data:} a single dry-season water-quality campaign cannot represent seasonal variability; flood design based on such data ignores rainfall variability and climate change, which alter flood magnitudes and dam safety margins.
\item \textbf{No greenhouse-gas assessment:} flooded vegetation in tropical reservoirs decomposes and emits \ce{CO2} and \ce{CH4}, which should be quantified in the project's climate balance.
\end{itemize}

\textbf{(c) Concrete improvements (6 marks — matched to gaps).}
\begin{itemize}
\item Commission a \textbf{site-specific seismic hazard assessment} (fault mapping, historical seismicity, probabilistic hazard curve) and design the dam for the appropriate design PGA; install strong-motion instruments.
\item Carry out a \textbf{geotechnical and slope-stability survey} of the reservoir rim; map unstable zones, plan drainage/anchoring or exclusion zones, and manage reservoir filling rates with slope monitoring (piezometers, inclinometers).
\item Perform a \textbf{sediment yield and reservoir sedimentation study}; implement catchment reforestation and erosion control upstream.
\item Extend baseline monitoring to \textbf{at least one full hydrological year} (both seasons); incorporate climate-change scenarios into flood design (larger design flood, revised spillway capacity).
\item Include a \textbf{reservoir greenhouse-gas study} (pre-flooding vegetation survey, emission estimates) and consider biomass clearing before impoundment.
\end{itemize}

\textbf{(d) Public participation and long-term monitoring (2 marks).}
Public participation ensures that affected communities (e.g. the downstream town) can voice concerns, improves the quality and legitimacy of the assessment, and aids acceptance of mitigation or resettlement measures. Long-term monitoring verifies that predictions were correct, detects unforeseen impacts (slope movements, water quality decline, seismicity) early, and allows management measures to be adjusted over the dam's lifetime.
\emph{Examiner expectation:} weaknesses must be \emph{justified from the scenario} (seismicity, deforested slopes, downstream town, single-season baseline), and each improvement must correspond to an identified gap.
\end{corrige}

\newpage

\coursec{Summary sheet}

\begin{popfigure}
\begin{center}\tcbox[colback=popink,colframe=popink,coltext=white,arc=9pt,boxsep=4pt,left=6pt,right=6pt]{\bfseries\small Environmental Geology (GEO08)}\end{center}
\vspace{-2pt}
\begin{tcbraster}[raster columns=3,raster equal height=rows,raster column skip=6pt,raster row skip=6pt,enhanced,blankest]
\begin{tcolorbox}[enhanced,colback=white,colframe=poporange,boxrule=0.9pt,arc=3pt,title={Foundations \& Geohazards},fonttitle=\bfseries\scriptsize,coltitle=white,colbacktitle=poporange,left=4pt,right=4pt,top=2pt,bottom=2pt,boxsep=1pt,toptitle=1.5pt,bottomtitle=1.5pt,halign=left]
\begin{itemize}[nosep,leftmargin=*,itemsep=1pt,font=\scriptsize]
\item Definition, scope
\item Earthquakes, volcanoes
\item Landslides, floods
\end{itemize}
\end{tcolorbox}
\begin{tcolorbox}[enhanced,colback=white,colframe=popgreen,boxrule=0.9pt,arc=3pt,title={Mitigation 1},fonttitle=\bfseries\scriptsize,coltitle=white,colbacktitle=popgreen,left=4pt,right=4pt,top=2pt,bottom=2pt,boxsep=1pt,toptitle=1.5pt,bottomtitle=1.5pt,halign=left]
\begin{itemize}[nosep,leftmargin=*,itemsep=1pt,font=\scriptsize]
\item Seismic \& volcanic effects
\item Prediction, zoning
\end{itemize}
\end{tcolorbox}
\begin{tcolorbox}[enhanced,colback=white,colframe=popblue!70!black,boxrule=0.9pt,arc=3pt,title={Mitigation 2},fonttitle=\bfseries\scriptsize,coltitle=white,colbacktitle=popblue!70!black,left=4pt,right=4pt,top=2pt,bottom=2pt,boxsep=1pt,toptitle=1.5pt,bottomtitle=1.5pt,halign=left]
\begin{itemize}[nosep,leftmargin=*,itemsep=1pt,font=\scriptsize]
\item Mass movement control
\item Flood \& coastal defences
\end{itemize}
\end{tcolorbox}
\begin{tcolorbox}[enhanced,colback=white,colframe=poppurple,boxrule=0.9pt,arc=3pt,title={Global Warming},fonttitle=\bfseries\scriptsize,coltitle=white,colbacktitle=poppurple,left=4pt,right=4pt,top=2pt,bottom=2pt,boxsep=1pt,toptitle=1.5pt,bottomtitle=1.5pt,halign=left]
\begin{itemize}[nosep,leftmargin=*,itemsep=1pt,font=\scriptsize]
\item Greenhouse effect
\item Sea level, climate shifts
\end{itemize}
\end{tcolorbox}
\begin{tcolorbox}[enhanced,colback=white,colframe=poppink,boxrule=0.9pt,arc=3pt,title={Pollution 1},fonttitle=\bfseries\scriptsize,coltitle=white,colbacktitle=poppink,left=4pt,right=4pt,top=2pt,bottom=2pt,boxsep=1pt,toptitle=1.5pt,bottomtitle=1.5pt,halign=left]
\begin{itemize}[nosep,leftmargin=*,itemsep=1pt,font=\scriptsize]
\item Air, water, soil sources
\item Transport in geosphere
\end{itemize}
\end{tcolorbox}
\begin{tcolorbox}[enhanced,colback=white,colframe=popgold,boxrule=0.9pt,arc=3pt,title={Pollution 2},fonttitle=\bfseries\scriptsize,coltitle=white,colbacktitle=popgold,left=4pt,right=4pt,top=2pt,bottom=2pt,boxsep=1pt,toptitle=1.5pt,bottomtitle=1.5pt,halign=left]
\begin{itemize}[nosep,leftmargin=*,itemsep=1pt,font=\scriptsize]
\item Waste management
\item Remediation, integration
\end{itemize}
\end{tcolorbox}
\end{tcbraster}

\legende{Mind map of Chapter GEO08: Environmental Geology.}
\end{popfigure}

\begin{aretenir}[Summary Sheet — Environmental Geology]
\textbf{Key definitions}
\begin{itemize}
\item \cle{Environmental geology}: the application of geological knowledge to the interactions between humans and the geological environment, in order to understand, predict and reduce \cle{geohazards}, manage resources and control pollution.
\item \cle{Geohazard}: a natural geological process or phenomenon (earthquake, volcanic eruption, landslide, flood, tsunami, subsidence) that threatens life, property or the environment. A hazard becomes a \emph{disaster} only when it affects a vulnerable population.
\item \cle{Risk} $=$ hazard $\times$ vulnerability $\times$ exposure. Reducing any factor reduces risk.
\item \cle{Greenhouse effect}: natural trapping of outgoing infrared radiation by atmospheric gases (\ce{CO2}, \ce{CH4}, \ce{H2O} vapour, \ce{N2O}); its \emph{enhancement} by human emissions drives \cle{global warming}.
\item \cle{Pollution}: introduction of harmful substances or energy into air, water or soil at rates exceeding the environment's capacity to dilute, disperse or degrade them.
\end{itemize}

\textbf{Major geohazards — essentials}
\begin{itemize}
\item \emph{Earthquakes}: sudden release of elastic strain along faults; measured by magnitude (energy, logarithmic Richter/moment scale) and intensity (Mercalli, observed effects). Secondary effects: ground shaking, liquefaction, tsunamis, fires.
\item \emph{Volcanic eruptions}: effusive (lava flows, low viscosity basalt — Mount Cameroon) versus explosive (viscous magma, pyroclastic flows, ash falls, lahars, gas emissions — Lake Nyos \ce{CO2} disaster, 1986, is a limnic eruption example).
\item \emph{Mass movements}: creep, slides, slumps, rockfalls, mudflows; controlled by slope angle, water content, rock type, vegetation and triggering events (rain, shaking).
\item \emph{Floods and subsidence}: river floods, flash floods, coastal flooding; subsidence from groundwater or mineral extraction.
\end{itemize}

\textbf{Mitigation methods (know at least two per hazard)}
\begin{itemize}
\item Earthquakes: seismic building codes, hazard microzonation maps, early-warning systems, education and drills. Prediction remains unreliable — \emph{forecasting} (probabilistic) is possible, exact prediction is not.
\item Volcanoes: monitoring (seismographs, gas analysis, ground deformation, thermal imaging), exclusion zones, evacuation plans, lava diversion barriers, degassing pipes (Lake Nyos).
\item Landslides: drainage of slopes, retaining walls, regrading (reducing slope angle), reforestation, avoiding construction on unstable slopes.
\item Floods: dams and levees, channelisation, floodplain zoning, wetland preservation, reforestation of catchments.
\end{itemize}

\textbf{Global warming — key points}
\begin{itemize}
\item Causes: burning of fossil fuels, deforestation, cement production, agriculture (\ce{CH4}, \ce{N2O}).
\item Effects: rising mean temperatures, melting glaciers and ice caps, thermal expansion of seawater $\Rightarrow$ \cle{sea-level rise}; shifting rainfall patterns, more frequent extreme events, ocean acidification, stress on ecosystems and agriculture.
\item Geological evidence: ice cores, sediment records, fossils, isotope ratios (\ce{^{18}O}/\ce{^{16}O}) reveal past climates.
\item Responses: mitigation (renewable energy, reforestation, carbon capture) and adaptation (coastal defences, resistant crops).
\end{itemize}

\textbf{Pollution — key points}
\begin{itemize}
\item \emph{Air}: \ce{SO2}, \ce{NO_x}, particulates $\Rightarrow$ acid rain, respiratory disease; \emph{water}: sewage, fertilisers (eutrophication), oil spills, heavy metals, industrial effluent; \emph{soil}: pesticides, landfill leachate, mining waste.
\item Transport: pollutants move through rivers, groundwater (controlled by \cle{permeability} and aquifer geometry) and the atmosphere; geology determines dispersal or concentration.
\item Management: waste reduction, recycling, sanitary landfills (lined, sited on impermeable rocks), water treatment, bioremediation, environmental impact assessment (EIA) before projects.
\end{itemize}

\textbf{Common mistakes to avoid}
\begin{itemize}
\item Confusing \emph{magnitude} (single value, energy) with \emph{intensity} (varies with location).
\item Saying earthquakes can be precisely predicted — they cannot; only hazard zones and probabilities are known.
\item Confusing the \emph{natural} greenhouse effect (essential for life) with its \emph{enhancement} (the problem).
\item Attributing Lake Nyos to a classic magmatic eruption: it was a \ce{CO2} limnic burst.
\item Writing ``pollution disappears'': it is transported, transformed or accumulated — trace the pathway.
\end{itemize}

\textbf{GCE tips}
\begin{itemize}
\item Always \emph{define the hazard}, describe its \emph{causes and effects}, then give \emph{mitigation measures with named examples} (Mount Cameroon, Lake Nyos, Benue trough floods earn credit).
\item In structured questions, link mitigation to the specific cause (e.g.\ slope drainage $\leftrightarrow$ water-triggered landslides).
\item For essays on global warming or pollution, structure: causes $\rightarrow$ mechanisms $\rightarrow$ effects $\rightarrow$ solutions, with a labelled diagram where possible.
\item Use correct terminology: vulnerability, permeability, eutrophication, liquefaction, microzonation.
\end{itemize}
\end{aretenir}

\findecours

\end{document}
